\documentclass[11pt]{article}

\usepackage[letterpaper,margin=1in,headheight=14pt]{geometry}
\usepackage[T1]{fontenc}
\usepackage[utf8]{inputenc}
\usepackage{lmodern}
\usepackage{microtype}
\usepackage{amsmath,amssymb,amsthm,mathtools,mathrsfs}
\usepackage{booktabs,tabularx,array}
\usepackage{xcolor}
\usepackage{enumitem}
\usepackage{fancyhdr}
\usepackage{titlesec}
\usepackage[hidelinks]{hyperref}
\usepackage[nameinlink,noabbrev]{cleveref}

\definecolor{ink}{HTML}{1F2937}
\definecolor{accent}{HTML}{1F5A7A}
\definecolor{soft}{HTML}{F3F6F8}
\definecolor{warn}{HTML}{8A4B20}
\hypersetup{
  pdftitle={Prime-Zeta Divisor Tomography: Local Zero Processes, Frozen-Edge Rigidity, and Derivative Caustics},
  pdfauthor={KEIO},
  colorlinks=true,
  linkcolor=accent,
  citecolor=accent,
  urlcolor=accent
}

\titleformat{\section}{\Large\bfseries\color{ink}}{\thesection}{0.7em}{}
\titleformat{\subsection}{\large\bfseries\color{ink}}{\thesubsection}{0.6em}{}
\titleformat{\subsubsection}{\normalsize\bfseries\color{ink}}{\thesubsubsection}{0.5em}{}
\setlist{itemsep=0.25em,topsep=0.5em}
\setlength{\parindent}{1.15em}
\setlength{\parskip}{0.15em}
\allowdisplaybreaks

\pagestyle{fancy}
\fancyhf{}
\fancyhead[L]{\small Prime-zeta divisor tomography and frozen-edge rigidity}
\fancyhead[R]{\small KEIO}
\fancyfoot[C]{\thepage}

\newtheorem{theorem}{Theorem}[section]
\newtheorem{proposition}[theorem]{Proposition}
\newtheorem{lemma}[theorem]{Lemma}
\newtheorem{corollary}[theorem]{Corollary}
\theoremstyle{definition}
\newtheorem{definition}[theorem]{Definition}
\newtheorem{problem}[theorem]{Problem}
\newtheorem{conjecture}[theorem]{Conjecture}
\newtheorem{heuristic}[theorem]{Heuristic}
\theoremstyle{remark}
\newtheorem{remark}[theorem]{Remark}

\newcommand{\PP}{\mathbb P}
\newcommand{\cP}{\mathcal P}
\newcommand{\dd}{\,\mathrm d}
\newcommand{\eps}{\varepsilon}
\newcommand{\C}{\mathbb C}
\newcommand{\R}{\mathbb R}
\newcommand{\N}{\mathbb N}
\newcommand{\one}{\mathbf 1}
\DeclareMathOperator{\Tr}{Tr}
\DeclareMathOperator{\supp}{supp}
\DeclareMathOperator{\divisor}{div}
\DeclareMathOperator{\type}{type}
\DeclareMathOperator{\PV}{PV}

\begin{document}
\hypersetup{pageanchor=false}

\begin{titlepage}
\centering
\vspace*{1.0in}
{\huge\bfseries\color{ink} Prime-Zeta Divisor Tomography:\par}
\vspace{0.10in}
{\huge\bfseries\color{ink} Local Zero Processes, Frozen-Edge Rigidity,\par}
\vspace{0.10in}
{\huge\bfseries\color{ink} and Derivative Caustics\par}
\vspace{0.36in}
{\Large Exact densities, phase recovery, small-target monodromy,\par
and inverse prime spectroscopy\par}
\vspace{0.8in}
{\large KEIO\par}
\vspace{0.25in}
{\large Final combined manuscript, Version 4.0\par}
\vspace{0.08in}
{\large 3 September 2026\par}
\vfill
\begin{minipage}{0.88\textwidth}
\small
\textbf{Scope.}
This paper concerns the absolute-convergence half-plane of prime-supported Dirichlet
series.  General almost-periodic projection and Jessen theories are used as foundations.
The principal results are smooth and marked divisor laws, complete local point-process
limits, deterministic hard-core criteria, an exact phase-fidelity frontier, quantized
small-target escape with divisor monodromy, inverse support and phase recovery,
an exact large-deviation profile at the outer edge, a universal rooted edge divisor,
and the complete internal-band geometry of the derivative tower.
Component masses and barycentres recover their boundary primes, while moving high
derivatives produce exponentially accurate zero lattices.  The exact projection edge is
not claimed as a new general phenomenon.  Numerical experiments are kept separate
from the unconditional statements, and uncertified decimal thresholds are not used as
hypotheses of the principal theorems.
\end{minipage}
\vfill
\end{titlepage}

\hypersetup{pageanchor=true}
\pagenumbering{roman}

\begin{abstract}
Let
\[
  F_c(s)=\sum_{p}c_pp^{-s},\qquad \Re s>1,
\]
under the coefficient hypotheses stated below, and write $F_\omega$ when
$|\omega_p|=1$.  We identify vertical translation with a uniquely ergodic
prime-phase hull.  Its one-point Kac--Rice density $G(\sigma,a)$ is smooth,
radial, and phase blind.  On every compact horizontal strip, multiplicity,
distinct-point, and simple-point counts have the same exact linear main term;
a marked law records the derivative at an $a$-point.  For the untwisted prime
zeta function the count extends to the full half-plane $\Re s>1$.  This proves
the literal distinct-point height law of Belovas, \v{C}epaityt\.{e}, and
Sabaliauskas and strengthens their linear-counting conjecture to an exact
asymptotic.  The projection edge itself is used from classical
almost-periodic theory and is not claimed as a new general phenomenon.

At every fixed derivative order, the empirical law of the vertically
translated $a$-point divisor of any deterministic twist converges to the same
stationary Haar divisor process.  All factorial $k$-point correlations are
therefore given by joint Kac--Rice densities, and compact-window zero bounds
turn them into finite formulae for the Laplace functional, count laws, void
probabilities, and Palm nearest-neighbour distributions.  The self-inclusive
smoothed two-point function is Bohr almost periodic; its factorial version is
an almost-periodic function plus a $C_0$ diagonal correction.  An explicit
variational lower bound for $|F_\omega'(\rho)|$ produces a deterministic
simple-zero region and a twist-independent planar hard core.  In the
complementary nondegenerate regime,
\[
 K_2^{(m,a)}(s,s+h)=\frac{|h|^2}{16}\,\mathcal C_m(\Re s,a)+O(|h|^3),
\]
and the full confluent hierarchy has the corresponding Vandermonde repulsion.

Anchored edge data restore precisely the information erased by stationary
averaging.  If $\rho=\sigma+it$ is a nonzero $a$-point and
$\varepsilon=A_c(\sigma)-|a|$, then the exact phase defect is
\[
 \sum_p |c_p|p^{-\sigma}
 \left|1-\frac{c_p}{|c_p|}\frac{|a|}{a}p^{-it}\right|^2
 =2\varepsilon.
\]
Hence a nonzero-target edge sequence recovers all coefficient phases, while a
zero-edge sequence recovers their relative phases, with common rotation as the
only ambiguity.  Together with sparse samples of the right-edge curve, this is
inverse complete for the coefficient system.  More strongly, after recentering
at any zero whose real part approaches the universal zero edge
$S=\sigma_{\rm e}$, the normalized analytic germ converges, uniformly over all
twists, to
\[
 H_{\rm e}(z)=2^{1-(S+z)}-\mathcal P(S+z).
\]
Thus the rooted divisor and every rooted Palm law collapse to one deterministic
profile.  Boundary-attaining twists and equal edge germs are classified exactly
by common rotation and the vertical Kronecker flow; in shrinking edge layers,
distinct zero ordinates separate at least logarithmically in the reciprocal
defect.  This gives a phase-readable, uniformly discrete edge divisor even
though its entire stationary local law is phase blind.

The phase defect also converts the outer-edge large-deviation theorem into an
exact density--fidelity frontier.  If $\mathcal C(d)$ is the vertical intensity
of zeros whose weighted relative-phase distortion is at most $d$, then, with
$\beta_{\rm e}=(\sigma_{\rm e}-1)^{-1}$ and
$\ell_d=\log(2/d)$,
\[
 -\log\mathcal C(d)
 =K_{\rm e}(d/2)^{-\beta_{\rm e}}\ell_d^{-(1+\beta_{\rm e})}
 \left(1+(1+\beta_{\rm e})\frac{\log\ell_d}{\ell_d}
 +O(\ell_d^{-1})\right).
\]
This is a limiting-intensity theorem, not a first-occurrence estimate.  Its
finite-prime consequence is stated only as a sufficient reconstruction
certificate.

Finally, the quantized $a\to0$ escape law is lifted from measures to divisors.
On every compact subband of the hard-core region, the bounded $a$-point divisor
consists of uniformly separated holomorphic continuations of the zero divisor;
a zero-free corridor separates these persistent branches from a far-right
logarithmic lattice.  A loop around $a=0$ fixes every persistent branch and
acts by a unit shift on the escaping lattice.  The escape mass and reciprocal lattice spacing---equivalently, the physical
displacement of one monodromy shift---recover the least support prime, while
the bounded mass has an even $C^\infty$ expansion in $|a|$.

The paper also classifies the internal dominance bands of the derivative tower
and its moving-order zero lattices.  It proves no statement across
$\Re s=1$ and no new case of the Riemann hypothesis, Montgomery's
pair-correlation conjecture, or a prime-gap conjecture.
\end{abstract}

\tableofcontents
\vspace{0.5em}
\noindent\textbf{Keywords.}
Prime zeta function; almost periodic Dirichlet series; local zero process; Palm law;
Kac--Rice formula; hard core; large deviations; phase recovery; rate--distortion;
divisor monodromy; derivative caustic; inverse spectral problem.

\medskip
\noindent\textbf{MSC 2020.}
11M41, 30B50, 30D20; secondary 11M26, 47B10.

\clearpage
\pagenumbering{arabic}

\section{Introduction and main conclusions}
\label{sec:introduction}

The prime zeta function
\[
 \cP(s)=\sum_{p\in\PP}p^{-s}
\]
is absolutely convergent in $\Re s>1$ and is connected to the Riemann zeta function by
M\"obius inversion,
\[
 \cP(s)=\sum_{k\ge1}\frac{\mu(k)}{k}\log\zeta(ks),
\]
on compatible branches.  Its continuation has a natural boundary at $\Re s=0$ and a
dense family of singularities generated by the zeros and pole of $\zeta$; see
Titchmarsh~\cite{Titchmarsh1986}, Fr\"oberg~\cite{Froberg1968}, and
Landau--Walfisz~\cite{LandauWalfisz1920}.
The present paper asks a different question: how much of the zero geometry is already
forced, and how much arithmetic can be recovered, inside the elementary half-plane
$\Re s>1$?

Belovas, \v{C}epaityt\.{e}, and Sabaliauskas~\cite{BCS2025} numerically studied the
rightmost zeros and formulated conjectures concerning their edge, their linear number,
and the uniform distribution of normalized heights.  The right projection itself is a
polygonal property of an almost-periodic Dirichlet series.  General criteria are due to
Jessen--Tornehave~\cite{JessenTornehave1945} and, in a direct projection form, to
Sepulcre--Vidal~\cite{SepulcreVidal2022}.  A recent preprint of
Camargo~\cite{Camargo2025} independently specializes equivalent-series methods to the
prime zeta function and obtains the same projection edge.  Accordingly, the new content
claimed here is not the mere existence of that boundary, nor the classical existence
of a multiplicity-weighted Jessen mean count.  It is the explicit phase-blind density,
the upgrade to distinct and simple divisors, marked slope laws, quantized escape,
endpoint transfer, inverse theorems, and derivative geometry built on the boundary.
In particular, we do not claim Conjecture~1 of \cite{BCS2025} as a new result; the
literal distinct-point forms and sharpenings of its metric Conjectures~2--3 are the
relevant application of the new package.

The conclusions may be organized into fourteen layers.  The first nine concern
the horizontal divisor, edge, derivative, spectral, and arithmetic structures;
the last five pass to the complete fixed-level local process, rooted edge
rigidity, the exact phase-fidelity frontier, and pointwise small-target surgery.

\begin{enumerate}[label=\textup{(\Roman*)},leftmargin=2.4em]
\item \textbf{Smooth density.}
Every unimodular twist has the same smooth $a$-point density $G(\sigma,a)$ in
$\Re s>1$.  Its zero density is positive exactly before the polygon edge and is
infinitely flat at that edge.  We determine the flatness by an explicit Cram\'er
large-deviation law whose power exponent is $(\sigma_{\rm e}-1)^{-1}$ and whose
logarithmic correction is forced by the prime number theorem.

\item \textbf{Simple and marked divisors.}
For $1<\gamma<\sigma_{\rm e}$,
\[
 M_{\omega,\gamma}(T)=c_\gamma T+o(T),\qquad
 c_\gamma=\int_\gamma^{\sigma_{\rm e}}G(\sigma,0)\dd\sigma>0.
\]
Multiplicity, distinct, and simple counts have this same main term.  A marked
Kac--Rice law records the derivative at each $a$-point and proves a quartic law for
near-critical slopes.  The normalized heights are uniformly distributed, giving the
literal distinct-point versions of the 2025 metric conjectures.

\item \textbf{The endpoint.}
For the untwisted $\cP$, the exponential envelope
$e^{\cP}=\zeta H$ controls the layer $1<\Re s<1+\delta$.  It yields
\[
 \#\{\cP(\rho)=a:\Re\rho>1,\ 0<\Im\rho<T\}
 =C_aT+o(T)
\]
for every fixed $a$ under each of the three counting conventions.  Simple points in
every occupied compact band recur with bounded vertical gaps.  In particular, the weak
endpoint prediction $M_1(T)=O(T)$ is sharpened to $M_1(T)\sim c_1T$.

\item \textbf{Quantized small-target escape.}
As $a\to0$ through nonzero values, exactly $\log2/(2\pi)$ units of horizontal
$a$-point mass escape to infinity.  The escaping points form a simple asymptotic
lattice.  For a general infinite prime support the jump is $\log p_0/(2\pi)$ and
therefore recovers its least prime $p_0$.

\item \textbf{Tomography and conservation.}
All nonzero $a$-point Riesz measures determine the coefficient moduli of a bounded
prime-supported Dirichlet series.  A single vertical limiting value law already recovers
the full prime sequence, and the nonzero $a$-point right-edge curve is an even smaller
complete datum.  Integrating over the target plane gives an exact prime-spectral mass
identity and, for character supports, one scalar recovers the modulus radical.  In the
derivative tower, one bounded zero component recovers the two prime labels on its
boundary from its mass and barycentre.

\item \textbf{Derivative caustics.}
The zero edge of $(-1)^m\cP^{(m)}$ is the pinch point at which the right-hand value gap
of $(-1)^{m-1}\cP^{(m-1)}$ changes topology.  The edge sequence is a lossless encoding
of the primes.  More generally every prime owns at most one dominance band, every prime
is born at a finite derivative order, and the bands are disjoint, ordered, and persist
through the tower under the explicit shift in (6.23).

\item \textbf{Crystalline moving-order limit.}
After the rescaling $\sigma\mapsto\sigma/m$, the horizontal zero-density measures
converge vaguely to a weighted prime-indexed picket fence.  Before vertical averaging,
the actual zeros in each fixed crossover window form an exponentially accurate simple
lattice.  Its horizontal velocity, vertical spacing, and phase offset recover the two
adjacent primes and their twist ratio.

\item \textbf{Equivalent spectral forms.}
The same density constant is a Weyl coefficient, a heat coefficient, an Abelian
spectral-zeta residue, a Dixmier trace, and the exponential type of a symmetric
canonical product.  These are Tauberian repackagings, not additional zero-distribution
hypotheses.

\item \textbf{Arithmetic transfer.}
Reciprocal prime series have matching pole and convergence edges.  Root-of-unity
filters split ordered prime factorizations into equidistributed length classes while
retaining a completely explicit family of off-axis poles.
\item \textbf{Complete fixed-level local divisor process.}
For every deterministic unimodular twist, vertical translation of the
$a$-point divisor has a unique stationary local limit, namely the divisor of the
Haar prime-phase model.  Its factorial moment densities are the joint Kac--Rice
functions $K_k$, for every $k\geq1$.  Uniform compact-window zero bounds make
the correlation hierarchy determinate: the Laplace functional, every
finite-window count law, void probabilities, and Palm nearest-neighbour laws
are finite inclusion--exclusion expressions in the $K_k$.  Thus the local law
is universal in the twist even though an individual anchored configuration is
not.  The self-inclusive smoothed pair function is Bohr almost periodic; the
factorial pair function lies in $AP+C_0$ and has product mean.

\item \textbf{Hard-core transition and diagonal law.}
Put
\[
 \mathfrak m(\sigma)=\sup_{\lambda>0}
 \frac{2^{-\sigma}-\sum_{p\geq3}p^{-\sigma}
       |1-\lambda\log(p/2)|}{\lambda}.
\]
Whenever $\mathfrak m(\sigma)>0$, every zero of every unimodular twist on
$\Re s=\sigma$ is simple and satisfies a twist-independent derivative lower
bound; Taylor's theorem gives an explicit positive separation radius.  In the
nondegenerate complementary regime the two-point density vanishes quadratically
on the diagonal, with coefficient determined by the third jet.  These are
deterministic and Kac--Rice forms of fixed-level repulsion.  Any quoted decimal
transition value is kept separate from the theorem unless supported by an
interval certificate.  The same dual construction gives symbolic hard-core
and ordinate-separation criteria at every fixed derivative level.

\item \textbf{Rooted edge rigidity and ordinate phase recovery.}
Horizontal densities, edge locations, and the entire stationary local law
depend only on the coefficient moduli.  Nevertheless, the defect identity at a
single nonzero $a$-point gives
\[
 \left|\frac{c_p}{|c_p|}-\frac{a}{|a|}p^{it}\right|
 \leq
 \left(\frac{2\{A_c(\sigma)-|a|\}}
 {|c_p|p^{-\sigma}}\right)^{1/2}.
\]
An edge-approaching sequence therefore recovers all phases.  For zeros it
recovers the ratios of the phases, modulo their unavoidable common rotation.
Together with the right-edge curve, which recovers the moduli, this gives a
hybrid inverse datum for the full coefficient system.  At the zero
edge, recentering at any near-edge root gives the universal analytic limit
\[
 H_{\rm e}(z)=2^{1-(\sigma_{\rm e}+z)}-\cP(\sigma_{\rm e}+z),
\]
uniformly over the twist.  The rooted divisors and rooted Palm laws collapse to
its deterministic divisor.  Equality of edge germs, with or without an
imaginary translation, is equivalent to equality of the twists modulo common
rotation and the corresponding vertical flow.  In a shrinking edge layer the
minimum ordinate separation grows at least logarithmically in the reciprocal
defect.

\item \textbf{Exact phase-fidelity frontier.}
For a zero $\rho=\sigma+it$, the weighted chordal distortion of the natural
relative-phase decoder is exactly
\[
 \mathcal D_\omega(\rho)
 =2\left(\sum_{p\ge3}p^{-\sigma}-2^{-\sigma}\right).
\]
Consequently a distortion threshold is exactly a horizontal edge layer.  The
outer-edge Cram\'er law therefore gives the complete logarithmic asymptotic for
its vertical intensity $\mathcal C(d)$ and, after inversion, an exact
density--fidelity frontier.  The induced condition for learning all phases at
primes $p\le Q$ is a sufficient certificate with explicitly computable
intensity; neither optimality nor a deterministic first-hit height is claimed.

\item \textbf{Small-target divisor surgery and monodromy.}
On compact subbands of the deterministic hard-core region, every zero continues
to a unique holomorphic $a$-point branch, uniformly in the twist and height,
and every bounded $a$-point arises this way after inserting fixed horizontal
buffers.  A growing zero-free corridor separates this persistent divisor from
the far-right escaping lattice.  When the target winds once around $0$, every
persistent branch has identity monodromy while the escaping lattice undergoes
a bilateral unit shift.  The escaping mass and reciprocal lattice step---equivalently, the physical
displacement of one shift---encode the least prime of the support; the bounded contribution has an even
$C^\infty$ expansion in $|a|$ to every finite order.

\end{enumerate}

\subsection{What is and is not settled}

This revision settles, in the following fixed-scale sense, the questions left
as Problems~3 and~6 in Version~2.0.  For Problem~3, at every fixed derivative
order the complete stationary local divisor process is now determined.  Its
factorial $k$-point functions, finite-window count and void laws, and Palm
nearest-neighbour and consecutive-gap laws are given by exact finite
Kac--Rice/inclusion--exclusion formulae.  The corrected pair statement is that
the self-inclusive smoothed pair function is Bohr almost periodic, whereas the
factorial pair function is an almost-periodic term plus a compactly supported
diagonal correction.  The local theory also gives a symbolic deterministic
hard core on one side of the support transition and Vandermonde, in particular
quadratic pair, repulsion in the nondegenerate bulk.  This is a qualitative
limit theorem on each fixed compact window; it supplies neither a height-error
term nor an elementary closed form for the gap density.

For Problem~6, one no longer needs the complete moving derivative tower in
order to recover phases.  At a single fixed level, the ordinates of supplied
near-edge nonzero $a$-points recover all coefficient phases.  Supplied
near-edge zeros recover all relative phases; the remaining common rotation is
unavoidable because it leaves the zero set unchanged.  The rooted zero-edge
profile is universal, and its anchored germ determines the twist modulo this
common rotation; after quotienting the germ by imaginary translation it
determines the corresponding vertical-flow orbit.  Horizontal divisors alone
remain phase blind, even when their entire stationary finite-order correlation
hierarchy is retained.

Three further refinements should be separated from those two closures.  First,
the outer-edge large-deviation law has a deterministic rooted counterpart:
every near-edge zero, after recentering and normalization, converges to the
same analytic profile $H_{\rm e}$, and the associated rooted divisors and Palm
laws collapse to its divisor.  Second, the exact zero phase defect turns the
edge law into an exact asymptotic for the limiting intensity of any prescribed
small distortion, together with its inverse density--fidelity frontier.  This
does not identify the first height at which that distortion occurs, and the
finite-prime reconstruction condition extracted from it is sufficient rather
than necessary.  Third, the small-target measure limit is upgraded to a
pointwise divisor theorem in the hard-core subbands: bounded zeros persist as
holomorphic branches, a zero-free corridor separates them from the escaping
logarithmic lattice, and the two families have respectively identity and
shift monodromy.  The corresponding bounded mass has an even smooth
asymptotic jet.

The paper proves no new case of the Riemann hypothesis and no case of the twin
prime, Polignac, Cram\'er, or Montgomery pair-correlation conjectures.  The zeros
studied here lie in the absolute-convergence half-plane and have fixed-scale
almost-periodic statistics; they must not be conflated with the high zeros of
the Riemann zeta function.  Nor do the exact prime-gap dictionaries later in the
paper estimate a prime gap: they are inverse identities only.

The exact projection edge is likewise not claimed as a new general theorem.
It follows from classical analytic almost-periodic projection theory, and an
independent prime-zeta specialization was given by Camargo.  The distinct- and
simple-divisor laws, local Kac--Rice hierarchy, deterministic hard core,
rooted edge profile, ordinate phase recovery, phase-fidelity law, small-target
divisor surgery, and inverse statements are the specialized contributions
asserted here, subject to the priority qualification at the end of the paper.
For arbitrary phases the universal zero-free statement remains the
\emph{open} half-plane $\Re s>\sigma_{\rm e}$: an adversarial twist can attain
its boundary, whereas the untwisted prime zeta function does not.

\section{Almost-periodic hulls and exact projection sets}
\label{sec:hulls}

Let $c=(c_p)_{p\in\PP}$ be nonzero and suppose that
$|c_p|\ll(\log p)^M$ for some fixed $M\ge0$.  Put
\[
 F_c(s)=\sum_pc_pp^{-s},\qquad
 A_c(\sigma)=\sum_p|c_p|p^{-\sigma},\qquad \sigma>1.
\]
For fixed $\sigma$, define the inner polygonal radius
\[
 B_c(\sigma)=
 \max\left\{2\max_p(|c_p|p^{-\sigma})-A_c(\sigma),0\right\}.
\]

\begin{proposition}[Phase annulus and real-part projection]
\label{prop:phase-annulus}
For every $\sigma>1$,
\[
 \overline{\{F_c(\sigma+it):t\in\R\}}
 =\{z\in\C:B_c(\sigma)\le |z|\le A_c(\sigma)\}.
\]
Consequently, for every $a\in\C$,
\[
 \boxed{
 \overline{\{\Re s:F_c(s)=a,\ \Re s>1\}}
 =\{\sigma>1:B_c(\sigma)\le |a|\le A_c(\sigma)\}.
 }
\tag{2.1}
\]
The closure in the second formula is relative to $(1,\infty)$.
\end{proposition}

\begin{proof}
Unique factorization implies that the logarithms of distinct primes are linearly
independent over $\mathbb Q$.  Kronecker--Weyl approximation, first on a finite set of
primes and then with the absolutely convergent tail restored, identifies the vertical
value closure with
\[
 \left\{\sum_p|c_p|p^{-\sigma}z_p:|z_p|=1\right\}.
\]
The infinite polygon inequality gives precisely the stated closed annulus.

If a phase choice gives an equivalent Dirichlet series $G$ with $G(\sigma)=a$, then
arbitrarily large translates of $F_c$ converge locally uniformly to $G$ in
$\Re s>1$.  Hurwitz's theorem supplies $a$-points of $F_c$ whose real parts tend to
$\sigma$.  Conversely, uniform differentiability on compact substrips shows that a
limit of real parts of actual $a$-points must lie in the vertical value closure.
This proves (2.1).  This last projection step is a standard principle in analytic
almost-periodic theory; compare~\cite{JessenTornehave1945,SepulcreVidal2022}.
\end{proof}

For $c_p=\omega_p$, $|\omega_p|=1$, write $F_\omega$.  At the zero target, the largest
term is $2^{-\sigma}$, and the right edge is the unique solution
\[
 2^{-\sigma_{\rm e}}=\sum_{p\ge3}p^{-\sigma_{\rm e}},
 \qquad
 \sigma_{\rm e}=1.7795446535469941164\ldots .
\tag{2.2}
\]

\begin{corollary}[Exact zero projection and its boundary]
\label{cor:zero-projection}
For every unimodular twist,
\[
 \overline{\{\Re s:F_\omega(s)=0,\ \Re s>1\}}
 =(1,\sigma_{\rm e}].
\tag{2.3}
\]
All such twists are zero-free in $\Re s>\sigma_{\rm e}$.  This cannot be changed to a
universal closed-half-plane statement: the twist
\[
 \omega_2=1,\qquad \omega_p=-1\quad(p\ge3)
\]
satisfies $F_\omega(\sigma_{\rm e})=0$.

For the untwisted prime zeta function,
\[
 \cP(s)\ne0\qquad(\Re s\ge\sigma_{\rm e}),
\tag{2.4}
\]
although zero real parts approach $\sigma_{\rm e}$ from the left.
\end{corollary}

\begin{proof}
Only the last assertion needs comment.  Equality at the polygon boundary would force
all phases $p^{-it}$, $p\ge3$, to oppose $2^{-it}$.  In particular,
$t\log(5/3)$ and $t\log(7/3)$ would both lie in $2\pi\mathbb Z$.  If $t\ne0$, their
ratio would be rational, contradicting unique factorization; $t=0$ is impossible
because all summands are positive.
\end{proof}

\begin{remark}[Priority boundary]
\label{rem:priority-boundary}
The polygon equation (2.2) is the natural exact form of the numerically observed edge.
The general mechanism predates the 2025 conjecture, and Camargo~\cite{Camargo2025}
subsequently obtained the phase-invariant projection closure by an equivalent-series
argument.  Later sections therefore use (2.2) as an input and concentrate on genuinely
stronger density and inverse statements.
\end{remark}

\section{Smooth weighted densities and exact horizontal laws}
\label{sec:smooth-density}

Let $(U_p)_{p\in\PP}$ be independent Haar variables on the unit circle.  For
$\sigma>1$ define
\[
 X_\sigma=\sum_pU_pp^{-\sigma},
 \qquad
 Y_\sigma=-\sum_p(\log p)U_pp^{-\sigma}.
\tag{3.1}
\]
Both sums are then absolutely convergent; this random-Euler-product viewpoint goes
back to Jessen--Wintner~\cite{JessenWintner1935}.  Write $q_\sigma(x,y)$ for the joint
density of $(X_\sigma,Y_\sigma)$, when it exists, and define the derivative-weighted
value density
\[
 G(\sigma,a)=\int_\C |y|^2q_\sigma(a,y)\dd A(y).
\tag{3.2}
\]
Here and throughout, $\dd A$ is ordinary (unnormalized) planar Lebesgue measure.

\begin{lemma}[Bessel-product smoothing]
\label{lem:bessel-smoothing}
The pair $(X_\sigma,Y_\sigma)$ has a $C^\infty$ compactly supported density on
$\C^2$.  On every compact $K\Subset(1,\infty)$, the function
$q_\sigma(x,y)$ and all its derivatives depend continuously and uniformly on
$\sigma\in K$.  Consequently, $G\in C^\infty((1,\infty)\times\C)$.

Moreover, with
\[
 R(\sigma)=\cP(\sigma),
 \qquad
 r(\sigma)=\max\{2^{1-\sigma}-\cP(\sigma),0\},
\tag{3.3}
\]
one has
\[
 G(\sigma,a)>0
 \quad\Longleftrightarrow\quad
 |a|<R(\sigma)\ \text{ and }\ \bigl(r(\sigma)=0\ \text{or}\ |a|>r(\sigma)\bigr),
\tag{3.4}
\]
and $G(\sigma,a)=0$ outside the interior of the support.  When $r(\sigma)>0$,
both boundary circles have zero weighted density; when $r(\sigma)=0$, the origin is
an interior point and has positive density.
The function is radial in $a$.
\end{lemma}

\begin{proof}
For $u,v\in\C$, the joint characteristic function is
\[
 \widehat q_\sigma(u,v)
 =\prod_pJ_0\!\left(p^{-\sigma}|u-(\log p)v|\right),
\tag{3.5}
\]
where $J_0$ is the Bessel function.  Fix as many disjoint pairs of primes as desired.
For every dual direction $(u,v)$, at least one of the two linear forms belonging to
each pair is bounded below by a fixed multiple of $|(u,v)|$; cancellation can occur
for at most one logarithmic slope.  Since
\[
 |J_0(x)|\ll(1+x)^{-1/2},
\tag{3.6}
\]
the factors from these fixed pairs give arbitrarily high polynomial decay.  After any
fixed number of differentiations in $(u,v,\sigma)$, unused pairs provide the same
decay.  Thus $\widehat q_\sigma$ is a Schwartz function, uniformly for $\sigma$ in a
compact subinterval.  Fourier inversion proves the first assertion.

The support of $X_\sigma$ is the Minkowski sum of circles of radii $p^{-\sigma}$;
the polygon theorem identifies it with the closed annulus in (3.3), interpreted as a
closed disc when $r(\sigma)=0$.  At a strict interior point, finitely many phase
variables can be chosen so that the phase-to-value
map has real rank two and its $\sigma$-derivative is nonzero.  The coarea formula gives
a positive local contribution to (3.2), and convolution with the unused variables
preserves positivity.  The genuine boundary circles are critical polygonal alignments
and have zero weighted density.  Common rotation of all $U_p$ proves radiality.
\end{proof}

For $a\in\C$ and a unimodular twist $F_\omega$, let
\[
 J_a(\sigma)=
 \lim_{T\to\infty}\frac1T\int_0^T
 \log|F_\omega(\sigma+it)-a|\dd t.
\tag{3.7}
\]
The following theorem is the metric step beyond the projection law.

\begin{theorem}[Phase-blind smooth $a$-point law]
\label{thm:smooth-a-law}
For every $a\in\C$, the limit (3.7) exists, is independent of $\omega$, and is
$C^\infty$ on $(1,\infty)$.  One has
\[
 \boxed{J_a''(\sigma)=2\pi G(\sigma,a).}
\tag{3.8}
\]
If multiplicities are denoted by $m_\rho$, then for
$1<\alpha<\beta<\infty$,
\[
 \boxed{
 \sum_{\substack{F_\omega(\rho)=a\\
                   \alpha<\Re\rho<\beta\\0<\Im\rho<T}}m_\rho
 =T\int_\alpha^\beta G(\sigma,a)\dd\sigma+o(T).
 }
\tag{3.9}
\]
The error may depend on $\alpha,\beta,a$, and $\omega$.
\end{theorem}

\begin{proof}
Kronecker--Weyl equidistribution identifies the vertical translation hull of every
$F_\omega$ with the Haar product model (3.1); the coefficient phases disappear under
Haar rotation.  Classical Jessen--Tornehave theory gives the existence and convexity
of (3.7).  The Poincar\'e--Lelong identity, averaged over the vertical flow, gives in
the distributional sense
\[
 J_a''(\sigma)=2\pi\,
 \mathbb E\bigl[|Y_\sigma|^2\delta_a(X_\sigma)\bigr].
\]
The right side is exactly (3.2).  The logarithmic uniform-integrability and
mean-convergence passage is the weighted-density theorem of
Borchsenius--Jessen; the convenient modern formulation in
Lee~\cite[Sections 2--3]{Lee2014} applies after finite-prime approximation.
Lemma~\ref{lem:bessel-smoothing} removes all exceptional differentiability points.
Integrating (3.8) over $(\alpha,\beta)$ gives (3.9).
\end{proof}

\subsection{Asymptotic simplicity and marked divisor laws}

The classical Jessen count is multiplicity weighted.  For the present prime-supported
family the loss incurred by that convention is sublinear.

\begin{theorem}[Asymptotic simplicity on compact strips]
\label{thm:asymptotic-simplicity}
Fix $a\in\C$ and a compact interval $I=[\alpha,\beta]\Subset(1,\infty)$.  If
$m(\rho)$ denotes the multiplicity of an $a$-point of $F_\omega$, then
\[
 E_{\omega,a,I}(T):=
 \sum_{\substack{F_\omega(\rho)=a\\
 \Re\rho\in I,\ 0<\Im\rho<T}}(m(\rho)-1)=o(T).
\]
Consequently the multiplicity count $N^{\rm mult}$, the distinct-point count
$N^{\rm dist}$, and the simple-point count $N^{\rm simp}$ satisfy
\[
 \boxed{
 N^{\rm mult}_{\omega,a,I}(T)
 =N^{\rm dist}_{\omega,a,I}(T)+o(T)
 =N^{\rm simp}_{\omega,a,I}(T)+o(T).
 }
\]
The assertion is for fixed $a$ and $I$; no uniformity as $\alpha\downarrow1$ is
claimed.
\end{theorem}

\begin{proof}
Let $\mathbb H=\mathbb T^{\PP}$ with Haar probability and put
\[
 F_u(s)=\sum_pu_pp^{-s},\qquad (R_tu)_p=u_pp^{-it}.
\]
Then $F_\omega(s+it)=F_{R_t\omega}(s)$.  Unique factorization makes the finite
subsets of $(\log p)_p$ linearly independent over $\mathbb Q$, so this Kronecker flow
is uniquely ergodic on $\mathbb H$.

For a compact $K\Subset\{\Re s>1\}$, let $\mathcal B_{a,K}$ be the set of $u$ for
which $F_u-a$ has a multiple zero in $K$.  It is closed and Haar-null.  To see the
second assertion, fix a prime $q$ and condition on all phases except $u_q$.  With
$H(s)=\sum_{p\ne q}u_pp^{-s}$, a multiple root must satisfy
\[
 H'(s)+(\log q)(H(s)-a)=0,
 \qquad u_q=q^s(a-H(s)).
\]
The first left-hand side is not identically zero, by uniqueness of absolutely
convergent Dirichlet series.  It has finitely many zeros in $K$, and each permits at
most one value of $u_q$.  Fubini's theorem proves Haar nullity.

There is also a uniform bound, over $u\in\mathbb H$, for the number of zeros in a
fixed compact box, counted with multiplicity.  Indeed the hull is compact in the
compact-open topology and contains no function identically equal to $a$; a failure of
the bound would contradict Hurwitz's and Rouch\'e's theorems on a slightly larger box.
Choose $K=I+i[-1,1]$.  The excess-zero functional is bounded by a constant times
$\one_{\mathcal B_{a,K}}$.  Unique ergodicity and the Portmanteau inequality for this
closed null set therefore give
\[
 \int_0^T Z_K^{\rm exc}(R_t\omega)\dd t=o(T).
\]
Every multiple point of $F_\omega-a$ with ordinate in $(1,T-1)$ is seen by this
moving box for two units of time.  The two end windows contribute only $O(1)$ by the
uniform local bound.  Hence $E_{\omega,a,I}(T)=o(T)$.  Finally
$N^{\rm mult}-N^{\rm dist}=E$ and
$0\le N^{\rm dist}-N^{\rm simp}\le E$.
\end{proof}

\begin{theorem}[Marked simple $a$-point law]
\label{thm:marked-kac-rice}
Let $I\Subset(1,\infty)$ and $\Psi\in C_c(I\times\C)$.  Then
\[
 \boxed{
 \frac1T\!\sum_{\substack{F_\omega(\rho)=a,\ F_\omega'(\rho)\ne0\\
                              0<\Im\rho<T}}
 \Psi(\Re\rho,F_\omega'(\rho))
 \longrightarrow
 \int_I\!\int_\C \Psi(\sigma,y)|y|^2q_\sigma(a,y)\dd A(y)\dd\sigma .
 }
\]
In particular, the simple-point count has the same main term as (3.9).
\end{theorem}

\begin{proof}
On a unit-height box, sum the displayed mark over the simple roots of $F_u-a$.
This bounded zero functional is continuous except when a root lies on the boundary or
is multiple.  The second exceptional set is Haar-null by the conditioning argument in
the preceding proof.  The first is Haar-null as well: the expected number of roots in
a shrinking collar is $O(\delta)$ by the Kac--Rice formula, and one lets
$\delta\downarrow0$.  Thus the functional is Haar-Riemann integrable, and unique
ergodicity transfers its Haar mean to every twist.  Complex Kac--Rice evaluates that
mean as
\[
 \int_I\mathbb E\!\left[
  \Psi(\sigma,Y_\sigma)|Y_\sigma|^2\delta_a(X_\sigma)
 \right]\dd\sigma,
\]
which is the asserted integral.  The unmarked statement follows by a compact cutoff
in the uniformly bounded derivative range and Theorem~\ref{thm:asymptotic-simplicity}.
\end{proof}

\begin{corollary}[Near-critical slopes and large residues]
\label{cor:near-multiple}
For fixed $a$ and $I\Subset(1,\infty)$,
\[
 \lim_{T\to\infty}\frac1T
 \#\{\rho:F_\omega(\rho)=a,\ \Re\rho\in I,\ 0<\Im\rho<T,
                   \ 0<|F_\omega'(\rho)|\le\eps\}
 =\int_I\int_{|y|\le\eps}|y|^2q_\sigma(a,y)\dd A(y)\dd\sigma
\]
and, as $\eps\downarrow0$, the right side equals
\[
 \frac{\pi\eps^4}{2}\int_Iq_\sigma(a,0)\dd\sigma+O_{I,a}(\eps^6).
\]
If $\mathcal D_z=(1-zF_\omega)^{-1}$, then in the same strip
\[
 \frac1T\#\{\rho:\rho\text{ is a simple pole of }\mathcal D_z,
                   \ |\operatorname {Res}_\rho\mathcal D_z|>R\}
 =\frac{\pi}{2|z|^4R^4}\int_Iq_\sigma(z^{-1},0)\dd\sigma
  +O_{I,z}(R^{-6})+o_T(1).
\]
The constants are not uniform as the target tends to zero.
\end{corollary}

\begin{corollary}[Complete simple valence in thin strips]
\label{cor:thin-strip-valence}
Let $S$ be an infinite prime support and suppose, for a fixed $m\ge0$, that
\[
 \sum_{p\in S}(\log p)^mp^{-\sigma}\longrightarrow\infty
 \qquad(\sigma\downarrow1).
\]
For every $a\in\C$, every $\delta>0$, and every fixed unimodular twist, there is a
compact interval $I\Subset(1,1+\delta)$ and a constant $\kappa>0$ such that
\[
 N^{\rm simp}_{m,a}(I;T)=\kappa T+o(T).
\]
Hence every value is assumed infinitely often, at simple points, in every boundary
strip.  The hypothesis holds for all primes, cofinite prime supports, and positive
Chebotarev-density supports.
\end{corollary}

\begin{proof}
Near $1$, the total polygonal radius diverges whereas its largest summand remains
bounded.  Every fixed target therefore lies in the strict interior of the value
annulus on some compact interval $I\Subset(1,1+\delta)$.  The weighted density is
positive there, and Theorems~\ref{thm:smooth-a-law} and
\ref{thm:asymptotic-simplicity} give the result.
\end{proof}

\section{The 2025 conjectures and the endpoint at \texorpdfstring{$\Re s=1$}{Re s=1}}
\label{sec:bcs-endpoint}

Put
\[
 \rho(\sigma)=G(\sigma,0).
\tag{4.1}
\]
By Lemma~\ref{lem:bessel-smoothing} and (2.2),
\[
 \rho(\sigma)>0\quad(1<\sigma<\sigma_{\rm e}),
 \qquad
 \rho(\sigma)=0\quad(\sigma\ge\sigma_{\rm e}).
\tag{4.2}
\]
Since $G$ is smooth through the boundary, the edge is infinitely soft:
\[
 \rho^{(j)}(\sigma_{\rm e})=0\qquad(j\ge0).
\tag{4.3}
\]

\subsection{Exact density, edge convergence, and height uniformity}

For $1<\gamma<\sigma_{\rm e}$ define
\[
 c_\gamma=\int_\gamma^{\sigma_{\rm e}}\rho(\sigma)\dd\sigma.
\tag{4.4}
\]

\begin{theorem}[Exact simple form of the BCS metric conjectures]
\label{thm:bcs}
For every unimodular twist and $1<\gamma<\sigma_{\rm e}$,
\[
 M^{\rm mult}_{\omega,\gamma}(T):=
 \sum_{\substack{F_\omega(\rho)=0\\
                  \Re\rho>\gamma\\0<\Im\rho<T}}m_\rho
 =c_\gamma T+o(T),
\tag{4.5}
\]
and
\[
 \boxed{
 M^{\rm dist}_{\omega,\gamma}(T)
 =M^{\rm simp}_{\omega,\gamma}(T)
 =c_\gamma T+o(T).
 }
\]
where $c_\gamma>0$.  Moreover, for $0\le x\le1$,
\[
 \frac{M^{\bullet}_{\omega,\gamma}(xT)}
      {M^{\bullet}_{\omega,\gamma}(T)}\longrightarrow x,
 \qquad \bullet\in\{\mathrm{mult},\mathrm{dist},\mathrm{simp}\}.
\tag{4.6}
\]
For the untwisted prime zeta function, if
\[
 \Sigma(T)=\sup\{\Re\rho:\cP(\rho)=0,\ 0<\Im\rho<T\},
\]
then
\[
 \Sigma(T)\longrightarrow\sigma_{\rm e}.
\tag{4.7}
\]
Thus (4.7) recovers BCS Conjecture~1 from the earlier general projection theory.
Equations (4.5)--(4.6) prove the literal distinct-point Conjecture~2 and strengthen
Conjecture~3 from two-sided linear bounds to an exact asymptotic, in their stated range
$1<\gamma<\sigma_{\rm e}$.
\end{theorem}

\begin{proof}
Choose any $\beta>\sigma_{\rm e}$ and apply (3.9) to $(\gamma,\beta)$; there
are no zeros with real part in $(\sigma_{\rm e},\beta)$, and
positivity follows from (4.2).  Theorem~\ref{thm:asymptotic-simplicity} transfers the
same main term to the distinct and simple counts.  Replacing $T$ by $xT$ proves
(4.6).  For every
$\eps>0$, (4.5) supplies zeros with real part in
$(\sigma_{\rm e}-\eps,\sigma_{\rm e})$ for all sufficiently large $T$, while
Corollary~\ref{cor:zero-projection} supplies the upper bound.  This proves (4.7).
\end{proof}

The result has a useful two-variable form.  For compactly supported continuous
$\Phi$ on $(1,\sigma_{\rm e})\times(0,1)$,
\[
 \boxed{
 \frac1T\sum_{\substack{F_\omega(\rho)=0\\0<\Im\rho<T}}
 m_\rho\,
 \Phi\!\left(\Re\rho,\frac{\Im\rho}{T}\right)
 \longrightarrow
 \int_1^{\sigma_{\rm e}}\int_0^1
 \Phi(\sigma,u)\rho(\sigma)\dd u\dd\sigma.
 }
\tag{4.8}
\]
Indeed, (4.8) holds first for rectangles by subtracting (4.5) at two heights and
then for all such $\Phi$ by approximation.  Thus the real part and normalized height
are asymptotically independent at the scale detected by the limiting divisor.  Also,
every fixed vertical line carries $o(T)$ zeros because the limiting horizontal measure
has no atoms.  The same limit holds with distinct or simple points, by
Theorem~\ref{thm:asymptotic-simplicity}.

\subsection{Syndetic recurrence of simple points}

The density law gives many points on average.  Almost periodicity upgrades existence
in any occupied band to a bounded-window statement.

\begin{theorem}[Bounded-window recurrence]
\label{thm:syndetic-recurrence}
Let $I\Subset(1,\infty)$ and suppose
$\int_I G(\sigma,a)\dd\sigma>0$.  For every fixed unimodular twist there is
$H=H(I,a,\omega)>0$ such that every sufficiently high interval $[u,u+H]$ contains
the ordinate of a simple $a$-point with real part in $I$.
\end{theorem}

\begin{proof}
Theorems~\ref{thm:smooth-a-law} and \ref{thm:asymptotic-simplicity} provide one
simple point in the interior of the band.  Choose a small disc about it on whose
boundary $F_\omega-a$ is nonzero.  The set of uniform almost periods on this compact
disc is relatively dense, because it is the return set of a minimal translation on
the compact prime-phase hull.  Rouch\'e's theorem transports the simple zero to every
corresponding translate.  Shrinking the disc keeps the real part inside $I$ and yields
the asserted bound $H$.
\end{proof}

\begin{corollary}[Syndetic approach to the zero edge]
\label{cor:edge-recurrence}
For every $\eta>0$ and every fixed unimodular twist, the strip
\[
 \sigma_{\rm e}-\eta<\Re s<\sigma_{\rm e}
\]
contains simple zeros whose ordinates have bounded gaps.  Thus the right edge is
approached in every sufficiently high bounded-height window, not merely along a sparse
sequence.
\end{corollary}

\begin{proof}
Choose $I\Subset(\max\{1,\sigma_{\rm e}-\eta\},\sigma_{\rm e})$.
Then $G(\sigma,0)>0$ on $I$, and Theorem~\ref{thm:syndetic-recurrence} applies.
\end{proof}

For $\sigma\ge\sigma_{\rm e}$, the $p=2$ term dominates and
\[
 J_0(\sigma)=-\sigma\log2.
\tag{4.9}
\]
Consequently,
\[
 c_\gamma=\frac{-\log2-J_0'(\gamma)}{2\pi}.
\tag{4.10}
\]

\subsection{Exponential transfer across the abscissa of convergence}

The remaining endpoint is special to $\omega_p\equiv1$.  In $\Re s>1$,
\[
 e^{\cP(s)}=\zeta(s)H(s),
 \qquad
 H(s)=\exp\!\left(-\sum_{k\ge2}\frac{\cP(ks)}k\right).
\tag{4.11}
\]
The function $H$ is holomorphic and nonvanishing in $\Re s>\tfrac12$.

\begin{lemma}[Exponential boundary majorant]
\label{lem:boundary-majorant}
Fix $a\in\C$ and set
\[
 E(s)=\zeta(s)H(s),\qquad Q_a(s)=E(s)-e^a.
\]
Then $Q_a$ is meromorphic in $\Re s>\tfrac12$, with only the pole at $s=1$.
Its multiplicity-weighted zero divisor has a locally bounded horizontal density near
$\sigma=1$.  In particular, for sufficiently small fixed $\delta>0$,
\[
 N_{Q_a}(1,1+\delta;T)\ll_a\delta T+o(T).
\tag{4.12}
\]
\end{lemma}

\begin{proof}
For $\Re s>1$,
\[
 E(s)=\prod_pe^{p^{-s}}=\sum_{n\ge1}\frac{b_n}{n^s},
 \qquad
 b_n=\prod_p\frac1{\nu_p(n)!}.
\tag{4.13}
\]
For $\sigma>\tfrac12$,
\[
 \sum_{n\ge1}\frac{b_n^2}{n^{2\sigma}}
 =\prod_p I_0(2p^{-\sigma})<\infty,
\tag{4.14}
\]
and the same is true after inserting any fixed logarithmic power.  Hence the Dirichlet
polynomials of (4.13) converge to the meromorphic continuation in the required local
mean-square sense.  The generalized Jessen transfer theorem of
Borchsenius--Jessen~\cite{BorchseniusJessen1948}, in the form used by
Lee~\cite[Theorem 3.1]{Lee2014}, applies after excising the isolated pole.

For completeness, the random series (3.1) are well defined in $L^2$ for
$\sigma>\tfrac12$.  The proof of Lemma~\ref{lem:bessel-smoothing}, now using prime
blocks rather than absolute support, gives a Schwartz joint density and a weighted
density $G_*(\sigma,z)$ smooth on this larger half-plane.  The random envelope is
$e^{X_\sigma}$, with derivative $e^{X_\sigma}Y_\sigma$.  At a preimage
$X_\sigma=a+2\pi ik$, the Jacobian of the exponential cancels the factor
$|e^{X_\sigma}|^2$.  The zero density of $Q_a$ is therefore
\[
 \widetilde G_a(\sigma)=
 \sum_{k\in\mathbb Z}G_*(\sigma,a+2\pi ik).
\tag{4.15}
\]
The series and all its derivatives converge locally uniformly because $G_*$ is
Schwartz in the value variable.  It is thus bounded near $\sigma=1$, and integration
over $(1,1+\delta)$ proves (4.12).
\end{proof}

\begin{theorem}[Exact endpoint $a$-point law]
\label{thm:endpoint}
For every fixed $a\in\C$,
\[
 M_a^{\rm mult}(T):=
 \sum_{\substack{\cP(\rho)=a\\\Re\rho>1\\0<\Im\rho<T}}m_\rho
 =C_aT+o(T),
 \qquad
 C_a=\int_1^\infty G(\sigma,a)\dd\sigma,
\tag{4.16}
\]
where $0<C_a<\infty$.  The normalized ordinates are uniformly distributed:
\[
 M_a^{\rm dist}(T)=M_a^{\rm simp}(T)=C_aT+o(T),
 \qquad
 \frac{M_a^\bullet(xT)}{M_a^\bullet(T)}\longrightarrow x
 \quad(\bullet\in\{\mathrm{mult},\mathrm{dist},\mathrm{simp}\}).
\tag{4.17}
\]
In particular,
\[
 \boxed{
 M_1(T):=
 \sum_{\substack{\cP(\rho)=0\\\Re\rho>1\\0<\Im\rho<T}}m_\rho
 =c_1T+o(T),
 \qquad
 c_1=\int_1^{\sigma_{\rm e}}\rho(\sigma)\dd\sigma>0.
 }
\tag{4.18}
\]
and the same asymptotic holds for the distinct and simple zero counts.
This strengthens the separate endpoint estimate $M_1(T)=O(T)$ in
\cite{BCS2025}.
\end{theorem}

\begin{proof}
For fixed $\delta>0$, Theorem~\ref{thm:smooth-a-law} gives the multiplicity asymptotic in
$\Re s>1+\delta$.  Every zero of $\cP-a$ is a zero of $Q_a=e^{\cP}-e^a$ with the same
multiplicity.  Lemma~\ref{lem:boundary-majorant} therefore bounds the omitted layer by
$O_a(\delta T)+o(T)$.  First let $T\to\infty$ and then $\delta\downarrow0$.
This squeezes $M_a^{\rm mult}(T)/T$ to the integral in (4.16).

For $a\ne0$, the integral is finite because no $a$-point is possible once
$\cP(\sigma)<|a|$; for $a=0$, its support ends at $\sigma_{\rm e}$.  Positivity follows
because, as $\sigma\downarrow1$, the strict polygon annulus contains every fixed
target on some nonempty interval.  On $\Re s\ge1+\delta$,
Theorem~\ref{thm:asymptotic-simplicity} makes the excess multiplicity $o(T)$; the
omitted boundary layer is again at most $O_a(\delta T)+o(T)$.  Letting
$\delta\downarrow0$ proves the distinct and simple formulas.  Height uniformity follows
by replacing $T$ by $xT$.
\end{proof}

\begin{remark}[The exponential is a majorant, not an equivalence]
\label{rem:aliases}
It is essential that
\[
 \cP(s)=a\quad\Longrightarrow\quad e^{\cP(s)}=e^a
\]
is only one direction.  The reverse implication gives
$\cP(s)=a+2\pi ik$ for some $k\in\mathbb Z$; indeed,
\[
 \divisor_0(E-e^a)
 =\sum_{k\in\mathbb Z}\divisor_0(\cP-a-2\pi ik)
 \qquad(\Re s>1).
\tag{4.19}
\]
Thus the envelope is used only for a boundary-layer upper bound.  No analogous
factorization is available for an arbitrary twist, so Theorem~\ref{thm:endpoint} is not
asserted for $F_\omega$.
\end{remark}

\subsection{Quantized small-target escape}

The target $a=0$ is singular from the viewpoint of the far right.  Nonzero targets
have a finite right edge, whereas zero points of the hull retain a terminal contribution
from the least supported prime.  The discrepancy is a quantized atom at infinity.

\begin{theorem}[Quantized escape of $a$-point mass]
\label{thm:mass-escape}
Let $S$ be an infinite prime support for which the smooth density conclusion of
Lemma~\ref{lem:bessel-smoothing} holds, let $p_0=\min S$, and put
\[
 F_{S,\omega}(s)=\sum_{p\in S}\omega_pp^{-s}.
\]
For fixed $\sigma_0>1$, write
\[
 C_{S,a}(\sigma_0)=\int_{\sigma_0}^{\infty}G_S(\sigma,a)\dd\sigma.
\]
Then
\[
 \boxed{
 \lim_{a\to0,\ a\ne0}C_{S,a}(\sigma_0)
 =C_{S,0}(\sigma_0)+\frac{\log p_0}{2\pi}.
 }
\tag{4.20}
\]
More precisely, on the one-point compactification $[\sigma_0,\infty]$,
\[
 \boxed{
 G_S(\sigma,a)\dd\sigma\ \Longrightarrow\
 G_S(\sigma,0)\dd\sigma+\frac{\log p_0}{2\pi}\delta_\infty .
 }
\tag{4.21}
\]
Thus the jump recovers the least support prime:
\[
 p_0=\exp\!\left(2\pi
  \left[\lim_{a\to0,\ a\ne0}C_{S,a}(\sigma_0)-C_{S,0}(\sigma_0)\right]
 \right).
\tag{4.22}
\]
For the full prime zeta function the escaping mass is $\log2/(2\pi)$.
\end{theorem}

\begin{proof}
Integrating $J_a''=2\pi G_S$ to $+\infty$ gives
\[
 2\pi C_{S,a}(\sigma_0)=-J_a'(\sigma_0)\quad(a\ne0),
 \qquad
 2\pi C_{S,0}(\sigma_0)=-\log p_0-J_0'(\sigma_0).
\]
Indeed, $J_a(\sigma)=\log|a|$ once the entire value support lies inside the circle
of radius $|a|$, whereas eventual $p_0$-dominance gives
$J_0'(\sigma)=-\log p_0$.  If $f_\sigma$ is the marginal density of $X_\sigma$,
then
\[
 J_a'(\sigma_0)=
 \int_\C\log|x-a|\,\partial_\sigma f_{\sigma_0}(x)\dd A(x).
\]
Smooth compact support of $f_\sigma$ and local integrability of the logarithmic kernel
give $J_a'(\sigma_0)\to J_0'(\sigma_0)$.  This proves (4.20).  Local uniform
convergence $G_S(\sigma,a)\to G_S(\sigma,0)$ on finite $\sigma$-intervals, together
with convergence of total masses, proves (4.21).
\end{proof}

\begin{proposition}[The escaping simple lattice]
\label{prop:escape-lattice}
Let $p_1$ be the next prime in $S$ and
$\eta=\log p_1/\log p_0-1>0$.  As $a\to0$, the remote $a$-points are simple and,
uniformly in their integer label $k$,
\[
 \rho_{a,k}=\frac{\log(1/|a|)}{\log p_0}
 +i\,\frac{\arg\omega_{p_0}-\arg a+2\pi k}{\log p_0}
 +O(|a|^\eta),
\tag{4.23}
\]
\[
 F_{S,\omega}'(\rho_{a,k})=-(\log p_0)a+O(|a|^{1+\eta}).
\tag{4.24}
\]
Their positive-height density is $\log p_0/(2\pi)$, exactly the atom in (4.21).
\end{proposition}

\begin{proof}
The displayed leading points are the simple solutions of
$\omega_{p_0}p_0^{-s}=a$.  On equal small circles around them, the sum over
$p\ge p_1$ is $O(|a|^{1+\eta})$.  Rouch\'e's theorem and one Newton step give
(4.23)--(4.24).  Consecutive leading ordinates differ by $2\pi/\log p_0$.
\end{proof}

\section{Prime spectral tomography from nonzero \texorpdfstring{$a$}{a}-points}
\label{sec:tomography}

We now reverse the preceding construction.  For a bounded coefficient system $c$, let
\[
 X_{c,\sigma}=\sum_p|c_p|p^{-\sigma}U_p,
 \qquad
 \mu_{c,\sigma}=\operatorname{Law}(X_{c,\sigma}),
\tag{5.1}
\]
and define
\[
 J_{c,a}(\sigma)=\mathbb E\log|X_{c,\sigma}-a|,
 \qquad
 \nu_{c,a}=\frac1{2\pi}J_{c,a}''
\tag{5.2}
\]
on $(1,\infty)$.  The measure $\nu_{c,a}$ is the limiting horizontal
multiplicity-weighted $a$-point measure.

\begin{theorem}[All-$a$ tomographic equivalence]
\label{thm:all-a-tomography}
Let $c=(c_p)$ and $d=(d_p)$ be bounded coefficient systems.  The following are
equivalent.
\begin{enumerate}[label=\textup{(\roman*)}]
\item $|c_p|=|d_p|$ for every prime $p$.
\item $\mu_{c,\sigma}=\mu_{d,\sigma}$ for every $\sigma>1$.
\item $J_{c,a}(\sigma)=J_{d,a}(\sigma)$ for every $\sigma>1$ and every $a\ne0$.
\item $\nu_{c,a}=\nu_{d,a}$ on $(1,\infty)$ for every $a\ne0$, with the common
normalization
\[
 J_{c,a}(\sigma),J_{d,a}(\sigma)\longrightarrow\log|a|,
 \qquad
 J_{c,a}'(\sigma),J_{d,a}'(\sigma)\longrightarrow0
\quad(\sigma\to\infty).
\]
\end{enumerate}
Thus the family of all nonzero $a$-point horizontal measures is a complete invariant
of the coefficient moduli and is maximally insensitive to their phases.
\end{theorem}

\begin{proof}
Condition (i) gives equality in law in (5.1), and hence (ii)$\Rightarrow$(iii)
as well.  Distributional differentiation gives (iii)$\Rightarrow$(iv).

Conversely, put $r=|a|$.  Rotational invariance and Jensen's circle formula yield
\[
 J_{c,a}(\sigma)=
 \mathbb E\log\max\{r,|X_{c,\sigma}|\}.
\tag{5.3}
\]
Once $A_c(\sigma)<r$, this equals $\log r$.  The normalization at infinity therefore
removes both affine ambiguities in the second derivative and gives
\[
 \boxed{
 J_{c,a}(\sigma)=\log r+
 2\pi\int_{(\sigma,\infty)}(u-\sigma)\nu_{c,a}(\dd u).
 }
\tag{5.4}
\]
Thus (iv)$\Rightarrow$(iii).

For fixed $\sigma$, logarithmic potential theory gives
\[
 \Delta_aJ_{c,a}(\sigma)=2\pi\mu_{c,\sigma}
\tag{5.5}
\]
in distributions.  Equality for all $a\ne0$, together with total mass one and the
normalization at infinity, also determines any mass at the origin; hence
(iii)$\Rightarrow$(ii).  Finally, equality of the laws gives equality of variances:
\[
 \sum_p|c_p|^2p^{-2\sigma}
 =\mathbb E|X_{c,\sigma}|^2
 =\mathbb E|X_{d,\sigma}|^2
 =\sum_p|d_p|^2p^{-2\sigma}.
\tag{5.6}
\]
Uniqueness of absolutely convergent Dirichlet series proves (i).
\end{proof}

\begin{corollary}[Explicit recovery of a vertical value law]
\label{cor:radial-recovery}
Writing $R_{c,\sigma}=|X_{c,\sigma}|$, the one-sided radial derivatives satisfy
\[
 r\,\partial_r^-J_{c,r}(\sigma)=\mathbb P(R_{c,\sigma}<r),
 \qquad
 r\,\partial_r^+J_{c,r}(\sigma)=\mathbb P(R_{c,\sigma}\le r).
\tag{5.7}
\]
Thus the all-$a$ horizontal divisor recovers the full vertical value law at every
$\sigma>1$.
\end{corollary}

\begin{proof}
Differentiate (5.3) under the Stieltjes integral.  The two one-sided derivatives record
a possible atom on the circle of radius $r$.  Rotational invariance then reconstructs
the planar law from its radial distribution.
\end{proof}

\subsection{Target-plane conservation and compressed tomography}

The horizontal density can also be integrated in the target variable.  This loses most
of the value distribution but leaves an exact spectral trace.

\begin{proposition}[Target-plane conservation]
\label{prop:target-mass}
For a prime support $S$, put $P_S(s)=\sum_{p\in S}p^{-s}$ and let $G_S$ be the
corresponding weighted density.  With ordinary planar Lebesgue measure $\dd A$,
\[
 \boxed{
 \int_\C G_S(\sigma,a)\dd A(a)
 =\sum_{p\in S}(\log p)^2p^{-2\sigma}=P_S''(2\sigma).
 }
\tag{5.7a}
\]
More generally, for $1<\alpha<\beta\le\infty$,
\[
 \int_\C\int_\alpha^\beta G_S(\sigma,a)\dd\sigma\dd A(a)
 =\frac12\sum_{p\in S}(\log p)
       \bigl(p^{-2\alpha}-p^{-2\beta}\bigr).
\tag{5.7b}
\]
In particular, if $C_S(a)=\int_1^\infty G_S(\sigma,a)\dd\sigma$, then
\[
 \boxed{
 \int_\C C_S(a)\dd A(a)
 =\frac12\sum_{p\in S}\frac{\log p}{p^2}=-\frac12P_S'(2).
 }
\tag{5.7c}
\]
For the $m$th derivative tower the right sides become, respectively,
$P_S^{(2m+2)}(2\sigma)$ and $-P_S^{(2m+1)}(2)/2$.
\end{proposition}

\begin{proof}
Tonelli's theorem and Haar orthogonality give
\[
 \int_\C G_S(\sigma,a)\dd A(a)
 =\mathbb E|Y_{S,\sigma}|^2
 =\sum_{p,q\in S}\frac{(\log p)(\log q)}{(pq)^\sigma}
     \mathbb E(U_p\overline{U_q})
 =P_S''(2\sigma).
\]
Integrating in $\sigma$ proves the remaining formulas.  In the $m$th derivative
case, $Y$ has coefficient $(\log p)^{m+1}p^{-\sigma}$.
\end{proof}

\begin{corollary}[One scalar hears the modulus radical]
\label{cor:modulus-tomography}
For an integer $q\ge1$, define
\[
 \mathfrak M(q)=\frac12\sum_{p\nmid q}\frac{\log p}{p^2}.
\]
Then
\[
 \boxed{
 \mathfrak M(q_1)=\mathfrak M(q_2)
 \quad\Longleftrightarrow\quad
 \operatorname{rad}(q_1)=\operatorname{rad}(q_2).
 }
\tag{5.7d}
\]
Equivalently, for Dirichlet-character prime series, the complete horizontal density,
the full right-edge curve, and this single target-integrated mass all classify exactly
the set of primes dividing the modulus.  They do not recover character phases or prime
powers in the modulus.
\end{corollary}

\begin{proof}
Only injectivity of the scalar needs proof.  Equality gives a finite relation
$\sum_{p\in E}\eps_p(\log p)/p^2=0$ with $\eps_p\in\{-1,1\}$.  Multiplication by
$D=\prod_{p\in E}p^2$ and exponentiation give
$\prod_{p\in E}p^{\eps_pD/p^2}=1$.  Unique factorization forces $E=\varnothing$.
The other equivalences follow from Theorem~\ref{thm:all-a-tomography} and
Theorem~\ref{thm:right-edge-tomography}.
\end{proof}

\subsection{One vertical line already contains every prime}

The preceding theorem used all vertical lines because it allowed arbitrary weights.
For the unit-coefficient prime series, one line suffices.

\begin{theorem}[One-line recovery of the prime spectrum]
\label{thm:one-line-recovery}
Fix $\sigma>1$.  The vertical limiting value law of
\[
 F_\omega(\sigma+it)=\sum_p\omega_pp^{-\sigma-it}
\]
determines the complete prime sequence $2,3,5,\ldots$.
\end{theorem}

\begin{proof}
For $X_\sigma$ as in (3.1), the law determines
\[
 \phi_\sigma(t)=\mathbb E e^{it\Re X_\sigma}
 =\mathbb E J_0(t|X_\sigma|)
 =\prod_pJ_0(tp^{-\sigma}).
\tag{5.8}
\]
If $j_{0,\ell}$ are the positive zeros of $J_0$, then near the origin
\[
 \log J_0(z)=-\sum_{m\ge1}b_mz^{2m},
 \qquad
 b_m=\frac1m\sum_{\ell\ge1}j_{0,\ell}^{-2m}>0.
\tag{5.9}
\]
It follows that
\[
 \log\phi_\sigma(t)
 =-\sum_{m\ge1}b_m\cP(2m\sigma)t^{2m}.
\tag{5.10}
\]
Hence the law determines every power sum
\[
 S_m=\sum_p(p^{-2\sigma})^m.
\]
Arrange $x_1>x_2>\cdots$ with $x_k=p_k^{-2\sigma}$.  Successive root peeling gives
\[
 x_1=\lim_{m\to\infty}S_m^{1/m},
\qquad
 x_{k+1}=\lim_{m\to\infty}
 \left(S_m-\sum_{j\le k}x_j^m\right)^{1/m},
\tag{5.11}
\]
and finally $p_k=x_k^{-1/(2\sigma)}$.
\end{proof}

\begin{remark}
For general $c_p$, one fixed line recovers only the unordered multiset
$\{|c_p|p^{-\sigma}:p\in\PP\}$.  Prime labels are recovered by using all $\sigma$,
as in Theorem~\ref{thm:all-a-tomography}, or by the edge curve below.
\end{remark}

\subsection{Minimal right-edge tomography}

For $a\ne0$ define
\[
 E_c(a)=
 \sup\overline{\{\Re s:F_c(s)=a,\ \Re s>1\}}.
\tag{5.12}
\]

\begin{theorem}[Right-edge tomography]
\label{thm:right-edge-tomography}
Assume that at least two coefficients $c_p$ are nonzero.  For every
$r=|a|$ in the range of $A_c$, one has
\[
 \boxed{E_c(a)=A_c^{-1}(r).}
\tag{5.13}
\]
Consequently, the following data determine one another:
\[
 \boxed{
 (|c_p|)_{p\in\PP}
 \quad\Longleftrightarrow\quad
 A_c(\sigma)
 \quad\Longleftrightarrow\quad
 E_c(r).
 }
\tag{5.14}
\]
The edge depends only on $|a|$ and the coefficient moduli.
\end{theorem}

\begin{proof}
The triangle inequality forbids an $a$-point whenever
$A_c(\sigma)<r$, proving the upper bound.  At
$\sigma_0=A_c^{-1}(r)$, align every hull phase with $a/r$.  The resulting equivalent
series takes the value $a$ at $\sigma_0$, and the zero of its difference from $a$ is
simple because
\[
 \frac{\dd}{\dd\sigma}
 \sum_p|c_p|p^{-\sigma}\ne0.
\]
Kronecker approximation and Hurwitz's theorem therefore produce actual $a$-points with
real parts approaching $\sigma_0$.  This proves (5.13).  Inverting the edge gives
$A_c$; uniqueness of positive Dirichlet series then recovers $|c_p|$ prime by prime.
\end{proof}

\begin{corollary}[Can one hear the primes?]
\label{cor:hear-primes}
For every unimodular twist and every $r>0$,
\[
 \boxed{E_\omega(r)=\cP^{-1}(r).}
\tag{5.15}
\]
Thus $E_\omega^{-1}(\sigma)=\cP(\sigma)$, and the primes are recovered recursively by
\[
 p_1=\lim_{\sigma\to\infty}\cP(\sigma)^{-1/\sigma},
\tag{5.16}
\]
\[
 p_{k+1}=\lim_{\sigma\to\infty}
 \left(\cP(\sigma)-\sum_{j\le k}p_j^{-\sigma}\right)^{-1/\sigma}.
\tag{5.17}
\]
For the untwisted function and positive target $a=r$, the right endpoint is the actual
real point $s=\cP^{-1}(r)$; for a general twist it is only asserted as a supremum of a
closure.
\end{corollary}

\begin{corollary}[Sparse edge sampling]
\label{cor:sparse-edge}
Suppose $E_c(r_n)=E_d(r_n)$ for distinct positive $r_n$ converging to an interior point
of the common edge range.  Then $|c_p|=|d_p|$ for every prime $p$.  Thus even an edge
sample with one interior accumulation point is inverse complete.
\end{corollary}

\begin{proof}
At the common edge values $\sigma_n$ one has
$A_c(\sigma_n)=r_n=A_d(\sigma_n)$.  The $\sigma_n$ accumulate in $\Re s>1$; the
identity theorem gives $A_c=A_d$, followed by Dirichlet-series uniqueness.
\end{proof}

\begin{remark}[What zero data cannot hear]
\label{rem:zero-not-complete}
The zero-target measure alone is not inverse complete.  Both $2^{-s}$ and
$2^{-s}+3^{-s}$ have no zeros in $\Re s>1$, although their modulus spectra differ.
Moreover, every limiting datum above is phase blind: it recovers $|c_p|$ but never the
individual phases $c_p/|c_p|$, and it does not identify the actual zero ordinates of
two twists.
\end{remark}


\section{Derivative edges and value caustics}
\label{sec:derivative-caustics}

For $m\ge0$ and a unimodular sequence $\omega$, put
\[
 \cP_{m,\omega}(s)=
 \sum_p\omega_p(\log p)^mp^{-s},
 \qquad \Re s>1.
\tag{6.1}
\]
Thus $\cP_{m,1}=(-1)^m\cP^{(m)}$.  Define
\[
 a_{m,p}(\sigma)=(\log p)^mp^{-\sigma},
 \qquad
 A_m(\sigma)=\sum_pa_{m,p}(\sigma),
\tag{6.2}
\]
and
\[
 R_m(\sigma)=
 \max\left\{0,\max_p\left(a_{m,p}(\sigma)-
 \sum_{q\ne p}a_{m,q}(\sigma)\right)\right\}.
\tag{6.3}
\]

\begin{lemma}[Fixed-order logarithmic-weight transfer]
\label{lem:fixed-order-transfer}
Fix $m\ge0$ and define
\[
 X_{m,\sigma}=\sum_pU_p(\log p)^mp^{-\sigma},\qquad
 Y_{m,\sigma}=-\sum_pU_p(\log p)^{m+1}p^{-\sigma}.
\]
The pair $(X_{m,\sigma},Y_{m,\sigma})$ has a $C^\infty$ compactly supported
density $q_{m,\sigma}$, locally uniformly with all derivatives for $\sigma>1$.
If
\[
 G_m(\sigma,a)=\int_{\C}|y|^2q_{m,\sigma}(a,y)\dd A(y),
\]
then the phase-annulus projection, the Jessen identity
\[
 J_{m,a}''(\sigma)=2\pi G_m(\sigma,a),
\]
the horizontal counting law, asymptotic simplicity, and the marked simple-point
law of Sections~2--3 all hold for $\cP_{m,\omega}$, for each fixed $m$.
All assertions are locally uniform in $\sigma$, but no uniformity in
$m\to\infty$ is asserted.
\end{lemma}

\begin{proof}
For every $j\ge0$ and $\delta>0$,
\[
 \sum_p(\log p)^{m+j}p^{-1-\delta}<\infty.
\]
Hence the series and every fixed $s$-derivative converge locally uniformly in
$\Re s>1$.  This is precisely the tail input used in the finite-prime
Kronecker--Weyl argument and proves the phase-annulus projection for the present
weights.

The joint characteristic function is
\[
 \prod_pJ_0\!\left((\log p)^mp^{-\sigma}
                 |u-(\log p)v|\right).
\]
The fixed-pair argument of Lemma~\ref{lem:bessel-smoothing} supplies arbitrary
polynomial decay.  Differentiation in $u,v,$ or $\sigma$ inserts only further
powers of $\log p$, which are absorbed by the displayed locally uniform
summability.  Fourier inversion therefore gives the asserted smooth density.
The Jessen and Kac--Rice arguments then apply without change.

For completeness, the simplicity proof also survives verbatim.  The weighted
translation hull is compact in the compact-open topology.  After conditioning on
all phases except the phase at a prime $q$, a multiple $a$-point must satisfy
\[
 H'(s)+(\log q)(H(s)-a)=0,
\]
and uniqueness of absolutely convergent Dirichlet series shows that this function
is not identically zero.  The conditional exceptional phase set is therefore
finite on compact boxes and Haar-null.  The compact-hull zero bound, unique
ergodicity, and the marked Kac--Rice argument from Section~3 complete the transfer.
\end{proof}

Lemma~\ref{lem:fixed-order-transfer} and
Proposition~\ref{prop:phase-annulus} give the exact global law
\[
 \overline{\{\Re s:\cP_{m,\omega}(s)=a\}}
 =\{\sigma>1:R_m(\sigma)\le |a|\le A_m(\sigma)\}.
\tag{6.4}
\]
For $m\ge1$, the largest term can change with $\sigma$; (6.4), rather than a single
interval slogan, is the safe global statement.

Set
\[
 D_m(\sigma)=
 (\log2)^m2^{-\sigma}-\sum_{p\ge3}(\log p)^mp^{-\sigma}.
\tag{6.5}
\]

\begin{theorem}[Derivative zero-edge tower]
\label{thm:derivative-edge}
For every $m\ge0$, there is a unique $\sigma_m>1$ such that
\[
 \boxed{
 \sum_{p\ge3}
 \left(\frac{\log p}{\log2}\right)^m
 \left(\frac2p\right)^{\sigma_m}=1.
 }
\tag{6.6}
\]
It is the exact right endpoint of every unimodular derivative twist:
\[
 \sup\overline{\{\Re s:\cP_{m,\omega}(s)=0,\ \Re s>1\}}
 =\sigma_m.
\tag{6.7}
\]
The untwisted $\cP_{m,1}$ has no actual zero on
$\Re s=\sigma_m$, although zero real parts approach it.  An adversarial twist can
attain the boundary, so the universal zero-free half-plane is
$\Re s>\sigma_m$.

The dominance functions satisfy
\[
 \boxed{D_m'(\sigma)=-D_{m+1}(\sigma),}
\tag{6.8}
\]
and
\[
 D_m(\sigma)<0\ (1<\sigma<\sigma_m),
 \qquad
 D_m(\sigma)>0\ (\sigma>\sigma_m).
\tag{6.9}
\]
\end{theorem}

\begin{proof}
After division by the $p=2$ term, the left side of (6.6) decreases strictly from
$+\infty$ to $0$ as $\sigma$ runs from $1^+$ to infinity.  This proves existence,
uniqueness, and (6.9).  On $[\sigma_m,\infty)$ the $p=2$ term is the largest one, so
$R_m=D_m$ and (6.4) gives (6.7).  Equality at the edge is attained in the hull by
putting the $p=2$ term opposite to every other term, and Hurwitz approximation puts the
edge in the closure of actual zero real parts.

For the untwisted series, an actual boundary zero would force equality in the triangle
inequality.  The phases at $p=3,5,7$ would then impose rational relations between
$\log(5/3)$ and $\log(7/3)$, which unique factorization forbids.  Termwise
differentiation gives (6.8).
\end{proof}

The differential recursion turns a list of separate boundaries into a single
geometric mechanism.

\begin{theorem}[Derivative--caustic correspondence]
\label{thm:caustic}
Define the right-hand $2$-dominance caustic radius
\[
 \boxed{r_m^*=D_m(\sigma_{m+1}).}
\tag{6.10}
\]
Then $r_m^*>0$, and $D_m$ increases strictly on
$(\sigma_m,\sigma_{m+1})$, decreases strictly on
$(\sigma_{m+1},\infty)$, and tends to $0$ at infinity.

If $0<r<r_m^*$, there are unique
\[
 \sigma_m<\ell_m(r)<\sigma_{m+1}<u_m(r)
\]
with $D_m(\ell_m(r))=D_m(u_m(r))=r$.  Let $\alpha_m(r)$ solve
$A_m(\alpha_m(r))=r$.  Then, for every $|a|=r$,
\[
 \boxed{
 \overline{\{\Re s:\cP_{m,\omega}(s)=a\}}
 \cap[\sigma_m,\infty)
 = [\sigma_m,\ell_m(r)]\cup[u_m(r),\alpha_m(r)].
 }
\tag{6.11}
\]
At $r=r_m^*$ the two positive-density bands pinch at $\sigma_{m+1}$; the density
vanishes at the pinch.  If
$r_m^*<r<A_m(\sigma_m)$, the terminal projection is the single interval
$[\sigma_m,\alpha_m(r)]$.  Equality $r=A_m(\sigma_m)$ leaves only
$\{\sigma_m\}$, and larger $r$ leaves no point in $[\sigma_m,\infty)$.

The pinch is quadratic.  For $a=r_m^*e^{i\phi}$, the hull function
\[
 H_{m,\phi}(s)=e^{i\phi}\left((\log2)^m2^{-s}
 -\sum_{p\ge3}(\log p)^mp^{-s}\right)
\]
satisfies
\[
 H_{m,\phi}(\sigma_{m+1})=a,
 \quad H_{m,\phi}'(\sigma_{m+1})=0,
 \quad H_{m,\phi}''(\sigma_{m+1})\neq 0.
\tag{6.12}
\]
Equivalently,
\[
 \boxed{
 \text{the next derivative zero edge is the current derivative's
 right-hand $a$-point caustic.}
 }
\tag{6.13}
\]
\end{theorem}

\begin{proof}
Equations (6.8)--(6.9) give the asserted monotonicity and unique maximum.  On the
terminal branch, (6.4) reduces the $a$-point condition to
$D_m(\sigma)\le r\le A_m(\sigma)$, which yields (6.11) and all endpoint cases.
Finally,
\[
 H_{m,\phi}'=-e^{i\phi}D_{m+1},
 \qquad
 H_{m,\phi}''=e^{i\phi}D_{m+2}.
\]
At $\sigma_{m+1}$ the first expression is zero, whereas the second is nonzero because
$\sigma_{m+1}<\sigma_{m+2}$.
\end{proof}

\subsection{Internal dominance bands and the crystalline limit}
\label{subsec:internal-bands}

The terminal $2$-band is only one part of the derivative geometry.  We now
classify all internal bands.  Besides closing the internal-band problem posed
in Version 2.0, the classification
will give exact masses and barycentres for every bounded zero component.
For a prime $p$, put
\[
 g_{m,p}(\sigma)
 :=2a_{m,p}(\sigma)-A_m(\sigma)
 =a_{m,p}(\sigma)-\sum_{q\ne p}a_{m,q}(\sigma),
\tag{6.14}
\]
and, for $r\ge0$, define
\[
 B_{m,p}(r):=\{\sigma>1:g_{m,p}(\sigma)>r\},\qquad
 B_{m,p}:=B_{m,p}(0),
\tag{6.15}
\]
\[
 \mathfrak D_m:=\{p:B_{m,p}\ne\varnothing\},\qquad
 m_p:=\min\{m\ge0:p\in\mathfrak D_m\}.
\tag{6.16}
\]
Thus $R_m(\sigma)=\max\{0,\max_p g_{m,p}(\sigma)\}$.  We use the
phase-independent density $G_m(\sigma,a)$ from
Lemma~\ref{lem:fixed-order-transfer}.
For $p<q$ write
\[
 \tau_{p,q}:=
 \frac{\log(\log q/\log p)}{\log(q/p)},
 \qquad \sigma^{(m)}_{p,q}:=m\tau_{p,q}.
\tag{6.17}
\]
If $p\ge3$ has neighbouring primes $p^-<p<p^+$, let
\[
 W_p=(\tau_{p,p^+},\tau_{p^-,p}),
 \qquad W_2=(\tau_{2,3},\infty).
\tag{6.18}
\]
Strict concavity of $u\mapsto\log u$ shows that these windows are
nonempty and tile $(0,\infty)$ up to their endpoints.
For later use put
\[
 \phi_{p,q}(\tau)=\log\frac{\log q}{\log p}
                  -\tau\log\frac qp.
\]

\begin{theorem}[Complete dominance-band classification]
\label{thm:internal-band-classification}
The following assertions hold.
\begin{enumerate}[label=\textup{(\roman*)},leftmargin=2.2em]
\item For every $m\ge0$, prime $p$, and $r\ge0$, the set $B_{m,p}(r)$
is an open interval, possibly empty.

\item If $m\ge1$, then $B_{m,p}\subset mW_p$.  If $p<q$ both belong
to $\mathfrak D_m$, then
\[
 \sup B_{m,q}<\sigma^{(m)}_{p,q}<\inf B_{m,p}.
\tag{6.19}
\]
Consequently the nonempty bands are pairwise disjoint and are ordered,
from left to right, by decreasing prime.

\item The set $\mathfrak D_m$ is finite, contains $2$, and
\[
 B_{m,2}=(\sigma_m,\infty).
\]
Every $p\in\mathfrak D_m\setminus\{2\}$ satisfies
\[
 \frac{p}{(\log p)^m}<\frac{2}{(\log2)^m}.
\tag{6.20}
\]

\item For each unimodular twist define the twist-dependent closed
projection
\[
 Z_{m,\omega}:=
 \overline{\{\Re s:\cP_{m,\omega}(s)=0,\ \Re s>1\}}.
\]
Then the closure, although not the underlying set of actual zero real
parts, is independent of $\omega$ and is given exactly by
\[
 \boxed{
 Z_{m,\omega}=(1,\sigma_m]\setminus
 \bigcup_{p\in\mathfrak D_m\setminus\{2\}}B_{m,p}.}
\tag{6.21}
\]
Each $B_{m,p}$ is a zero-free vertical strip for every unimodular twist.
For the untwisted derivative its closed band $\overline{B_{m,p}}$ is
zero-free.  At every interior point of $Z_{m,\omega}$ the zero density
$G_m(\sigma,0)$ is positive.
\end{enumerate}
\end{theorem}

\begin{proof}
Multiplication by $p^\sigma$ gives
\[
 p^\sigma\bigl(g_{m,p}(\sigma)-r\bigr)
 = (\log p)^m-
 \sum_{q\ne p}(\log q)^m e^{-\sigma\log(q/p)}
 -r e^{\sigma\log p}.
\tag{6.22}
\]
The right side is strictly concave; hence its positive set is an open
interval.  Dominance by $p$ implies dominance over each individual term.
Comparison with the two neighbouring terms and monotonicity of secant
slopes of $\log$ give $B_{m,p}\subset mW_p$ for $m\ge1$ and (6.19).
Comparison with the $p=2$ term gives (6.20), hence finiteness.
The assertion for $p=2$ is Theorem~\ref{thm:derivative-edge}.

Finally, (6.4) at $a=0$ says that the projection closure is
$\{\sigma>1:\max_p g_{m,p}(\sigma)\le0\}$, which is (6.21).
The triangle inequality gives zero-freeness in an open band.  Equality
on a band endpoint for the untwisted series would force simultaneous
phase alignment for at least three distinct primes, contradicting unique
factorization.  The smooth-density theorem, applied to the weights
$(\log p)^m$, gives positivity at every strict interior point.
\end{proof}

The internal bands persist once born, and every prime eventually appears.

The following explicit estimate is useful both for proof and for certified
computation.  For $p\ge3$ set
\[
 r_+=\log\frac{p^+}{p},\quad r_-=\log\frac{p}{p^-},\quad
 c_+=\log\frac{\log p^+}{\log p},\quad
 c_-=\log\frac{\log p}{\log p^-},
\]
\[
 \tau_p^*=\frac{c_++c_-}{r_++r_-},\qquad
 \phi_p^*=\frac{c_+r_--c_-r_+}{r_++r_-}<0,
 \qquad \rho_p=\frac{r_+}{r_-},
\]
and
\[
 \mathcal K(\rho)=(1+\rho)\rho^{-\rho/(1+\rho)}\in(1,2].
\]
Also put
\[
 \delta_+=\tau_p^*-\frac1{\log p^+}>0,
 \qquad
 \delta_-=\frac1{\log p^-}-\tau_p^*>0,
\]
\[
 T_+(m)=
 \frac{p^+}{m\delta_+-1}
 \left(\frac{p^+}{p^++1}\right)^{m\delta_+-1}
 \quad(m\delta_+>1),
 \qquad
 T_-(m)=\frac{p^-}{m\delta_-+1},
\]
with $T_+(m)=\infty$ when $m\delta_+\le1$.

\begin{lemma}[Explicit birth bounds]
\label{lem:explicit-birth-bounds}
For every prime $p\ge3$,
\[
 \boxed{
 \frac{\log\mathcal K(\rho_p)}{|\phi_p^*|}<m_p
 \le
 \min\left\{m\ge1:
 m|\phi_p^*|>\log\bigl(2+T_+(m)+T_-(m)\bigr)\right\}.}
\tag{6.22a}
\]
The lower bound remains valid if $(p^-,p^+)$ is replaced by any
pair of primes $q_1<p<q_2$, with the corresponding parameters.
\end{lemma}

\begin{proof}
At $\tau=\tau_p^*+s$ the two neighbouring exponentials equal
\[
 e^{m\phi_p^*}
 \bigl(e^{-mr_+s}+e^{mr_-s}\bigr).
\]
Their minimum over $s\in\R$ is
$\mathcal K(\rho_p)e^{m\phi_p^*}$.  Dominance forces this number to
be $<1$, proving the lower bound.

For the converse, evaluate (6.26) at $\tau_p^*$.  Concavity of
$u\mapsto\log u$ gives
\[
 \sum_{q>p^+}e^{m\phi_{p,q}(\tau_p^*)}
 \le e^{m\phi_p^*}T_+(m),
 \qquad
 \sum_{q<p^-}e^{m\phi_{p,q}(\tau_p^*)}
 \le e^{m\phi_p^*}T_-(m),
\]
where the displayed expressions are elementary integral bounds for
the two tails.  Adding the two neighbouring terms gives
\[
 \sum_{q\ne p}e^{m\phi_{p,q}(\tau_p^*)}
 \le e^{m\phi_p^*}\bigl(2+T_+(m)+T_-(m)\bigr).
\]
The right side is $<1$ under the condition in (6.22a), and hence
$p\in\mathfrak D_m$.
\end{proof}

\begin{theorem}[Monotonicity, birth, and the prime-gap scale]
\label{thm:band-birth}
For every $m\ge0$ and every prime $p$,
\[
 \boxed{
 B_{m,p}+\frac1{\log p}\subset B_{m+1,p}.}
\tag{6.23}
\]
Thus $\mathfrak D_m\subset\mathfrak D_{m+1}$, every $m_p$ is finite,
and $p\in\mathfrak D_m$ if and only if $m\ge m_p$.

For $p\ge3$ let $g_-=p-p^-$, $g_+=p^+-p$, and
$g_{\max}=\max(g_-,g_+)$.  As $p\to\infty$,
\[
 \left(\frac{p\log p}{g_{\max}}\right)^2(1-o(1))
 \le m_p
 \ll \log(2+2g_{\max})
       \frac{p^2\log^2p}{g_-g_+}.
\tag{6.24}
\]
In particular, Cram\'er's gap conjecture would imply $m_p=p^{2+o(1)}$.
Unconditionally the normalized quantity
\[
 Y_p:=\frac{m_p}{p^2\log^2p}
\]
satisfies, along every sequence $p_j\to\infty$,
\[
 Y_{p_j}\longrightarrow0
 \quad\Longleftrightarrow\quad
 \max\{g_-(p_j),g_+(p_j)\}\longrightarrow\infty.
\]
In particular,
\[
 \liminf_{p\to\infty}Y_p=0<\limsup_{p\to\infty}Y_p<\infty.
\tag{6.25}
\]
More precisely, $\limsup Y_p>0$ is equivalent to
$\liminf_n(p_{n+2}-p_n)<\infty$.  The latter holds by Maynard's theorem on
infinitely many bounded clusters containing three primes \cite{Maynard2015}.
No assertion that $m_pg_-g_+/(p^2\log^2p)$ is unbounded is made.
\end{theorem}

\begin{proof}
Let $\sigma'=\sigma+1/\log p$.  The function
$u\mapsto ue^{-u/\log p}$ is uniquely maximized at $u=\log p$.
Termwise comparison at $u=\log q$ therefore proves (6.23).
For $\tau\in W_p$, write
\[
 \frac{\sum_{q\ne p}a_{m,q}(m\tau)}{a_{m,p}(m\tau)}
 =\sum_{q\ne p}\exp\{m\phi_{p,q}(\tau)\}.
\tag{6.26}
\]
Every exponent is negative on $W_p$, and dominated convergence makes
the sum tend to zero; this proves finite birth.

For completeness, put $L=\log p$.  Uniformly for the two neighbouring
gaps, the prime number theorem and Taylor expansion give
\[
 \phi_p^*=-\frac{r_+r_-}{2L^2}(1+o(1)),\qquad
 \delta_\pm=\frac{r_++r_-}{2L^2}(1+o(1)),\qquad
 r_\pm=\frac{g_\pm}{p}(1+o(1)).
\]
The elementary inequality
\[
 2\log\mathcal K(\rho)\ge\min\{\rho,\rho^{-1}\}
 \qquad(\rho>0)
\]
and the lower half of (6.22a) therefore yield
\[
 m_p>\frac{\log\mathcal K(\rho_p)}{|\phi_p^*|}
 \ge(1-o(1))\left(\frac{pL}{g_{\max}}\right)^2.
\]
For the converse choose, with a sufficiently large absolute $C$,
\[
 m=\left\lceil
 C\log(2+2g_{\max})\frac{p^2L^2}{g_-g_+}
 \right\rceil.
\]
Then
\[
 m|\phi_p^*|=\left(\frac C2+o(1)\right)\log(2+2g_{\max}),
 \qquad
 T_+(m)+T_-(m)\ll
 \frac{g_-g_+}{(g_-+g_+)\log(2+2g_{\max})}.
\]
Thus the upper condition in (6.22a) holds and proves (6.24).  If
$g_{\max}\to\infty$, the upper bound is
$O(\log(2+2g_{\max})/g_{\max})$, because the other neighbouring gap is at
least $2$; hence $Y_p\to0$ along such a sequence.  On the other hand,
within each bounded three-prime cluster choose its middle prime.  Both neighbouring
gaps are then bounded, and the lower bound in (6.24) keeps $Y_p$ away from zero.
The upper bound in (6.24), split according as $g_{\max}$ is bounded or tends to
infinity, also bounds $Y_p$ globally and proves the displayed sequential
equivalence.  A positive limsup is equivalent to bounded $g_{\max}$ along a
subsequence, which is the same as $\liminf_n(p_{n+2}-p_n)<\infty$.  This proves
(6.25).
\end{proof}

The differential recursion is prime-by-prime, not merely terminal.
When $B_{m,p}=(\ell_{m,p},u_{m,p})$, we allow $u_{m,2}=\infty$.

\begin{theorem}[Prime-by-prime internal caustics]
\label{thm:internal-caustics}
If $p\in\mathfrak D_m$, then $g_{m,p}$ has a unique global maximum at
$\ell_{m+1,p}$.  With
\[
 r^*_{m,p}:=g_{m,p}(\ell_{m+1,p}),
\tag{6.27}
\]
one has $r^*_{m,p}>0$ and
\[
 g_{m,p}'(\ell_{m+1,p})=0,
 \qquad
 g_{m,p}''(\ell_{m+1,p})<0.
\tag{6.28}
\]
For $0\le r<r^*_{m,p}$ the set $B_{m,p}(r)$ is a nonempty open
interval containing $\ell_{m+1,p}$; its closure shrinks to that point as
$r\uparrow r^*_{m,p}$, and it is empty for $r\ge r^*_{m,p}$.  Moreover,
\[
 \ell_{m,p}<\ell_{m+1,p}\le
 \ell_{m,p}+\frac1{\log p}.
\tag{6.29}
\]
For $p\ge3$ one also has
\[
 u_{m+1,p}\ge u_{m,p}+\frac1{\log p}.
\]
Thus every dominant prime has its own derivative--caustic
correspondence: the pinch of its level-$m$ $a$-point band is the left
edge of its level-$(m+1)$ zero-free band.
\end{theorem}

\begin{proof}
Termwise differentiation gives
\[
 g_{m,p}'=-g_{m+1,p},\qquad g_{m,p}''=g_{m+2,p}.
\tag{6.30}
\]
Theorem~\ref{thm:band-birth} ensures that the next two $p$-bands exist.
For $p\ge3$, since $g_{m+1,p}$ is nonpositive outside its band and positive
inside, $g_{m,p}$ increases on $(1,\ell_{m+1,p})$, decreases on
$(\ell_{m+1,p},u_{m+1,p})$, and increases again after $u_{m+1,p}$ toward
$0$ from below.  Because $B_{m,p}$ is nonempty, its positive interval must
therefore contain $\ell_{m+1,p}$, where the unique global maximum is
attained; beyond $u_{m+1,p}$ the function cannot rival this positive value.
For $p=2$, the same conclusion follows from the sign change of
$D_{m+1}$ at $\sigma_{m+1}$ and the fact that $D_m(\sigma)\to0^+$.
At the left endpoint of the positive set of
$g_{m+1,p}$ one has $g'_{m+1,p}(\ell_{m+1,p})>0$, so
\[
 g_{m+2,p}(\ell_{m+1,p})=-g'_{m+1,p}(\ell_{m+1,p})<0.
\]
This proves the transverse quadratic pinch.
The remaining assertions follow from the interval property and (6.23).
\end{proof}

We next describe the high-derivative limit.  Fix consecutive primes
$p<q=p^+$, set $d_{p,q}=\log(q/p)$ and $\tau=\tau_{p,q}$, and put
\[
 \eta_{m,p}:=
 \sum_{r\notin\{p,q\}}
 \frac{a_{m,r}(m\tau)}{a_{m,p}(m\tau)}.
\tag{6.31}
\]
At $m\tau$ the $p$- and $q$-terms are equal.  Furthermore
\[
 \eta_{m,p}=\Theta_p(e^{-m\mu_p}),\qquad
 \mu_p:=\min_{r\notin\{p,q\}}
 \left[-\log\frac{\log r}{\log p}+\tau\log\frac rp\right]>0.
\tag{6.32}
\]
The leading constant in (6.32) is the number, one or two, of minimizers in
the definition of $\mu_p$.

\begin{theorem}[Thin zero windows and the picket-fence limit]
\label{thm:picket-fence}
For each fixed consecutive pair $p<q=p^+$, as $m\to\infty$,
\[
 \ell_{m,p}-m\tau_{p,q}
  =(1+o(1))\frac{\eta_{m,p}}{d_{p,q}},
 \qquad
 m\tau_{p,q}-u_{m,q}
  =(1+o(1))\frac{\eta_{m,p}}{d_{p,q}}.
\tag{6.33}
\]
In particular every band fills its window:
\[
 \frac{\ell_{m,p}}m\longrightarrow\tau_{p,p^+},
 \qquad
 \frac{u_{m,p}}m\longrightarrow\tau_{p^-,p}\quad(p\ge3).
\tag{6.34}
\]
Locally on $(0,\infty)$ the closed sets $m^{-1}Z_{m,\omega}$ converge in the
Fell topology to the crossover set.  Equivalently, for every compact interval
$K\Subset(0,\infty)$ whose boundary contains no crossover point,
\[
 \left(m^{-1}Z_{m,\omega}\right)\cap K
 \longrightarrow
 \{\tau_{p,p^+}:p\in\PP\}\cap K
\tag{6.35}
\]
in Hausdorff distance, uniformly in the choice of unimodular twist.
The component near $m\tau_{p,p^+}$ has width
\[
 (2+o(1))\frac{\eta_{m,p}}{\log(p^+/p)}.
\tag{6.36}
\]
\end{theorem}

\begin{proof}
Divide the dominance inequality by the $p$-term and expand at
$\sigma=m\tau_{p,q}$.  The $q/p$ ratio contributes
$e^{-d_{p,q}s}=1-d_{p,q}s+O(s^2)$, while the sum of all remaining
ratios is $\eta_{m,p}(1+o(1))$ uniformly for
$|s|=O(\eta_{m,p})$.  Solving the two endpoint equations proves
(6.33).  Formula (6.34) follows from the window inclusion.  On a compact
$K\Subset(0,\infty)$ only finitely many windows occur; every associated
prime has been born for all sufficiently large $m$.  Equations
(6.21) and (6.33) then give (6.35)--(6.36).
\end{proof}

The limiting set alone forgets how much divisor mass lies at each
picket.  The Jessen law restores that information exactly.  Let
$J_{m,a}$ denote the analogue of (3.7) for $\cP_{m,\omega}$, retaining
the notation $G_m$ fixed above.  If $q>p$ are adjacent members of
$\mathfrak D_m$,
write
\[
 C_{m;q,p}:=[u_{m,q},\ell_{m,p}]
\tag{6.37}
\]
for the bounded zero-projection component between their bands.

\begin{theorem}[Exact component mass, barycentre, and local recovery]
\label{thm:component-tomography}
For every adjacent dominant pair $p<q$,
\[
 \boxed{
 \int_{C_{m;q,p}}G_m(\sigma,0)\dd\sigma
 =\frac1{2\pi}\log\frac qp,}
\tag{6.38}
\]
\[
 \boxed{
 \int_{C_{m;q,p}}\sigma G_m(\sigma,0)\dd\sigma
 =\frac m{2\pi}\log\frac{\log q}{\log p}.}
\tag{6.39}
\]
Consequently its density barycentre is exactly $m\tau_{p,q}$, and for
each of the multiplicity, distinct, and simple counting conventions, with $m$ and
the component fixed before $T\to\infty$,
\[
 N_{m;q,p}^{\bullet}(T)
 =\frac{\log(q/p)}{2\pi}T+o(T).
\tag{6.40}
\]
Conversely, this one component recovers its two boundary primes.  If
\[
 D=2\pi\int_CG_m(\sigma,0)\dd\sigma,
 \qquad
 E=\frac{2\pi}{m}\int_C\sigma G_m(\sigma,0)\dd\sigma,
\]
then
\[
 \boxed{
 p=\exp\!\left(\frac{D}{e^E-1}\right),
 \qquad q=pe^D.}
\tag{6.41}
\]
\end{theorem}

\begin{proof}
On a $p$-dominance band, Jensen's formula gives
\[
 J_{m,0}(\sigma)=\log a_{m,p}(\sigma)
 =m\log\log p-\sigma\log p,
\]
so $J_{m,0}'=-\log p$.  Integrating
$J_{m,0}''=2\pi G_m(\sigma,0)$ across (6.37) proves (6.38).
One integration by parts gives (6.39).  Asymptotic simplicity gives
(6.40).  Finally $D=\log(q/p)$ and
$E=\log(\log q/\log p)$, from which (6.41) follows algebraically.
\end{proof}

\begin{theorem}[All-target topology and caustic conservation]
\label{thm:all-target-band-topology}
Fix $m\ge1$ and $r=|a|>0$, and let $\alpha_m(r)$ be the unique solution of
$A_m(\alpha_m(r))=r$.  Put
\[
 \mathfrak D_m(r)=\{p\in\mathfrak D_m:r<r^*_{m,p}\}.
\]
Then
\[
 \boxed{
 \overline{\{\Re s:\cP_{m,\omega}(s)=a,\ \Re s>1\}}
 =(1,\alpha_m(r)]\setminus
 \bigcup_{p\in\mathfrak D_m(r)}B_{m,p}(r),}
\tag{6.42}
\]
and the number of connected components of this closed support is
\[
 \boxed{b_0(m,r)=1+\#\mathfrak D_m(r).}
\tag{6.43}
\]
At $r=0$ the corresponding number is $\#\mathfrak D_m$.

If $q>p$ label consecutive active bands from left to right and $C$ is the
bounded component between them, then the two formulas (6.38)--(6.39), and hence
the recovery formula (6.41), remain valid with $G_m(\sigma,0)$ replaced by
$G_m(\sigma,a)$.  If the $2$-band is active, the terminal component between its
right endpoint and $\alpha_m(r)$ has
\[
 \int_CG_m(\sigma,a)\dd\sigma=\frac{\log2}{2\pi},
 \qquad
 \int_C\sigma G_m(\sigma,a)\dd\sigma
 =\frac{m\log\log2-\log r}{2\pi}.
\tag{6.44}
\]
As $r$ crosses a caustic value $r^*_{m,p}$, the $p$-band disappears and the two
neighbouring components merge.  Their zeroth and first moments add exactly to
those of the merged component.  The same telescoping law holds for simultaneous
caustics.
\end{theorem}

\begin{proof}
Equation (6.42) is the phase-annulus criterion (6.4): the outer inequality is
$r\le A_m(\sigma)$ and every failure of the inner inequality is precisely one of
the disjoint sets $B_{m,p}(r)$.  Removing $k$ disjoint open intervals from
$(1,\alpha_m(r)]$ gives $k+1$ components, proving (6.43).  At the critical radius
the strict superlevel set is empty and its two endpoints have coalesced; the
support closure has therefore merged through a density-zero pinch.

On $B_{m,p}(r)$, conditioning on all phases except $U_p$ and using Jensen's
circle formula gives the same affine plateau as at target zero:
\[
 J_{m,a}(\sigma)=m\log\log p-\sigma\log p,
 \qquad J_{m,a}'(\sigma)=-\log p.
\]
For $A_m(\sigma)<r$, the same formula after a common Haar rotation gives
$J_{m,a}(\sigma)=\log r$ and $J_{m,a}'(\sigma)=0$.  Integrating
$J_{m,a}''=2\pi G_m(\sigma,a)$ once and then by parts proves all moment
formulas.  The merger laws reduce to
$\log(q/p)+\log(p/h)=\log(q/h)$ and the analogous identity for logarithms of
logarithms.
\end{proof}

The component masses have a useful scaling limit.  Define
\[
 \nu_m:=(\sigma\mapsto\sigma/m)_*
        \bigl(G_m(\sigma,0)\dd\sigma\bigr).
\]
Then, vaguely on $(0,\infty)$,
\[
 \boxed{
 \nu_m\Longrightarrow
 \frac1{2\pi}\sum_p\log\frac{p^+}{p}\,
       \delta_{\tau_{p,p^+}}.}
\tag{6.45}
\]
An atom $(\tau,w)$ of the limiting measure recovers its consecutive
prime pair by setting $D=2\pi w$ and using
\[
 p=\exp\!\left(\frac{D}{e^{\tau D}-1}\right),
 \qquad p^+=pe^D.
\tag{6.46}
\]
Moreover, if $w_p=(2\pi)^{-1}\log(p^+/p)$, then
\[
 \sum_{p\le p_N}w_p
 =\frac{\log p_{N+1}-\log2}{2\pi},
 \qquad
 \sum_{p\le p_N}\tau_{p,p^+}w_p
 =\frac{\log\log p_{N+1}-\log\log2}{2\pi}.
\tag{6.47}
\]
These exact telescoping laws show that the picket measure is a locally
finite, lossless encoding of the consecutive-prime chain, not a source
of information unavailable from all other edge data.

Indeed, on a compact subset of $(0,\infty)$ only finitely many prime windows
occur.  Every associated prime has been born for all sufficiently large $m$;
(6.33) collapses the intervening component to its crossover, while (6.38)
keeps its total mass exactly equal to the displayed atomic weight.  Testing
against a compactly supported continuous function proves the vague convergence
in (6.45).

\begin{remark}[Nonrigorous numerical orientation]
Floating-point minimization suggests, for example, the initial birth pattern
$m_3=10,m_5=22,m_7=44,m_{11}=103,\ldots$.
Such values are not used above.  Ordinary double-precision optimization,
including floating-point tail estimates, is not an exact certificate.  Any
numerical table promoted to theorem status must instead be verified by
directed interval arithmetic with archived code and data.
\end{remark}


\section{Convex edge interpolation and inverse prime spectroscopy}
\label{sec:edge-spectroscopy}

Interpolate the integer derivative order by the unique solution $\sigma(u)$, $u\ge0$,
of
\[
 \sum_{p\ge3}
 \left(\frac{\log p}{\log2}\right)^u
 \left(\frac2p\right)^{\sigma(u)}=1.
\tag{7.1}
\]
Thus $\sigma(m)=\sigma_m$.

\begin{proposition}[Strictly convex edge curve]
\label{prop:convex-edge}
The function $\sigma(u)$ is real analytic, strictly increasing, and strictly convex.
If
\[
 x_p=\log\frac{\log p}{\log2},
 \qquad y_p=\log\frac p2,
 \qquad w_p=e^{ux_p-\sigma(u)y_p},
\]
then $\sum_{p\ge3}w_p=1$ and
\[
 \sigma'(u)=\frac{\sum_pw_px_p}{\sum_pw_py_p}>0,
\tag{7.2}
\]
\[
 \boxed{
 \sigma''(u)=
 \frac{\sum_pw_p(x_p-\sigma'(u)y_p)^2}{\sum_pw_py_p}>0.
 }
\tag{7.3}
\]
\end{proposition}

\begin{proof}
The implicit-function theorem applies to (7.1).  First and second implicit
differentiation give (7.2)--(7.3).  Strictness follows because $x_p/y_p$ is not constant
in $p$.
\end{proof}

Set
\[
 L=\log\frac32,
 \qquad
 \kappa=\frac{\log(\log3/\log2)}{\log(3/2)}
 =1.1358825679163036\ldots,
\tag{7.4}
\]
and, for $p\ge3$,
\[
 q_p=\frac{\log p}{\log3}\left(\frac3p\right)^\kappa.
\tag{7.5}
\]
Then $q_3=1$, $q_p$ decreases strictly with $p$, and
\[
 q_5=0.8200411528584832\ldots,
 \qquad q_7=0.6765500255651028\ldots .
\tag{7.6}
\]

For orientation, direct prime summation gives the following initial values; the digits
shown are stable under enlargement of the prime cutoff.
\[
\begin{array}{c@{\qquad}c@{\qquad}c}
\toprule
m&\sigma_m&r_m^*\\
\midrule
0&1.77954465&0.08048045\\
1&2.60762384&0.02861372\\
2&3.50440655&0.00995714\\
3&4.44616892&0.00340929\\
4&5.42063867&0.00115243\\
5&6.42040617&0.00038547\\
\bottomrule
\end{array}
\tag{7.6a}
\]

\begin{theorem}[Asymptotic edge law]
\label{thm:edge-asymptotic}
As $u\to\infty$,
\[
 \boxed{
 \sigma(u)=\kappa u+\frac{q_5^u}{L}+O(q_7^u).
 }
\tag{7.7}
\]
In particular,
\[
 0<\sigma_{m+1}-\sigma_m
 <\sigma_{m+2}-\sigma_{m+1}<\kappa,
 \qquad
 \sigma_{m+1}-\sigma_m\longrightarrow\kappa.
\tag{7.8}
\]
Furthermore,
\[
 r_m^*=C B^m(1+O(q_5^m)),
\tag{7.9}
\]
where
\[
 B=(\log2)2^{-\kappa}=0.3154210960683259\ldots,
 \qquad
 C=2^{-\kappa}\left(1-\frac{\log2}{\log3}\right)
 =0.1679477965425128\ldots .
\tag{7.10}
\]
\end{theorem}

\begin{proof}
The function
\[
 q(x)=\frac{\log x}{\log3}\left(\frac3x\right)^\kappa
\]
is strictly decreasing for $x\ge3$, because
$(\log q)'=(1/\log x-\kappa)/x<0$.  Put
$h(u)=\sigma(u)-\kappa u$.  Separating the $p=3$ term in (7.1) gives the exact identity
\[
 \boxed{
 e^{Lh(u)}-1
 =\sum_{p\ge5}q_p^u\left(\frac3p\right)^{h(u)}.
 }
\tag{7.11}
\]
It follows first that $h(u)>0$ and $h(u)=O(q_5^u)$.  Separating $p=5$ gives
\[
 e^{Lh(u)}-1=q_5^u+O(q_5^{2u}+q_7^u).
\]
Since $q_5^2<q_7$, expansion of the logarithm proves (7.7).  Formula (7.2) is a
weighted average of $x_p/y_p$, whose strict maximum is the $p=3$ value $\kappa$.
Together with strict convexity this proves (7.8).  Substitution of
$\sigma_{m+1}=\kappa(m+1)+O(q_5^m)$ into
$D_m(\sigma_{m+1})$ gives (7.9)--(7.10).
\end{proof}

The small error in (7.7) successively records $5,7,11,\ldots$, not merely the first
correction.

\begin{theorem}[All primes from all derivative zero edges]
\label{thm:edge-recovery}
Given the exact sequence $(\sigma_m)_{m\ge0}$ and the anchors $2,3$, every prime is
recovered recursively.  Conversely, the prime sequence determines every $\sigma_m$.
Thus
\[
 \boxed{
 \{\text{complete derivative zero-edge tower}\}
 \quad\Longleftrightarrow\quad
 \{\text{complete prime spectrum}\}.
 }
\tag{7.12}
\]
\end{theorem}

\begin{proof}
The forward construction from primes to edges is (6.6).  For the inverse direction,
put
\[
 h_m=\sigma_m-\kappa m,
 \qquad E_m=e^{Lh_m}-1.
\]
Equation (7.11) is exact:
\[
 E_m=\sum_{p\ge5}q_p^m\left(\frac3p\right)^{h_m}.
\tag{7.13}
\]
Since $h_m\to0$ and $q_p$ decreases strictly,
\[
 q_5=\lim_{m\to\infty}E_m^{1/m}.
\tag{7.14}
\]
The strictly decreasing real function $q(x)$ then identifies $5$.  If the primes
$5,7,\ldots,p_j$ have been recovered, subtract their exact terms and put
\[
 E_m^{(j)}=E_m-
 \sum_{5\le p\le p_j}q_p^m\left(\frac3p\right)^{h_m}.
\]
Then
\[
 q_{p_{j+1}}=\lim_{m\to\infty}(E_m^{(j)})^{1/m}.
\tag{7.15}
\]
The upper bound follows by fixing $m_0\ge1$ and using
\[
 E_m^{(j)}\le q_{p_{j+1}}^{m-m_0}
 \sum_{p>p_j}q_p^{m_0},
\]
whose last sum is finite.  Induction recovers all primes.
\end{proof}


\section{Equivalent encodings of the zero-density law}
\label{sec:equivalences}

This section records what else closes once the exact density theorem is known.  The
statements are not new zero estimates; they are exact Tauberian, moment, and operator
encodings of the same measure.

\subsection{Regular variation and the logical strength of the BCS statements}

For fixed $\gamma$, the BCS height assertion
\[
 \frac{M_\gamma(xT)}{M_\gamma(T)}\longrightarrow x
 \qquad(0<x<1)
\tag{8.1}
\]
is exactly regular variation of index one:
$M_\gamma(T)=T\ell_\gamma(T)$ with $\ell_\gamma$ slowly varying.  Their separate
two-sided linear bound says $M_\gamma(T)=\Theta(T)$.  Even together these assertions do
not force $\ell_\gamma(T)$ to converge.  Theorem~\ref{thm:bcs} supplies the strictly
stronger conclusion
\[
 \ell_\gamma(T)\longrightarrow c_\gamma.
\tag{8.2}
\]

Now fix a compact interval $I\Subset(1,\sigma_{\rm e})$ and put
\[
 \kappa_I=\int_I\rho(\sigma)\dd\sigma>0.
\tag{8.3}
\]
List the positive ordinates of the simple zeros whose real parts lie in $I$:
\[
 0<\tau_1\le\tau_2\le\cdots,
\]
discarding a finite initial set if convenient.  Theorem~\ref{thm:asymptotic-simplicity}
shows that this choice has the same Weyl coefficient as either other counting
convention.  Let $D_Ie_n=\tau_ne_n$.

\begin{theorem}[Weyl--heat equivalence and its consequences]
\label{thm:jessen-dixmier}
The following three assertions are equivalent:
\[
 N_I(T)\sim\kappa_IT,
\tag{8.4}
\]
\[
 \tau_n\sim\frac n{\kappa_I},
\tag{8.5}
\]
\[
 \Tr(e^{-tD_I})\sim\frac{\kappa_I}{t}
 \qquad(t\downarrow0).
\tag{8.6}
\]
They hold for the prime-zeta divisor by Theorem~\ref{thm:bcs}.  They imply
\[
 \lim_{s\downarrow1}(s-1)\Tr(D_I^{-s})=\kappa_I,
\tag{8.7}
\]
\[
 \sum_{\tau_n\le T}\frac1{\tau_n}
 =\kappa_I\log T+o(\log T),
\tag{8.8}
\]
and
\[
 D_I^{-1}\in\mathcal L^{1,\infty}\setminus\mathcal S_1,
 \qquad
 \Tr_\Omega(D_I^{-1})=\kappa_I
\tag{8.9}
\]
for every Dixmier trace $\Tr_\Omega$.  Furthermore,
\[
 \mathcal S_I(z)=\det(I-z^2D_I^{-2})
 =\prod_n\left(1-\frac{z^2}{\tau_n^2}\right)
\tag{8.10}
\]
is an entire function of order one and exponential type $\pi\kappa_I$, and
\[
 \log\det(I+x^2D_I^{-2})\sim\pi\kappa_Ix.
\tag{8.11}
\]
Thus
\[
 \boxed{
 \kappa_I
 =\frac1{2\pi}\int_IJ_0''(\sigma)\dd\sigma
 =\lim_{T\to\infty}\frac{N_I(T)}T
 =\lim_{t\downarrow0}t\Tr(e^{-tD_I})
 =\Tr_\Omega(D_I^{-1})
 =\frac1\pi\type\mathcal S_I.
 }
\tag{8.12}
\]
\end{theorem}

\begin{proof}
The equivalence of (8.4) and (8.5) follows by inversion of a monotone counting
function.  The equivalence of (8.4) and (8.6) is Karamata's Laplace--Stieltjes theorem
\cite{BGT1987}.
Equations (8.7)--(8.8) follow by Abelian partial summation.  Relation (8.5) gives
$s_n(D_I^{-1})\sim\kappa_I/n$, proving measurability and (8.9); see
\cite{Connes1994,Simon2005} for the operator-ideal background.

Finally, $D_I^{-2}$ is trace class.  Partial summation of
$\sum_n\log(1+x^2/\tau_n^2)$ against $N_I$ gives
\[
 \kappa_I\int_0^\infty\log(1+x^2/t^2)\dd t
 =\pi\kappa_Ix+o(x),
\]
which proves (8.11); the standard canonical-product estimate
\cite{Levin1996} gives the stated type.
\end{proof}

\begin{remark}[No false converses]
A real-axis limit such as (8.7), a logarithmic harmonic asymptotic, or a single Dixmier
trace generally records only an averaged density and does not by itself imply the
pointwise Weyl law (8.4).  Only the Weyl--heat equivalence is invoked in both directions
without extra regular-variation assumptions.
\end{remark}

The two-variable limit (4.8) has the operator form
\[
 \frac1\Lambda\Tr\!\left[g(B)f(D/\Lambda)\right]
 \longrightarrow
 \left(\int g(\sigma)\rho(\sigma)\dd\sigma\right)
 \left(\int_0^\infty f(u)\dd u\right),
\tag{8.13}
\]
where $B$ is the diagonal operator of zero real parts, $D$ that of positive heights,
$g$ is continuous with compact support in $(1,\sigma_{\rm e})$, and
$f\in C_c([0,\infty))$.  Hence the Jessen measure is also a heat-trace tomography:
\[
 \lim_{t\downarrow0}t\Tr\bigl[g(B)e^{-tD}\bigr]
 =\int g(\sigma)\rho(\sigma)\dd\sigma.
\tag{8.14}
\]

\subsection{Moment tomography and the infinitely soft edge}

Put
\[
 \mu_n=\int_1^{\sigma_{\rm e}}\sigma^n\rho(\sigma)\dd\sigma.
\tag{8.15}
\]
Since $\rho$ is positive on an interval, its Hankel matrices are strictly positive:
\[
 \sum_{i,j=0}^Nc_i\overline{c_j}\mu_{i+j}>0
 \qquad((c_0,\ldots,c_N)\ne0).
\tag{8.16}
\]
In particular,
\[
 \mu_n^2<\mu_{n-1}\mu_{n+1},
 \qquad
 \lim_{n\to\infty}\mu_n^{1/n}
 =\lim_{n\to\infty}\frac{\mu_{n+1}}{\mu_n}
 =\sigma_{\rm e}.
\tag{8.17}
\]
Thus the edge can be recovered from the moment sequence alone.

The centered Jessen excess
\[
 K(\sigma)=J_0(\sigma)+\sigma\log2
\]
obeys the exact Green formula
\[
 \boxed{
 K(\sigma)=2\pi\int_\sigma^{\sigma_{\rm e}}
 (u-\sigma)\rho(u)\dd u.
 }
\tag{8.18}
\]
It is positive, strictly decreasing, and strictly convex before the edge, but every
derivative vanishes at $\sigma_{\rm e}$.  Hence $J_0$ is smooth but not real analytic
at the edge.  Equivalently, the top-edge Laplace profile
\[
 \Lambda(t)=\int_1^{\sigma_{\rm e}}
 e^{-t(\sigma_{\rm e}-\sigma)}\rho(\sigma)\dd\sigma
\tag{8.19}
\]
satisfies
\[
 \Lambda(t)=O_A(t^{-A})\quad\text{for every }A>0,
\tag{8.20}
\]
but is not $O(e^{-ct})$ for any $c>0$.  Spectrally, the dimension in every fixed
nonempty edge window remains one while its Weyl coefficient
$\kappa_{(\sigma_{\rm e}-\delta,\sigma_{\rm e})}$ vanishes faster than every power of
$\delta$ as $\delta\downarrow0$.

\section{Consequences for factorization series and prime subsets}
\label{sec:applications}

\subsection{Fr\"oberg-type ordered factorization series}

For $z\ne0$, consider
\[
 \mathcal D_z(s)=\frac1{1-z\cP(s)}.
\tag{9.1}
\]
In the half-plane where the geometric expansion is absolute,
\[
 \mathcal D_z(s)=\sum_{n\ge1}\frac{A_z(n)}{n^s},
 \qquad
 A_z(n)=z^{\Omega(n)}
 \frac{\Omega(n)!}{\prod_p\nu_p(n)!}.
\tag{9.2}
\]
Thus $A_z(n)$ counts ordered words in the prime factors of $n$, with weight $z$ per
letter.  For $z=1$, the identity, its Tauberian main term, and the resulting
oscillation problem were studied by Montgomery and Tenenbaum
\cite{MontgomeryTenenbaum2017}.  The poles of $\mathcal D_z$ are the $1/z$-points
of $\cP$.

\begin{corollary}[Exact pole density and sign blindness]
\label{cor:froberg-density}
Put
\[
 \beta_z=\cP^{-1}(|z|^{-1}).
\]
For $1<\gamma<\beta_z$,
\[
 \sum_{\substack{\mathcal D_z\text{ has a pole at }\rho\\
                  \Re\rho>\gamma\\0<\Im\rho<T}}
 m_\rho^{\rm pole}
 =T\int_\gamma^{\beta_z}G(\sigma,1/z)\dd\sigma+o(T).
\tag{9.3}
\]
The density depends on $z$ only through $|z|$.  In particular, the pole divisors of
$(1-\cP)^{-1}$ and $(1+\cP)^{-1}$ have identical limiting horizontal measures.
Their common right edge is
\[
 \beta_1=\cP^{-1}(1)=1.3994333287\ldots .
\tag{9.4}
\]
\end{corollary}

\begin{proof}
Apply Theorem~\ref{thm:smooth-a-law} to $a=1/z$ and use radiality.  The right edge is
Theorem~\ref{thm:right-edge-tomography}.
\end{proof}

For $z>0$, the real pole at $\beta_z$ is simple and has residue
$-1/(z\cP'(\beta_z))$.  A Tauberian main term, when applicable, therefore has constant
\[
 C_z=-\frac1{\beta_zz\cP'(\beta_z)}.
\tag{9.5}
\]
There can be no fixed power-saving remainder
\[
 \sum_{n\le x}A_z(n)=C_zx^{\beta_z}+O(x^{\beta_z-\delta})
\tag{9.6}
\]
for any $\delta>0$: such a remainder would continue $\mathcal D_z$ meromorphically
to $\Re s>\beta_z-\delta$ with no pole there except $\beta_z$, whereas (9.3) supplies
linearly many further poles in every strip $(\beta_z-\eps,\beta_z)$.  For $z=1$ this
recovers the obstruction in \cite{MontgomeryTenenbaum2017}; (9.3) adds the exact
phase-blind pole density and asymptotic simplicity.

\subsection{Reciprocal abscissae and character phases}

Let $\chi$ be a Dirichlet character modulo $q$ and write
\[
 \cP_\chi(s)=\sum_p\chi(p)p^{-s},\qquad
 \cP_q(\sigma)=\sum_{p\nmid q}p^{-\sigma}.
\]
For $a\ne0$, let $\beta_{q,a}>1$ be the unique solution of
$\cP_q(\beta_{q,a})=|a|$.

\begin{theorem}[Reciprocal--abscissa equivalence]
\label{thm:reciprocal-abscissa}
Put $c(n)=\Omega(n)!/\prod_p\nu_p(n)!$.  In its half-plane of absolute convergence,
\[
 \boxed{
 \frac1{a-\cP_\chi(s)}
 =\frac1a\sum_{\substack{n\ge1\\(n,q)=1}}
 \frac{\chi(n)c(n)}{a^{\Omega(n)}n^s}.
 }
\tag{9.6a}
\]
The abscissa of ordinary convergence, the abscissa of absolute convergence, and the
right edge of its pole set all coincide:
\[
 \boxed{
 \sigma_c=\sigma_a
 =\sup\{\Re\rho:\cP_\chi(\rho)=a\}
 =\beta_{q,a}.
 }
\tag{9.6b}
\]
On each compact strip to the left of this edge, all but $o(T)$ of the poles are simple.
\end{theorem}

\begin{proof}
The multinomial theorem gives
$\cP_\chi(s)^k=\sum_{\Omega(n)=k}\chi(n)c(n)n^{-s}$, and the geometric series gives
(9.6a).  Taking absolute values gives exactly
\[
 \frac1{|a|}\sum_{k\ge0}\left(\frac{\cP_q(\sigma)}{|a|}\right)^k,
\]
so $\sigma_a=\beta_{q,a}$.  The triangle inequality excludes poles to its right.
Immediately to its left, the target $a$ lies in the strict interior of the value
annulus: its outer radius exceeds $|a|$, while its inner radius is strictly smaller.
The smooth $a$-point law supplies actual poles with real parts arbitrarily close to
$\beta_{q,a}$ from the left.  If $\sigma_c<\beta_{q,a}$, the Dirichlet series would be
holomorphic in a half-plane containing one of those poles, a contradiction.  Thus
$\sigma_c=\sigma_a$, and Theorem~\ref{thm:asymptotic-simplicity} gives the final claim.
\end{proof}

\subsection{Ordered factorization length modulo \texorpdfstring{$m$}{m}}

Let $\beta=\cP^{-1}(1)$ and, for $m\ge2$ and $0\le r<m$, put
\[
 \mathcal C_{m,r}(s)=
 \sum_{\substack{n\ge1\\\Omega(n)\equiv r\ ({\rm mod}\ m)}}\frac{c(n)}{n^s},
 \qquad \Omega(1)=0,\quad c(1)=1.
\]

\begin{theorem}[Modular prime-word decomposition]
\label{thm:modular-factorizations}
For $\omega=e^{2\pi i/m}$,
\[
 \boxed{
 \mathcal C_{m,r}(s)
 =\frac1m\sum_{j=0}^{m-1}\frac{\omega^{-jr}}{1-\omega^j\cP(s)}
 =\frac{\cP(s)^r}{1-\cP(s)^m}.
 }
\tag{9.6c}
\]
All $m$ functions have exactly the same pole divisor in $\Re s>1$.  For every
$I\Subset(1,\beta)$, their simple, distinct, and multiplicity pole counts satisfy
\[
 N_{m,r}^{\bullet}(I;T)
 =mT\int_I G(\sigma,1)\dd\sigma+o(T),
 \qquad \bullet\in\{\mathrm{simp},\mathrm{dist},\mathrm{mult}\}.
\tag{9.6d}
\]
Their summatory functions are equidistributed:
\[
 \boxed{
 \sum_{\substack{n\le x\\\Omega(n)\equiv r\ ({\rm mod}\ m)}}c(n)
 \sim-\frac{x^\beta}{m\beta\cP'(\beta)}.
 }
\tag{9.6e}
\]
\end{theorem}

\begin{proof}
Root-of-unity filtering of
$\sum_{\Omega(n)=k}c(n)n^{-s}=\cP(s)^k$ gives (9.6c).  At a pole,
$\cP(s)^m=1$, so the numerator $\cP(s)^r$ never vanishes; the pole divisors are
therefore identical.  The $m$ distinct unit targets have the same radial density,
which gives (9.6d) and asymptotic simplicity.  At $s=\beta$ the residue is
$-1/(m\cP'(\beta))$.  Triangle equality and rational independence of the prime
logarithms show that no other point on $\Re s=\beta$ is singular.  Wiener--Ikehara
then proves (9.6e).
\end{proof}

\begin{corollary}[Exact cancellation and parity overtaking]
\label{cor:modular-cancellation}
For $r\ne r'\pmod m$, exactly $\gcd(m,r-r')$ of the $m$ unit-target pole families
cancel from $\mathcal C_{m,r}-\mathcal C_{m,r'}$; the other
$m-\gcd(m,r-r')$ remain.  Consequently, for every $\eps>0$, the difference of the
two summatory functions is
\[
 \Omega_\pm(x^{\beta-\eps}),
\tag{9.6f}
\]
and changes sign infinitely often.  In particular, ordered prime words have
asymptotically even and odd lengths in equal proportions, but the leading class
overtakes the other infinitely many times and the error admits no fixed power saving.
\end{corollary}

\begin{proof}
At a unit target $\xi$, cancellation is equivalent to $\xi^{r-r'}=1$, giving the
gcd count.  The remaining nonreal poles approach $\Re s=\beta$ arbitrarily closely by
the $a$-point law.  The Phragm\'en--Landau oscillation argument used in
\cite{MontgomeryTenenbaum2017} now applies verbatim to the real coefficient
difference and gives (9.6f).
\end{proof}

\subsection{Arithmetic progressions and Dirichlet phases}

Let
\[
 \cP_{q,a}(s)=\sum_{p\equiv a\, (\mathrm{mod}\,q)}p^{-s},
 \qquad (a,q)=1.
\tag{9.7}
\]
The prime number theorem in arithmetic progressions supplies enough prime blocks for
the smoothing proof.  Consequently, the projection, density, and tomography theorems
hold with the support restricted to this residue class.  In particular, the nonzero
$a$-point right-edge curve of $\cP_{q,a}$ reconstructs the complete prime subset
$\{p:p\equiv a\pmod q\}$.

For a Dirichlet character $\chi$,
\[
 \cP_\chi(s)=\sum_p\chi(p)p^{-s}
\tag{9.8}
\]
has coefficient modulus one away from primes dividing its modulus and zero at those
primes.  Hence all characters omitting the same finite prime set have identical
limiting $a$-point measures in $\Re s>1$, although their individual zero ordinates can
differ.  All-$a$ tomography recovers precisely the omitted primes and no phase data.
This gives a concrete classification of what the horizontal law can and cannot see.

There is also a field-theoretic consequence.  For a number field $K$, let
$S_K^{\rm spl}$ be the rational primes that split completely in $K$ and form the
prime-supported series $\sum_{p\in S_K^{\rm spl}}p^{-s}$.

\begin{corollary}[Complete-split tomography]
\label{cor:spectral-bauer}
The full $a$-point density, or equivalently the full right-edge curve, determines
$S_K^{\rm spl}$.  It therefore determines the Galois closure of $K$ up to
$\mathbb Q$-isomorphism.  If $K/\mathbb Q$ is Galois, it determines $K$ itself.
\end{corollary}

\begin{proof}
Target-plane conservation for all $\sigma$ gives the Dirichlet series
$\sum_{p\in S_K^{\rm spl}}(\log p)^2p^{-2\sigma}$, and uniqueness recovers its
support.  The same follows directly from edge tomography.  Bauer's theorem
\cite{Bauer1916} says that the completely split primes, up to finitely many exceptions,
determine the normal closure; the final assertion is immediate for a Galois field.
\end{proof}

Finally, in $\Re s>1$,
\[
 \log L(s,\chi)=\sum_{k\ge1}\frac1k\cP_{\chi^k}(ks).
\tag{9.9}
\]
Every Euler-logarithm layer therefore has a scaled phase-blind density and edge law.
The layers are not independent, so (9.9) is a structural decomposition rather than a
claim about zeros of $L(s,\chi)$ or about GRH.

\section{The exact logarithmic profile at the outer edge}
\label{sec:exact-edge-profile}

The smoothness argument proves that the zero density is flat at
$\sigma_{\rm e}$, but smoothness alone does not determine the scale of that
flatness.  We now compute it.  This section solves the edge-profile problem left
open in the preceding version of the paper.  All constants below are defined by
convergent integrals; decimal values are included only for orientation.

Put
\[
 A_{\ge3}(\sigma)=\sum_{p\ge3}p^{-\sigma},\qquad
 \eps(\sigma)=A_{\ge3}(\sigma)-2^{-\sigma}=-D_0(\sigma),
\tag{10.1}
\]
and write
\[
 S_\sigma=\sum_{p\ge3}p^{-\sigma}U_p,\qquad x_0(\sigma)=2^{-\sigma}.
\]
Thus $\eps>0$ for $1<\sigma<\sigma_{\rm e}$ and
\[
 \eps(\sigma)=|D_1(\sigma_{\rm e})|
 (\sigma_{\rm e}-\sigma)(1+O(\sigma_{\rm e}-\sigma)).
\tag{10.2}
\]
Let $I_0$ be the modified Bessel function and define
\[
 h(x)=x-\log I_0(x),\qquad
 H_j(\sigma)=\int_0^\infty h(x)(\log x)^j x^{-1/\sigma-1}\dd x,
 \qquad j\in\mathbb Z_{\ge0}.
\tag{10.3}
\]
The integrals converge because
$h(x)=x-x^2/4+O(x^4)$ at zero and
$h(x)=\tfrac12\log(2\pi x)+O(x^{-1})$ at infinity.
For $\lambda\ge0$ set
\[
 \Lambda_\sigma(\lambda)=\sum_{p\ge3}\log I_0(\lambda p^{-\sigma}),
 \qquad
 \Psi_\sigma(\lambda)=\lambda A_{\ge3}(\sigma)-\Lambda_\sigma(\lambda)
 =\sum_{p\ge3}h(\lambda p^{-\sigma}).
\tag{10.4}
\]
For $1<\sigma<\sigma_{\rm e}$ there is a unique $\lambda_\sigma>0$ satisfying
\[
 \Lambda_\sigma'(\lambda_\sigma)=x_0(\sigma),
 \quad\text{equivalently}\quad
 \Psi_\sigma'(\lambda_\sigma)=\eps(\sigma).
\tag{10.5}
\]
Moreover, $\lambda_\sigma\to\infty$ as $\sigma\uparrow\sigma_{\rm e}$.
We use the nonconflicting notation
\[
 \mathscr I_{\rm e}(\sigma)
 =\Lambda_\sigma^*(x_0)
 =\lambda_\sigma x_0-\Lambda_\sigma(\lambda_\sigma)
 =\Psi_\sigma(\lambda_\sigma)-\lambda_\sigma\eps,
 \qquad \mathcal L_\sigma=\log\lambda_\sigma.
\tag{10.6}
\]

\begin{lemma}[Radial reduction and derivative weight]
\label{lem:edge-radial-reduction}
Let $f_{S,\sigma}$ be the planar density of $S_\sigma$.  As
$\sigma\uparrow\sigma_{\rm e}$,
\[
 \boxed{
 \rho(\sigma)=D_1(\sigma_{\rm e})^2
  \bigl(1+O(\sqrt\eps)\bigr)f_{S,\sigma}(x_0).}
\tag{10.7}
\]
In particular $-\log\rho(\sigma)=-\log f_{S,\sigma}(x_0)+O(1)$.
\end{lemma}

\begin{proof}
Rotation invariance gives $f_{X,\sigma}(0)=f_{S,\sigma}(x_0)$.  On the fibre
$X_\sigma=0$, write $S_\sigma=|S_\sigma|e^{i\psi}$; then
\[
 \sum_{p\ge3}p^{-\sigma}|1-U_pe^{-i\psi}|^2=2\eps.
\]
After substituting $2^{-\sigma}U_2=-S_\sigma$ into $Y_\sigma$ and applying
Cauchy--Schwarz,
\[
 \left||Y_\sigma|-\bigl(|D_1(\sigma)|-(\log2)\eps\bigr)\right|
 \le
 \left(2\eps\sum_{p\ge3}\log^2(p/2)p^{-\sigma}\right)^{1/2}.
\]
The last prime sum is locally bounded at $\sigma_{\rm e}$.  Integrating
$|Y_\sigma|^2$ over the conditional density on the fibre proves (10.7).
\end{proof}

\begin{lemma}[Mellin--PNT expansion]
\label{lem:mellin-pnt-edge}
Uniformly for $\sigma$ in a compact subset of $(1,\infty)$, for every fixed
$K\in\mathbb Z_{\ge0}$ and $\lambda\to\infty$,
\[
 \Psi_\sigma(\lambda)
 =\lambda^{1/\sigma}\left(
   \sum_{j=0}^{K}\frac{H_j(\sigma)}{(\log\lambda)^{j+1}}
   +O_K((\log\lambda)^{-K-2})\right).
\tag{10.8}
\]
The differentiated expansion is obtained by termwise differentiation; in
particular its first two coefficients are $H_0/\sigma$ and
$H_1/\sigma-H_0$.
\end{lemma}

\begin{proof}
Summation by parts and the prime number theorem with de la Vall\'ee Poussin
remainder reduce the prime sum, up to
$O(\lambda^{1/\sigma}e^{-c\sqrt{\log\lambda}})$, to
\[
 M_\sigma(\lambda)=\int_2^\infty
   \frac{h(\lambda t^{-\sigma})}{\log t}\dd t.
\]
The substitution $x=\lambda t^{-\sigma}$ gives
\[
 M_\sigma(\lambda)=\lambda^{1/\sigma}
 \int_0^{\lambda2^{-\sigma}}
 \frac{h(x)x^{-1/\sigma-1}}{\log\lambda-\log x}\dd x.
\]
On $\lambda^{-1/2}\le x\le\lambda^{1/2}$ expand the denominator in powers
of $(\log x)/(\log\lambda)$.  The two tails are negligible by the endpoint
estimates for $h$, and extending each coefficient integral to $(0,\infty)$
gives (10.8).  Repeating the argument with
$\widetilde h(x)=xh'(x)$ gives the derivative formula.  Here
$\widetilde h(x)=x+O(x^2)$ at zero and
$\widetilde h(x)=\tfrac12+O(x^{-1})$ at infinity, so the same tail split is
valid.  Integration by parts yields, for $j\ge1$,
\[
 \int_0^\infty \widetilde h(x)(\log x)^j x^{-1/\sigma-1}\dd x
 =\frac{H_j(\sigma)}\sigma-jH_{j-1}(\sigma).
\]
For $j=0$ the corresponding coefficient is $H_0(\sigma)/\sigma$.
\end{proof}

\begin{lemma}[Exponential-tilt density bound]
\label{lem:tilted-edge-density}
For $\sigma<\sigma_{\rm e}$ sufficiently close to $\sigma_{\rm e}$, with
$\mathscr I_{\rm e}$ as in (10.6),
\[
 \mathscr I_{\rm e}(\sigma)-2\mathcal L_\sigma-O(1)
 \le -\log f_{S,\sigma}(x_0)
 \le \mathscr I_{\rm e}(\sigma)+O(\mathscr I_{\rm e}(\sigma)^{1/2}).
\tag{10.9}
\]
\end{lemma}

\begin{proof}
Tilt the independent phases by
\[
 \frac{\dd\mathbb P_\lambda}{\dd\mathbb P}
 =\exp(\lambda\Re S_\sigma-\Lambda_\sigma(\lambda)).
\]
Under this law the phase of the $p$th summand has von Mises density
$e^{\lambda p^{-\sigma}\cos\theta}/(2\pi I_0(\lambda p^{-\sigma}))$.
Five primes of size comparable with $\lambda^{1/\sigma}$, together with the
van der Corput estimate for their circle integrals, give
\[
 \|f^{(\lambda)}_{S,\sigma}\|_\infty\ll\lambda^2.
\]
Indeed, Bertrand's postulate supplies five primes $q_j$ with
$\lambda q_j^{-\sigma}\asymp1$.  The characteristic function of each
corresponding tilted circle is bounded by
$\min\{1,C(q_j^{-\sigma}|u|)^{-1/2}\}$, uniformly in the direction of
$u\in\mathbb R^2$.  The product of five such bounds is integrable, and the
change of scale $u=\lambda v$ gives the displayed $L^1$ estimate by Fourier
inversion.  Convolution with all remaining summands cannot increase the
$L^\infty$ norm.
The exact change-of-measure identity
\[
 f_{S,\sigma}(z)=e^{-\lambda\Re z+\Lambda_\sigma(\lambda)}
 f^{(\lambda)}_{S,\sigma}(z)
\]
at $z=x_0$ and $\lambda=\lambda_\sigma$ proves the left inequality.

For the reverse inequality put $a_p=p^{-\sigma}$ and isolate the two circles
belonging to $3$ and $5$.
Their convolution has density
\[
 \frac{\one_{|a_3-a_5|<|z|<a_3+a_5}}
 {\pi^2\sqrt{(|z|^2-(a_3-a_5)^2)((a_3+a_5)^2-|z|^2)}}
 \ge \frac1{2\pi^2a_3a_5}
\]
on its open annulus.  Put
\[
 S' =\sum_{p\ge7}a_pU_p,\qquad
 A_7=\sum_{p\ge7}a_p,\qquad
 \xi=A_7-\Re S',\qquad \eta=\Im S'.
\]
Tilt $S'$ with parameter $\lambda_7$ chosen so that
$\mathbb E_{\lambda_7}\xi=(1-2\vartheta)\eps$, where
$\vartheta=C(\log\lambda_\sigma)^{1/2}\lambda_\sigma^{-1/(2\sigma)}$
and $C$ is a sufficiently large absolute constant.  Consider the event
\[
 \mathcal E=\left\{
 (1-3\vartheta)\eps\le\xi\le(1-\vartheta)\eps,
 \quad \eta^2<\vartheta\eps(a_3+a_5)
 \right\}.
\]
For $\sigma$ sufficiently close to $\sigma_{\rm e}$, elementary geometry
shows that $x_0-S'$ lies in the open annulus on $\mathcal E$.

The standard von Mises moment estimates, summed over $p\ge7$, give
\[
 \operatorname{Var}_{\lambda_7}(\xi)
 \ll\frac{\lambda_7^{1/\sigma-2}}{\log\lambda_7},
 \qquad
 \mathbb E_{\lambda_7}\eta^2\ll\lambda_7^{-1}.
\]
Lemma~\ref{lem:mellin-pnt-edge}, differentiated once, implies
$\lambda_7\asymp\lambda_\sigma$ and
$\eps\asymp\lambda_\sigma^{1/\sigma-1}/\log\lambda_\sigma$.
In particular,
\[
 \frac{\operatorname{Var}_{\lambda_7}(\xi)}
      {\vartheta^2\eps^2}\ll C^{-2},
 \qquad
 \frac{\mathbb E_{\lambda_7}\eta^2}
      {\vartheta\eps(a_3+a_5)}=o(1).
\]
Chebyshev's inequality therefore gives
$\mathbb P_{\lambda_7}(\mathcal E)\ge c>0$ when $C$ is chosen large.
Define the deficit rate functions
\[
 \mathscr R_\sigma(u)=
 \sup_{\lambda\ge0}\{\Psi_\sigma(\lambda)-\lambda u\},
 \qquad
 \mathscr R_{\sigma,\ge7}(u)=
 \sup_{\lambda\ge0}\left\{
  \sum_{p\ge7}h(\lambda p^{-\sigma})-\lambda u\right\}.
\]
Thus $\mathscr R_\sigma(\eps)=\mathscr I_{\rm e}(\sigma)$.  The inverse
change of measure yields
\[
 \mathbb P(\mathcal E)
 \ge c\exp\{-\mathscr R_{\sigma,\ge7}((1-2\vartheta)\eps)
                    -\lambda_7\vartheta\eps\}.
\]
Deleting the two fixed primes changes the deficit function $\Psi_\sigma$ by
$h(\lambda a_3)+h(\lambda a_5)=O(\log\lambda)$.  The envelope identity
$\mathscr R_\sigma'(u)=-\lambda_\sigma(u)$, where $\lambda_\sigma(u)$ is the
unique maximizing parameter in the definition of $\mathscr R_\sigma(u)$, and the
differentiated form of
Lemma~\ref{lem:mellin-pnt-edge} consequently give
\[
 \mathscr R_{\sigma,\ge7}((1-2\vartheta)\eps)
 =\mathscr I_{\rm e}(\sigma)
  +O(\vartheta\lambda_\sigma\eps+\log\lambda_\sigma).
\]
Since
$\lambda_\sigma\eps\asymp\mathscr I_{\rm e}$ and
$\vartheta\mathscr I_{\rm e}\asymp\mathscr I_{\rm e}^{1/2}$, integrating
the two-circle density over $\mathcal E$ proves the right inequality.
\end{proof}

Define
\[
 \beta_{\rm e}=\frac1{\sigma_{\rm e}-1},\qquad
 h_{\rm e}=\frac{H_1(\sigma_{\rm e})}{H_0(\sigma_{\rm e})},
\tag{10.10}
\]
\[
 C_{\rm e}=\frac{\sigma_{\rm e}-1}{\sigma_{\rm e}}
 H_0(\sigma_{\rm e})^{1+\beta_{\rm e}}
 \sigma_{\rm e}^{-\beta_{\rm e}},\qquad
 K_{\rm e}=C_{\rm e}(1+\beta_{\rm e})^{-(1+\beta_{\rm e})}.
\tag{10.11}
\]

\begin{theorem}[Exact logarithmic edge profile]
\label{thm:exact-edge-profile}
As $\sigma\uparrow\sigma_{\rm e}$,
\[
 \boxed{-\log\rho(\sigma)
 =\mathscr I_{\rm e}(\sigma)+O(\mathscr I_{\rm e}(\sigma)^{1/2}).}
\tag{10.12}
\]
More precisely, with $\mathcal L=\mathcal L_\sigma$,
\[
 \mathscr I_{\rm e}(\sigma)
 =C_{\rm e}\eps^{-\beta_{\rm e}}\mathcal L^{-(1+\beta_{\rm e})}
 \left(1+\frac{(1+\beta_{\rm e})h_{\rm e}}{\mathcal L}
 +O(\mathcal L^{-2})\right),
\tag{10.13}
\]
where
\[
 \mathcal L=(1+\beta_{\rm e})\left(
 \log\frac1\eps-\log\mathcal L
 +\log\frac{H_0(\sigma_{\rm e})}{\sigma_{\rm e}}\right)
 +O(\mathcal L^{-1}).
\tag{10.14}
\]
Consequently, if $\ell=\log(1/\eps)$, then
\[
 \boxed{
 -\log\rho(\sigma)
 =K_{\rm e}\eps^{-\beta_{\rm e}}\ell^{-(1+\beta_{\rm e})}
 \left(1+(1+\beta_{\rm e})\frac{\log\ell}{\ell}
 +O(\ell^{-1})\right).}
\tag{10.15}
\]
Equivalently, for $\delta=\sigma_{\rm e}-\sigma$ and
$\ell_\delta=\log(1/\delta)$,
\[
 -\log\rho(\sigma)
 =K'_{\rm e}\delta^{-\beta_{\rm e}}
 \ell_\delta^{-(1+\beta_{\rm e})}
 \left(1+(1+\beta_{\rm e})\frac{\log\ell_\delta}{\ell_\delta}
 +O(\ell_\delta^{-1})\right),
\tag{10.16}
\]
where $K'_{\rm e}=K_{\rm e}|D_1(\sigma_{\rm e})|^{-\beta_{\rm e}}$.
\end{theorem}

\begin{proof}
Lemma~\ref{lem:edge-radial-reduction} and
Lemma~\ref{lem:tilted-edge-density} prove (10.12).  Apply
Lemma~\ref{lem:mellin-pnt-edge} and its derivative at the saddle point
$\Psi_\sigma'(\lambda)=\eps$.  Solving the derivative equation gives
\[
 \lambda^{1-1/\sigma}
 =\frac{H_0(\sigma)}{\sigma\eps\log\lambda}
 \left(1+\frac{H_1(\sigma)/H_0(\sigma)-\sigma}{\log\lambda}
 +O((\log\lambda)^{-2})\right).
\]
Substitution into $\Psi_\sigma(\lambda)-\lambda\eps$ yields (10.13)--(10.14),
first with $\sigma$ in the constants.  Since
$\sigma_{\rm e}-\sigma\asymp\eps$, replacing $\sigma$ by $\sigma_{\rm e}$
changes the expression only within the stated error.  Iterating (10.14) once
and using (10.12) gives (10.15).  Finally (10.2) gives (10.16).
\end{proof}

Numerically,
\[
 \beta_{\rm e}\doteq1.28280015,\qquad
 K_{\rm e}\doteq1.38465461,\qquad
 K'_{\rm e}\doteq6.35167631.
\tag{10.17}
\]
These decimals are high-precision evaluations, not interval certificates and not part
of the logical statement.  The theorem explains the source of the infinite softness:
the edge position is fixed by the initial prime spectrum, while the cost of approaching
it is paid by aligning primes of size about $\lambda_\sigma^{1/\sigma_{\rm e}}$.

\begin{conjecture}[Gaussian prefactor]
\label{conj:edge-prefactor}
As $\sigma\uparrow\sigma_{\rm e}$, let $v_1(\sigma)$ and $v_2(\sigma)$
be the two coordinate variances of
$S_\sigma$ under the saddle-point tilt.  Then
\[
 \rho(\sigma)=
 \frac{D_1(\sigma_{\rm e})^2}{2\pi\sqrt{v_1(\sigma)v_2(\sigma)}}
 e^{-\mathscr I_{\rm e}(\sigma)}(1+o(1)).
\]
This requires a local central limit theorem for the tilted triangular array and is
not used elsewhere in the paper.
\end{conjecture}

\section{Picket transforms and an exact prime-gap dictionary}
\label{sec:picket-transforms}

The weighted limit (6.45) admits a useful transform.  Write $p_n$ for the
$n$th prime and set
\[
 w_n=\frac1{2\pi}\log\frac{p_{n+1}}{p_n},\qquad
 \tau_n=\frac{\log(\log p_{n+1}/\log p_n)}
 {\log(p_{n+1}/p_n)}.
\tag{11.1}
\]

\begin{theorem}[Mellin abscissa of the weighted picket fence]
\label{thm:picket-mellin}
The ordered series
\[
 Z_{\rm pf}(s)=\sum_{n\ge1}w_n\tau_n^s
\tag{11.2}
\]
converges, indeed absolutely, exactly in the half-plane $\Re s>1$.  For
$\Re s\le1$ its ordered partial sums are not Cauchy.
\end{theorem}

\begin{proof}
Put $x_n=\log p_n$ and $\Delta_n=x_{n+1}-x_n$.  Then
$w_n=\Delta_n/(2\pi)$ and the mean-value theorem gives
\[
 \tau_n=\frac{\log x_{n+1}-\log x_n}{x_{n+1}-x_n}
 =\frac1{\xi_n},\qquad x_n<\xi_n<x_{n+1}.
\tag{11.3}
\]
For $\sigma=\Re s>1$, the sum
$\sum\Delta_n\xi_n^{-\sigma}$ is a Riemann-sum comparison with
$\int_{\log2}^\infty x^{-\sigma}\dd x$, proving absolute convergence.

Conversely fix $s=\sigma+it$ with $\sigma\le1$ and choose $c>1$ so close to
$1$ that $|t|\log c<\pi/4$.  In each sufficiently large block
$X\le x_n<x_{n+1}\le cX$, the phases $\tau_n^{it}$ lie in one arc of length
less than $\pi/3$.  After a common rotation, the real part of the block sum is
bounded below by a constant multiple of
\[
 X^{-\sigma}\sum_{X\le x_n<x_{n+1}\le cX}\Delta_n
 \asymp X^{1-\sigma}.
\]
Here the omitted boundary interval is $O(1)$ because Bertrand's postulate gives
$\Delta_n\le\log2$.  Thus the block tails fail to tend to zero for $\sigma=1$
and grow for $\sigma<1$.
\end{proof}

\begin{theorem}[Topology growth in the derivative tower]
\label{thm:component-growth}
Let $b_m$ be the number of connected components of the closed zero projection
$Z_{m,\omega}$.  Then
\[
 b_m=\#\mathfrak D_m,\qquad
 b_m-b_{m-1}=\#\{p:m_p=m\}\quad(m\ge1),\qquad
 b_m\longrightarrow\infty.
\tag{11.4}
\]
Unconditionally,
\[
 \boxed{
 \frac{\sqrt m}{(\log m)^2}\ll b_m
 \ll\sqrt m.}
\tag{11.5}
\]
Under Cram\'er's prime-gap conjecture,
\[
 \boxed{b_m=m^{1/2+o(1)}.}
\tag{11.6}
\]
\end{theorem}

\begin{proof}
The first line follows from (6.21), monotone birth, and finite birth for every
prime.  Since $g_-g_+\ge2g_{\max}$ and
$\log(2+2g_{\max})/g_{\max}$ is bounded, the upper bound in (6.24) gives
$m_p\ll p^2\log^2p$.  Hence all primes
$p\ll\sqrt m/\log m$ have appeared by level $m$, and the prime number theorem
gives the lower half of (11.5).  For the upper half, the lower birth estimate
shows that a prime $p\in\mathfrak D_m$ must have an adjacent gap of length
\[
 \gg \frac{p\log p}{\sqrt m}.
\]
In a dyadic interval $x<p\le2x$, the disjoint consecutive-prime gaps have total
length $O(x)$ and each gap can be charged to at most two endpoints.  Hence that
interval contains $O(\sqrt m/\log x)$ dominant primes.  The
Baker--Harman--Pintz estimate $g_{\max}\ll p^{0.525}$
\cite{BakerHarmanPintz2001} first restricts every such prime to
$p\le m^{20/19+o(1)}$; summing the dyadic bounds above $\sqrt m$, and using
$\pi(\sqrt m)\ll\sqrt m/\log m$ below it, gives $b_m\ll\sqrt m$.
Under Cram\'er, the two sides of (6.24) give
$m_p=p^{2+o(1)}$, and inversion with the prime number theorem gives (11.6).
\end{proof}

The same atoms provide exact coordinates for standard prime-gap questions.  The
next corollary is a dictionary, not a proof of any of the conjectures named in it.

\begin{corollary}[Exact prime-gap dictionary]
\label{cor:gap-dictionary}
Let $g_n=p_{n+1}-p_n$.  Then
\[
 2\pi p_nw_n=p_n\log\left(1+\frac{g_n}{p_n}\right)
 =g_n(1+o(1)).
\tag{11.7}
\]
Consequently:
\begin{enumerate}[label=\textup{(\roman*)},leftmargin=2.2em]
\item The twin-prime conjecture is equivalent to
\[
 \liminf_{n\to\infty}p_nw_n=\frac1\pi,
\]
and also to $1/\pi$ being a subsequential limit of $(p_nw_n)$.
\item For every $h\ge1$, the Polignac assertion that $g_n=2h$ infinitely often
is equivalent to $h/\pi$ being a subsequential limit of $(p_nw_n)$.
\item Cram\'er's conjecture $g_n=O((\log p_n)^2)$ is equivalent to
\[
 w_n=O\left(\frac{(\log p_n)^2}{p_n}\right),
\tag{11.8}
\]
and also to $|W_p|=O(1/p)$ over the primes.
\item For every fixed $k\ge1$,
\[
 \liminf_n(p_{n+k}-p_n)
 =\liminf_n2\pi p_n\sum_{j=0}^{k-1}w_{n+j}.
\tag{11.9}
\]
\end{enumerate}
\end{corollary}

\begin{proof}
Formula (11.7) is the definition of $w_n$ and Taylor expansion, using
$p_{n+1}/p_n\to1$.  A bounded subsequence of $p_nw_n$ forces $g_n=O(1)$;
parity and integrality then provide a constant-gap subsequence.  This proves
(i)--(ii).  The first equivalence in (iii) follows again from (11.7).  Expansion
of the two secant slopes defining $W_p$ gives
\[
 |W_p|=\frac{g_-+g_+}{2p(\log p)^2}(1+o(1)),
\tag{11.10}
\]
which proves the second.  Finally the sum in (11.9) telescopes to
$p_n\log(p_{n+k}/p_n)$, differing from $p_{n+k}-p_n$ by $o(1)$ along every
subsequence on which either has finite limit.
\end{proof}

For reference, (11.7) also gives exact, purely reformulatory versions of two
other classical inequalities:
\[
 \sqrt{p_{n+1}}-\sqrt{p_n}<1
 \quad\Longleftrightarrow\quad
 w_n<\frac1\pi\log(1+p_n^{-1/2}),
\tag{11.11}
\]
and Firoozbakht's inequality is equivalent to
\[
 w_n<\frac{\log p_n}{2\pi n}.
\tag{11.12}
\]
No new case of either conjecture follows from these identities.

\subsection{Crystalline zero lattices at consecutive-prime crossovers}
\label{subsec:crystalline-lattices}

The picket limit describes horizontal divisor mass after vertical averaging.  At a
fixed consecutive-prime crossover one can retain the ordinates and obtain an actual
zero lattice.  Fix consecutive primes $p<q=p^+$ and put
\[
 d=d_{p,q}=\log\frac qp,\qquad
 \tau=\tau_{p,q}=\frac{\log(\log q/\log p)}{\log(q/p)}.
\tag{11.13}
\]
For a prime $r\notin\{p,q\}$ define
\[
 \psi_{p,q}(r)=\log\frac{\log r}{\log p}-\tau\log\frac rp,
 \qquad
 \mu_{p,q}=-\max_{r\notin\{p,q\}}\psi_{p,q}(r)>0.
\tag{11.14}
\]
Strict concavity shows that the maximum is attained among $p^-$ and $q^+$,
with only $q^+$ present when $p=2$.  Choose a real lift
$\theta_{p,q}=\arg(\omega_q/\omega_p)$ and set
\[
 t_k=\frac{\theta_{p,q}+(2k+1)\pi}{d},\qquad
 s^0_{m,k}=m\tau+it_k.
\tag{11.15}
\]
Changing the lift only reindexes $k$.

\begin{theorem}[Uniform two-prime zero lattice]
\label{thm:two-prime-lattice}
Fix $p<q=p^+$ and $A>0$.  For all sufficiently large integers $m$, so that
$m\tau-A>1$, uniformly over all unimodular twists and all $k\in\mathbb Z$,
the rectangle
\[
 \mathscr R_{m,k}(A)=\left\{s:
 |\Re s-m\tau|<A,\quad |\Im s-t_k|<\frac\pi d\right\}
\tag{11.16}
\]
contains exactly one zero $\rho_{m,k}$ of $\cP_{m,\omega}$, counted with
multiplicity.  It is simple; there are no zeros on $\partial\mathscr R_{m,k}(A)$;
and these points exhaust
the zeros in $|\Re s-m\tau|\le A$.  Moreover,
\[
 \boxed{
 \rho_{m,k}=m\tau+i\frac{\theta_{p,q}+(2k+1)\pi}{d}
 +O_{p,q,A}(e^{-\mu_{p,q}m}).}
\tag{11.17}
\]
If
\[
 E_{m,\omega}(z)=
 \sum_{r\notin\{p,q\}}\frac{\omega_r}{\omega_p}
 e^{m\psi_{p,q}(r)}e^{-z\log(r/p)},
\tag{11.18}
\]
then the first correction is
\[
 \boxed{
 \rho_{m,k}=s^0_{m,k}-\frac1dE_{m,\omega}(it_k)
 +O_{p,q,A}(e^{-2\mu_{p,q}m}).}
\tag{11.19}
\]
\end{theorem}

\begin{proof}
Write $s=m\tau+z$ and divide by the nonvanishing $p$-term.  The defining
identity for $\tau$ gives
\[
 \frac{\cP_{m,\omega}(m\tau+z)}
 {\omega_p(\log p)^mp^{-m\tau-z}}
 =1+\frac{\omega_q}{\omega_p}e^{-dz}+E_{m,\omega}(z).
\tag{11.20}
\]
The strictly concave function $x\mapsto\log x-\tau x$ takes the same value at
$x=\log p$ and $x=\log q$, and its maximum lies strictly between them.  Since
$p,q$ are consecutive, $\psi_{p,q}(r)<0$ for $r\notin\{p,q\}$ and
$\mu_{p,q}>0$.  For every fixed $j\ge0$,
\[
 \sup_{|\Re z|\le A}|\partial_z^jE_{m,\omega}(z)|
 \ll_{p,q,A,j}e^{-\mu_{p,q}m}.
\tag{11.21}
\]
Indeed, factor out $e^{-(m-m_0)\mu_{p,q}}$ with
$m_0\tau-A>1$; the remaining majorant is a convergent log-weighted prime
Dirichlet series.  This also proves uniformity in $k$ and in the twist.

Writing $z=it_k+w$, the two-term part becomes $1-e^{-dw}$.  On the vertical
sides of (11.16) its modulus is at least $1-e^{-dA}$, while on the horizontal
sides it is at least $1+e^{-dA}$.  Equation (11.21) and Rouch\'e's theorem give
one zero, counted with multiplicity, in every rectangle and no boundary zero.
The rectangles tile the fixed-width strip, so they exhaust it.  Rouch\'e on
$|w|=Ce^{-\mu_{p,q}m}$ gives (11.17).  Finally
\[
 1-e^{-dw}=dw+O(w^2),\qquad
 E_{m,\omega}(it_k+w)=E_{m,\omega}(it_k)
 +O(e^{-\mu_{p,q}m}|w|),
\]
which yields (11.19).
\end{proof}

\begin{corollary}[Exact strip count and crystalline spacing]
\label{cor:crystalline-count}
For fixed $p,q,A$ and all sufficiently large $m$,
\[
 \boxed{
 N_{m;p,q,A}(T)=\frac{\log(q/p)}{2\pi}T+O_{p,q,A}(1),}
\tag{11.22}
\]
uniformly over the twists, where the count is taken in
$|\Re s-m\tau|<A$ and $0<\Im s\le T$.  Furthermore,
\[
 \rho_{m,k+1}-\rho_{m,k}
 =\frac{2\pi i}{\log(q/p)}+O_{p,q,A}(e^{-\mu_{p,q}m}),
\tag{11.23}
\]
and, for a fixed twist and a consistent indexing,
\[
 \rho_{m+1,k}-\rho_{m,k}
 =\tau_{p,q}+O_{p,q,A}(e^{-\mu_{p,q}m}).
\tag{11.24}
\]
\end{corollary}

\begin{corollary}[Inverse recovery of the primes and phases]
\label{cor:lattice-inversion}
Suppose the moving lattice belonging to a consecutive pair is known and put
\[
 H=\lim_{m\to\infty}\Im(\rho_{m,k+1}-\rho_{m,k}),\qquad
 \tau=\lim_{m\to\infty}\frac{\Re\rho_{m,k}}m,\qquad d=\frac{2\pi}{H}.
\]
Then
\[
 \boxed{
 \log p=\frac{d}{e^{\tau d}-1},\qquad q=pe^d.}
\tag{11.25}
\]
The ordinate offset recovers the relative phase:
\[
 \boxed{
 -e^{id\Im\rho_{m,k}}=\frac{\omega_q}{\omega_p}
 +O_{p,q,A}(e^{-\mu_{p,q}m}).}
\tag{11.26}
\]
Thus all consecutive-prime crossover lattices recover the whole twist up to one
common unimodular factor.
\end{corollary}

\begin{proof}
Equations (11.17), (11.23), and the algebra used in (6.41) give (11.25).
Equation (11.15) gives (11.26), and consecutive ratios recover all phases after
one phase is chosen as a normalization.
\end{proof}

The leading defect also detects which external neighbour is closest in exponential
scale.  Let
\[
 \mathcal R_{p,q}=\{r\notin\{p,q\}:-\psi_{p,q}(r)=\mu_{p,q}\}
\]
If $p=2$, then $\mathcal R_{p,q}\subseteq\{q^+\}$; if $p\ge3$, then
$\mathcal R_{p,q}\subseteq\{p^-,q^+\}$.

\begin{corollary}[Second-neighbour defect profile]
\label{cor:defect-profile}
Uniformly in $k$,
\[
 e^{\mu_{p,q}m}(\rho_{m,k}-s^0_{m,k})
 =-\frac1d\sum_{r\in\mathcal R_{p,q}}
 \frac{\omega_r}{\omega_p}e^{-it_k\log(r/p)}+o(1).
\tag{11.27}
\]
Define
\[
 \mathscr K_m=\overline{\left\{
 e^{\mu_{p,q}m}(\rho_{m,k}-s^0_{m,k}):k\in\mathbb Z
 \right\}}.
\]
If $|\mathcal R_{p,q}|=1$, then
\[
 d_{\rm H}\bigl(\mathscr K_m,\{z:|z|=1/d\}\bigr)\longrightarrow0;
\]
if $|\mathcal R_{p,q}|=2$, then
\[
 d_{\rm H}\bigl(\mathscr K_m,\{z:|z|\le2/d\}\bigr)\longrightarrow0.
\]
\end{corollary}

\begin{proof}
Multiply (11.19) by $e^{\mu_{p,q}m}$; every nonmaximal tail term decays
exponentially.  With one surviving neighbour, a rational relation between
$\log(r/p)$ and $\log(q/p)$ would exponentiate to a nontrivial equality between
products of distinct primes, contradicting unique factorization.  Thus Kronecker's
theorem gives the circle.  With two neighbours, any rational relation among
$1$, $\log(p^-/p)/d$, and $\log(q^+/p)/d$ likewise exponentiates to a
nontrivial prime-product identity.  Unique factorization makes the three numbers
rationally independent; their torus orbit is dense, and sums of two unit vectors
fill the closed disc.  The error in (11.27) is uniform in $k$, which proves the
stated Hausdorff convergence.
\end{proof}

\begin{remark}[Scope and prior geometry]
The theorem fixes $p<q$ and then lets $m\to\infty$; it is not uniform for a prime
pair growing with $m$.  Its phase-recovery datum is the complete
ordinate-resolved moving tower, not phase-blind data at one derivative level.
The distinct anchored fixed-level recovery mechanism is developed below.
Binder, Pauli, and Saidak proved an analogous two-term rectangle theorem for high
derivatives of the Riemann zeta function \cite{BinderPauliSaidak2015}; hence the
Rouch\'e lattice mechanism itself is not claimed as new.  Prime-specific additions
here are the decay scale $\mu_{p,q}$, the uniform first correction, arbitrary
unimodular twists, inverse recovery, and the second-neighbour defect.  A 2017 lecture
of Nakamura was entitled ``On zeros of derivatives of the prime zeta-function''
\cite{Nakamura2017PrimeZeta}; priority for the bare prime-zeta specialization is
therefore left open.
\end{remark}

% The following four sections are maintained as separate source modules during
% editorial verification and are flattened into the distributed standalone file.
% ============================================================================
% Integrated replacement section: local zero processes and spacing laws
% Intended for insertion after the fixed-order transfer lemma and before the
% final limitations section of KEIO_Prime_Zeta_Tomography_Revised_v3.0_EN.tex.
%
% Bibliography-key candidates (not invoked here, so this file compiles with the
% present bibliography unchanged):
%   AzaisWschebor2009  -- multi-point Kac--Rice formula
%   DaleyVereJones2008 -- point processes, factorial moments, Palm measures
%   JessenTornehave1945 -- almost-periodic zero distributions
% ============================================================================

\section{Universal local zero processes, Palm laws, and hard cores}
\label{sec:local-process}

The horizontal laws proved above determine more than one-point densities.  At
every fixed derivative level they determine the complete local zero process,
including all finite-window counting laws and all nearest-neighbour laws.  The
point-process formulation is also the cleanest way to separate two genuinely
different small-scale regimes: quartic planar repulsion in the bulk and an exact
deterministic hard core in a dominance region.

Throughout this section, fix an integer $m\geq0$ and write
\begin{equation}
 F_{m,u}(s):=\sum_pu_p(\log p)^mp^{-s},
 \qquad u=(u_p)_p\in\mathbb H:=\mathbb T^{\PP},
 \qquad \Re s>1.
 \label{eq:lp-family}
\end{equation}
Thus $F_{m,1}=(-1)^m\cP^{(m)}$.  We also fix $a\in\C$.  The zero divisor of
$F_{m,u}-a$, with algebraic multiplicity, is denoted by
\begin{equation}
 \Xi_{m,u,a}:=\sum_{F_{m,u}(\rho)=a}m_\rho\,\delta_\rho .
 \label{eq:lp-divisor}
\end{equation}
It is a locally finite integer-valued measure on
$D:=\{s\in\C:\Re s>1\}$.  The vertical translation of a divisor is
$(\vartheta_t\Xi)(B)=\Xi(B+it)$, and the Kronecker flow on $\mathbb H$ is
\begin{equation}
 (R_tu)_p=u_pp^{-it}.
 \label{eq:lp-flow}
\end{equation}
The elementary but decisive covariance identity is
\begin{equation}
 F_{m,R_tu}(s)=F_{m,u}(s+it),
 \qquad
 \Xi_{m,R_tu,a}=\vartheta_t\Xi_{m,u,a}.
 \label{eq:lp-covariance}
\end{equation}

\subsection{Compactness of the hull and smooth joint jets}

\begin{lemma}[Continuous divisor map and a uniform local zero bound]
\label{lem:lp-compact-divisor}
The map
\[
 u\longmapsto F_{m,u}
\]
from $\mathbb H$ to the holomorphic functions on $D$, equipped with the
compact-open topology, is continuous.  The induced divisor map
\begin{equation}
 \mathbb H\longrightarrow\mathcal N_{\rm lf}(D),
 \qquad u\longmapsto\Xi_{m,u,a},
 \label{eq:lp-divisor-map}
\end{equation}
is continuous for the vague topology on locally finite divisors, where zeros
are counted with multiplicity.  Moreover, for every compact $K\Subset D$,
\begin{equation}
 M_{m,a}(K):=\sup_{u\in\mathbb H}\Xi_{m,u,a}(K)<\infty.
 \label{eq:lp-uniform-count}
\end{equation}
\end{lemma}

\begin{proof}
If $K\Subset D$ and $\Re s\geq1+\eta$ on $K$, then for every $j\geq0$,
\[
 \sum_p(\log p)^{m+j}p^{-1-\eta}<\infty.
\]
The tails of the series for $F_{m,u}$ and all its fixed derivatives are
therefore uniform in $u$.  Coordinatewise convergence in $\mathbb H$ implies
compact-open convergence of the corresponding functions.

No member of the hull is identically equal to $a$.  For $a\neq0$ this follows
by letting $\Re s\to\infty$, and for $a=0$ it follows from uniqueness of
absolutely convergent Dirichlet series.  Hurwitz's theorem, equivalently the
argument principle on contours avoiding the limiting divisor, now gives
continuity of the divisor map with multiplicities.

Suppose that \eqref{eq:lp-uniform-count} failed.  There would be $u_n\in
\mathbb H$ for which the number of zeros in $K$ tends to infinity.  After
passing to a subsequence, compactness of $\mathbb H$ gives $u_n\to u$.  Choose
a bounded domain $V\Subset D$ containing $K$ whose boundary contains no zero
of $F_{m,u}-a$.  Rouch\'e's theorem makes the number of zeros of
$F_{m,u_n}-a$ in $V$ equal to the finite number for $F_{m,u}-a$ for all large
$n$, a contradiction.
\end{proof}

For pairwise distinct $s_1,\ldots,s_k\in D$ and an integer $r\geq0$, set
\begin{equation}
 \mathcal J_{k,r}(u;\mathbf s)
 :=\bigl(F_{m,u}^{(\nu)}(s_j)\bigr)_
       {1\leq j\leq k,\ 0\leq\nu\leq r}
 \in\C^{k(r+1)}.
 \label{eq:lp-jet}
\end{equation}

\begin{lemma}[Smooth joint jet densities]
\label{lem:lp-jets}
Every vector \eqref{eq:lp-jet} has a $C^\infty$ density with bounded
derivatives of all orders.  The density depends continuously, in every
$C^r$ norm, on the points $s_j$ away from their diagonals.  If points coalesce,
the same assertion holds uniformly after replacing their values by Newton
divided differences; the limiting variables are the corresponding normalized
jets.
\end{lemma}

\begin{proof}
Put $N=k(r+1)$.  The vector \eqref{eq:lp-jet} is a sum of independent Haar
phases multiplied by coefficient vectors
\[
 c_p=\bigl((-\log p)^\nu(\log p)^m p^{-s_j}\bigr)_{j,\nu}\in\C^N.
\]
The functions $x\mapsto x^{m+\nu}e^{-s_jx}$ are linearly independent.
Consequently the vectors $c_p$ span $\C^N$.  Indeed, otherwise a nonzero
exponential polynomial would vanish at $x=\log p$ for every sufficiently large
prime.  Such a function has only $O(T)$ real zeros in $[0,T]$, whereas
$\#\{p:\log p\leq T\}\asymp e^T/T$.  The same argument after removing any
finite set of primes supplies arbitrarily many disjoint spanning $N$-tuples.

The characteristic function of \eqref{eq:lp-jet} is a product of Bessel
factors
\[
 \prod_pJ_0\bigl(|\langle\xi,c_p\rangle|\bigr).
\]
On each spanning tuple at least one scalar product has size
$\gg_{\mathcal T}|\xi|$.  Since $|J_0(x)|\ll(1+x)^{-1/2}$, using arbitrarily
many disjoint tuples gives decay faster than any prescribed negative power of
$|\xi|$.  Differentiating in $\xi$ or in the parameters inserts only fixed
powers of $\log p$, which remain summable on compact subsets of $D$.
Fourier inversion proves the smoothness, boundedness, and parameter
continuity.

Newton divided differences are linear combinations of the values at distinct
nodes and converge to normalized derivatives as the nodes merge.  Their
coefficient vectors vary continuously through this confluent limit.  Applying
the preceding characteristic-function argument to these vectors proves the
last assertion.
\end{proof}

\subsection{All correlation functions and the complete local process}

Let $q^{(m)}_{\mathbf s}$ denote the joint density of
$(F_{m,U}(s_j),F'_{m,U}(s_j))_{j=1}^k$ for Haar-distributed $U$.  For distinct
$s_j$ define
\begin{align}
 K_k^{(m,a)}(s_1,\ldots,s_k)
 &:={\mathbb E}\!\left[
   \prod_{j=1}^k|F'_{m,U}(s_j)|^2
   \prod_{j=1}^k\delta_a(F_{m,U}(s_j))
 \right]                                                     \notag\\
 &=\int_{\C^k}\prod_{j=1}^k|y_j|^2
 q^{(m)}_{\mathbf s}(a,y_1,\ldots,a,y_k)
 \prod_{j=1}^k\dd A(y_j).
 \label{eq:lp-kac-rice}
\end{align}
Here and below the delta notation is shorthand for evaluation of the smooth
density.  The multi-point complex Kac--Rice formula identifies
\eqref{eq:lp-kac-rice} with the density of the $k$th factorial moment measure
of the Haar zero divisor.  Haar-nullity of multiple $a$-points, already proved
for the fixed-order family, is exactly the nondegeneracy needed here.  For
$k=1$ this agrees with the one-point density used above:
\begin{equation}
 K_1^{(m,a)}(s)=G_m(\Re s,a).
 \label{eq:lp-one-point}
\end{equation}

\begin{theorem}[Full local point-process universality]
\label{thm:lp-process}
For every $m\geq0$, $a\in\C$, and unimodular twist $\omega$, one has weak
convergence on $\mathcal N_{\rm lf}(D)$:
\begin{equation}
 \boxed{
 \frac1T\int_0^T\delta_{\vartheta_t\Xi_{m,\omega,a}}\,\dd t
 \ \Longrightarrow\ 
 \operatorname{Law}(\Xi_{m,U,a}),
 \qquad U\sim\operatorname{Haar}(\mathbb H).
 }
 \label{eq:lp-process-limit}
\end{equation}
In particular, for every $k\geq1$ and every compactly supported continuous
test function $\Psi$ on $D^k$, the vertical average of the $k$-fold factorial
sum converges to
\begin{equation}
 \int_{D^k}\Psi(s_1,\ldots,s_k)
 K_k^{(m,a)}(s_1,\ldots,s_k)\prod_{j=1}^k\dd A(s_j).
 \label{eq:lp-factorial-limit}
\end{equation}
After anchoring the first ordinate, this is the relative-coordinate all-$k$
correlation law.
\end{theorem}

\begin{proof}
Unique factorization makes every nontrivial finite integer combination of the
$\log p$ nonzero.  Hence the flow $R_t$ is uniquely ergodic on the compact
group $\mathbb H$, with Haar measure as its unique invariant probability.
Thus
\[
 \frac1T\int_0^T\delta_{R_t\omega}\,\dd t
 \Longrightarrow\operatorname{Haar}(\mathbb H)
\]
for every starting point $\omega$.  Push this convergence forward under the
continuous divisor map of Lemma~\ref{lem:lp-compact-divisor}, and use
\eqref{eq:lp-covariance}; this proves \eqref{eq:lp-process-limit}.

For a compactly supported $\Psi$, the associated factorial sum is a continuous
functional of a divisor with multiplicities, including when several simple
zeros coalesce.  It is bounded by Lemma~\ref{lem:lp-compact-divisor}.  Taking
its expectation in \eqref{eq:lp-process-limit} and applying
\eqref{eq:lp-kac-rice} proves \eqref{eq:lp-factorial-limit}.
\end{proof}

The preceding theorem admits a useful exact converse: the complete process is
recovered from its correlation functions by a finite, rather than merely
formal, expansion.

\begin{theorem}[Exact Laplace functional and finite-window laws]
\label{thm:lp-laplace}
Let $f\geq0$ be continuous with compact support $K\Subset D$.  Then
\begin{equation}
 \boxed{
 {\mathbb E}\exp\{-\langle f,\Xi_{m,U,a}\rangle\}
 =\sum_{k=0}^{M_{m,a}(K)}\frac1{k!}
 \int_{K^k}\prod_{j=1}^k\bigl(e^{-f(s_j)}-1\bigr)
 K_k^{(m,a)}(\mathbf s)\prod_j\dd A(s_j).
 }
 \label{eq:lp-laplace}
\end{equation}
Let $B\Subset D$ be Borel with planar-null boundary and put
\[
 \alpha_k(B):=\int_{B^k}K_k^{(m,a)}(\mathbf s)\prod_j\dd A(s_j),
 \qquad \alpha_0(B):=1.
\]
Then, for $0\leq j\leq M_{m,a}(\overline B)$,
\begin{equation}
 \boxed{
 {\mathbb P}\{\Xi_{m,U,a}(B)=j\}
 =\sum_{k=j}^{M_{m,a}(\overline B)}
 \frac{(-1)^{k-j}}{j!(k-j)!}\,\alpha_k(B).
 }
 \label{eq:lp-count-law}
\end{equation}
In particular,
\begin{equation}
 \boxed{
 {\mathbb P}\{\Xi_{m,U,a}(B)=0\}
 =\sum_{k=0}^{M_{m,a}(\overline B)}\frac{(-1)^k}{k!}\alpha_k(B).
 }
 \label{eq:lp-void}
\end{equation}
The same right sides are the limiting vertical frequencies for every fixed
twist $\omega$.
\end{theorem}

\begin{proof}
For a divisor $\xi$ having at most $M$ points in $K$, with labels repeated
according to multiplicity,
\[
 e^{-\langle f,\xi\rangle}
 =\prod_{x\in\xi}\bigl[1+(e^{-f(x)}-1)\bigr].
\]
Expanding the product into ordered factorial sums and dividing the $k$th sum
by $k!$ gives \eqref{eq:lp-laplace}.  Likewise
\[
 z^{\xi(B)}
 =\sum_{k=0}^{M}\frac{(z-1)^k}{k!}\,(\xi(B))_k.
\]
Expectation and coefficient extraction give \eqref{eq:lp-count-law}; setting
$j=0$ gives \eqref{eq:lp-void}.  The boundary condition makes the event
$\{\xi(B)=j\}$ a continuity event for the Haar law: its only possible boundary
configurations have a zero on $\partial B$, and the expected number of those
is zero because the one-point measure has a density.  The final assertion is
therefore an application of \eqref{eq:lp-process-limit} and the Portmanteau
theorem.
\end{proof}

\subsection{The corrected almost-periodic pair structure}

For $\psi\in C_c(D)$ let
\[
 L_\psi(u):=\langle\psi,\Xi_{m,u,a}\rangle.
\]
This is continuous and bounded on $\mathbb H$.  For
$\psi_1,\psi_2\in C_c(D)$ put
\begin{align}
 A_{\psi_1,\psi_2}(h)
 &:=\int_{\mathbb H}L_{\psi_1}(u)L_{\psi_2}(R_hu)\,\dd u,
 \label{eq:lp-all-pairs}\\
 D_{\psi_1,\psi_2}(h)
 &:=\int_D\psi_1(s)\psi_2(s-ih)
 G_m(\Re s,a)\,\dd A(s),
 \label{eq:lp-diagonal-correction}
\end{align}
and let $C_{\psi_1,\psi_2}(h)$ denote the corresponding factorial, or
distinct-label, pair correlation.

\begin{proposition}[Almost periodic plus a local diagonal correction]
\label{prop:lp-apc0}
The all-pairs function $A_{\psi_1,\psi_2}$ is Bohr almost periodic,
$D_{\psi_1,\psi_2}\in C_c(\R)$, and
\begin{equation}
 \boxed{
 C_{\psi_1,\psi_2}(h)
 =A_{\psi_1,\psi_2}(h)-D_{\psi_1,\psi_2}(h)
 \in AP(\R)+C_c(\R)\subset AP(\R)+C_0(\R).
 }
 \label{eq:lp-apc0}
\end{equation}
Its Bohr mean is
\begin{equation}
 \mathcal M(C_{\psi_1,\psi_2})
 =\left(\int_D\psi_1(s)G_m(\Re s,a)\dd A(s)\right)
  \left(\int_D\psi_2(s)G_m(\Re s,a)\dd A(s)\right).
 \label{eq:lp-pair-mean}
\end{equation}
In general $C_{\psi_1,\psi_2}$ itself is not Bohr almost periodic.
\end{proposition}

\begin{proof}
The first expression in \eqref{eq:lp-all-pairs} is a convolution of continuous
functions on the compact abelian group $\mathbb H$, evaluated on the
one-parameter subgroup $R_h$.  Peter--Weyl approximation by finite character
sums makes it a uniform limit of trigonometric polynomials in $h$, hence Bohr
almost periodic.  The repeated-label terms in the product
$L_{\psi_1}(u)L_{\psi_2}(R_hu)$ are exactly
\eqref{eq:lp-diagonal-correction}, so \eqref{eq:lp-apc0} follows.  Compact
support of the two test functions makes $D$ compactly supported in $h$, and its
Ces\`aro mean is zero.  Unique ergodicity computes the mean of $A$ as the
product of the two Haar means, which are the displayed one-point integrals.
Finally, a nonzero compactly supported continuous function cannot be Bohr
almost periodic; thus subtracting the diagonal term does not in general
preserve $AP(\R)$.
\end{proof}

\subsection{Projected processes and exact Palm gap laws}

Fix a compact interval $I\Subset(1,\infty)$ and restrict the divisor to the
strip $S_I=I+i\R$.  Its ordinate projection is
\begin{equation}
 \eta_{m,u,a,I}:=\sum_{\substack{F_{m,u}(\rho)=a\\\Re\rho\in I}}
 m_\rho\,\delta_{\Im\rho}.
 \label{eq:lp-projected}
\end{equation}
Under Haar measure this is a stationary point process.  Its factorial
correlation densities are
\begin{equation}
 g_{I,k}^{(m,a)}(t_1,\ldots,t_k)
 :=\int_{I^k}K_k^{(m,a)}(\sigma_1+it_1,\ldots,
                         \sigma_k+it_k)\prod_j\dd\sigma_j,
 \label{eq:lp-projected-k}
\end{equation}
and depend only on the differences $t_j-t_1$.  Its intensity is
\begin{equation}
 c_I^{(m,a)}=g_{I,1}^{(m,a)}=\int_I G_m(\sigma,a)\,\dd\sigma.
 \label{eq:lp-intensity}
\end{equation}
Assume in what follows that $c_I^{(m,a)}>0$.

The projected Haar process is simple almost surely.  Planar multiple roots are
Haar-null, and the expected number of two distinct planar roots having exactly
the same ordinate is zero: the latter pairs lie on a Lebesgue-null hypersurface
while their factorial moment measure has the locally integrable density
$K_2^{(m,a)}$.

Let ${\mathbb P}_I^0$ denote the reduced Palm law obtained by anchoring a
typical ordinate at the origin and deleting that anchored point.  Under this
law define
\[
 G_+:=\inf\{t>0:t\in\eta\},
 \qquad
 R_{\rm ord}:=\inf\{|t|:t\in\eta\}.
\]

\begin{theorem}[Exact Palm nearest-neighbour and consecutive-gap laws]
\label{thm:lp-palm}
For every $u>0$,
\begin{align}
 {\mathbb P}_I^0\{G_+>u\}
 &=\frac1{c_I^{(m,a)}}
   \sum_{k\geq0}\frac{(-1)^k}{k!}
   \int_{(0,u)^k}g_{I,k+1}^{(m,a)}(0,t_1,\ldots,t_k)
   \prod_j\dd t_j,                                      
 \label{eq:lp-forward-gap}\\
 {\mathbb P}_I^0\{R_{\rm ord}>u\}
 &=\frac1{c_I^{(m,a)}}
   \sum_{k\geq0}\frac{(-1)^k}{k!}
   \int_{(-u,u)^k}g_{I,k+1}^{(m,a)}(0,t_1,\ldots,t_k)
   \prod_j\dd t_j.
 \label{eq:lp-two-sided-gap}
\end{align}
Both sums terminate locally.  At every $u$ at which the boundary has zero Palm
probability, the same laws are obtained by choosing an $a$-point of any fixed
twist uniformly among those with $\Re\rho\in I$ and $0<\Im\rho<T$, and then
letting $T\to\infty$.

The forward-gap distribution is absolutely continuous away from a possible
hard-core endpoint, with density
\begin{equation}
 \boxed{
 f_{+,I}^{(m,a)}(u)
 =\frac1{c_I^{(m,a)}}
  \sum_{k\geq0}\frac{(-1)^k}{k!}
  \int_{(0,u)^k}
  g_{I,k+2}^{(m,a)}(0,t_1,\ldots,t_k,u)
  \prod_j\dd t_j.
 }
 \label{eq:lp-gap-density}
\end{equation}
\end{theorem}

\begin{proof}
Campbell's formula says that the $k$th factorial moment density under the
reduced Palm law is
\[
 \frac{g_{I,k+1}^{(m,a)}(0,t_1,\ldots,t_k)}{c_I^{(m,a)}}.
\]
Apply the finite void expansion \eqref{eq:lp-void} under this Palm law to
$(0,u)$ and $(-u,u)$; this gives \eqref{eq:lp-forward-gap} and
\eqref{eq:lp-two-sided-gap}.  Differentiating the finite sum in
\eqref{eq:lp-forward-gap} and using symmetry among the $k$ integration
variables gives \eqref{eq:lp-gap-density}.

For the deterministic assertion, apply Theorem~\ref{thm:lp-process} to the
bounded local Campbell functional which sums a bounded continuous mark over
the roots in a unit-height window.  The exceptional configurations having a
root on a boundary or a multiple root form a Haar-null set.  The numerator of
the empirical Palm average converges to its Campbell expectation, while the
denominator is $c_I^{(m,a)}T+o(T)$.  Approximation from above and below handles
the gap indicators at every Palm continuity point.
\end{proof}

The unconditioned projected void law is, similarly,
\begin{equation}
 V_I(u):={\mathbb P}\{\eta_{m,U,a,I}(0,u)=0\}
 =\sum_{k\geq0}\frac{(-1)^k}{k!}
   \int_{(0,u)^k}g_{I,k}^{(m,a)}(\mathbf t)\prod_j\dd t_j.
 \label{eq:lp-projected-void}
\end{equation}

\subsection{Confluent repulsion and the planar nearest-neighbour law}

The following estimate both controls all diagonal singularities and strengthens
the pair-repulsion statement to arbitrary order.

\begin{theorem}[Vandermonde $k$-repulsion]
\label{thm:lp-vandermonde}
Let $K\Subset D$ and $k\geq2$.  There are $r_K>0$ and $C_{K,k}<\infty$ such
that whenever $z_1,\ldots,z_k\in K$ lie in a disc of radius $r_K$ and are
pairwise distinct,
\begin{equation}
 \boxed{
 K_k^{(m,a)}(z_1,\ldots,z_k)
 \leq C_{K,k}\prod_{1\leq i<j\leq k}|z_i-z_j|^2.
 }
 \label{eq:lp-vandermonde}
\end{equation}
Thus $K_k^{(m,a)}$ extends continuously by zero to the complete diagonal.  The
same argument applied separately to each collision cluster gives a locally
bounded continuous extension across every partial diagonal.  In particular,
if $k$ points are constrained to a disc of radius $r$, their $k$th factorial
mass is $O(r^{k(k+1)})$.
\end{theorem}

\begin{proof}
Let
\[
 \Delta(\mathbf z)=\prod_{i<j}(z_j-z_i).
\]
Replace the value vector $(F(z_1),\ldots,F(z_k))$ by its Newton divided
differences.  The complex determinant of this linear transformation is
$\Delta(\mathbf z)^{-1}$.  Lemma~\ref{lem:lp-jets} shows that the density of
the divided-difference vector remains uniformly bounded as the nodes merge.
Consequently the density of the value vector at $(a,\ldots,a)$ is
$O(|\Delta(\mathbf z)|^{-2})$.

On the event $F(z_j)=a$ for every $j$, factor
\[
 F(z)-a=\prod_{j=1}^k(z-z_j)H(z)
\]
in a fixed disc slightly larger than the cluster.  The family $F_{m,u}-a$ is
uniformly bounded on that larger disc.  Since its boundary stays a fixed
distance from every $z_j$, the maximum principle applied to the quotient gives
$|H(z_j)|\leq C_{K,k}$.  Hence
\[
 \prod_{j=1}^k|F'(z_j)|^2
 =\prod_{j=1}^k\left|
    H(z_j)\prod_{\ell\neq j}(z_j-z_\ell)
  \right|^2
 \ll_{K,k}|\Delta(\mathbf z)|^4.
\]
Multiplying this conditional derivative bound by the value-density bound and
using Kac--Rice proves \eqref{eq:lp-vandermonde}.  If only a proper subset of
the points coalesces, perform the divided-difference change only inside that
subset; the remaining value coordinates stay uniformly nondegenerate.  This
proves the assertion concerning partial diagonals.  Scaling
$z_j=z_0+rw_j$ gives $r^{k(k-1)}$ from the squared Vandermonde and $r^{2k}$
from planar volume, which proves the final estimate.
\end{proof}

For the exact leading pair coefficient define
\begin{equation}
 \mathcal C_m(\sigma,a)
 :={\mathbb E}\!\left[
 |F''_{m,U}(\sigma)|^4
 \delta_a(F_{m,U}(\sigma))
 \delta_0(F'_{m,U}(\sigma))
 \right].
 \label{eq:lp-confluent-C}
\end{equation}

\begin{theorem}[Confluent pair law]
\label{thm:lp-confluent}
Uniformly for $s$ in compact subsets of $D$, as $h\to0$ in $\C$,
\begin{equation}
 \boxed{
 K_2^{(m,a)}(s,s+h)
 =\frac{|h|^2}{16}\mathcal C_m(\Re s,a)+O(|h|^3).
 }
 \label{eq:lp-confluent}
\end{equation}
Moreover $\mathcal C_m(\sigma,a)>0$ if and only if $(a,0)$ is an interior
point of the support of $(F_{m,U}(\sigma),F'_{m,U}(\sigma))$.
\end{theorem}

\begin{proof}
By vertical invariance take $s=\sigma$.  Put
\[
 X=F_{m,U}(s),\qquad Y=F'_{m,U}(s),
\]
and
\[
 W_h=\frac2{h^2}\bigl(F_{m,U}(s+h)-F_{m,U}(s)-hF'_{m,U}(s)\bigr).
\]
Lemma~\ref{lem:lp-jets} gives a smooth density $q_h(x,y,w)$ for $(X,Y,W_h)$,
with $q_h\to q_0$, together with its first derivatives, where $q_0$ is the
density of $(F,F',F'')$ at $s$.  Taylor's formula gives, uniformly in $U$,
\begin{align*}
 F(s+h)&=X+hY+\tfrac12h^2W_h,\\
 F'(s+h)&=Y+hW_h+O(|h|^2).
\end{align*}
On the two-point zero slice $F(s)=F(s+h)=a$, one therefore has
\[
 Y=-\tfrac12hW_h,
 \qquad
 F'(s+h)=\tfrac12hW_h+O(|h|^2).
\]
The complex delta transformation
\[
 \delta_a(F(s+h))
 =|h|^{-2}\delta_0\bigl(Y+\tfrac12hW_h\bigr)
\]
on the slice $X=a$ yields
\begin{align*}
 K_2^{(m,a)}(s,s+h)
 &=|h|^{-2}\int_\C
   \left|\frac{hw}{2}\right|^2
   \left|\frac{hw}{2}+O(|h|^2)\right|^2
   q_h\left(a,-\frac{hw}{2},w\right)\dd A(w)\\
 &=\frac{|h|^2}{16}\int_\C|w|^4q_0(a,0,w)\dd A(w)
   +O(|h|^3),
\end{align*}
which is \eqref{eq:lp-confluent}.  The marginal density of $(F,F')$ at
$(a,0)$ is $\int_\C q_0(a,0,w)\dd A(w)$.  Smoothness and nonnegativity show
that this marginal is positive exactly when the weighted integral with
$|w|^4$ is positive.  The standard coarea description of the phase map
identifies positivity of the smooth density with the interior of its support.
\end{proof}

Let ${\mathbb P}_{I,\mathrm{pl}}^0$ be the Palm law obtained by choosing a
typical zero per unit height with real part in $I$, translating its ordinate to
zero, and retaining only zeros whose real parts also lie in $I$.  Let
$R_{\rm pl}$ be the Euclidean distance from the anchor to the nearest remaining
zero.

\begin{corollary}[Quartic planar nearest-neighbour law]
\label{cor:lp-quartic-nn}
If $c_I^{(m,a)}>0$, then
\begin{equation}
 \boxed{
 {\mathbb P}_{I,\mathrm{pl}}^0\{R_{\rm pl}\leq r\}
 =\frac{\pi r^4}{32c_I^{(m,a)}}
   \int_I\mathcal C_m(\sigma,a)\dd\sigma
   +O_I(r^5).
 }
 \label{eq:lp-quartic-nn}
\end{equation}
When the integral is positive, the planar nearest-neighbour repulsion exponent
is therefore exactly four.
\end{corollary}

\begin{proof}
The Palm expected number of other zeros within distance $r$ of the anchor is
\[
 \frac1{c_I^{(m,a)}}\int_I
 \int_{\substack{|h|<r\\\sigma+\Re h\in I}}
 K_2^{(m,a)}(\sigma,\sigma+h)\dd A(h)\dd\sigma.
\]
Theorem~\ref{thm:lp-confluent} and
$\int_{|h|<r}|h|^2\dd A(h)=\pi r^4/2$ give the displayed main term; anchors
within $r$ of an endpoint of $I$ and the cubic remainder contribute $O(r^5)$.
The difference between the expected number and the probability of at least one
neighbour is bounded by the Palm second factorial moment.  By
Theorem~\ref{thm:lp-vandermonde}, the latter is
\[
 O\!\left(\int_{|h_1|,|h_2|<r}
 |h_1h_2(h_1-h_2)|^2\dd A(h_1)\dd A(h_2)\right)=O(r^{10}),
\]
which does not alter the stated error term.
\end{proof}

\subsection{Symbolic deterministic hard cores}

We now specialize to the unweighted zero problem $m=0$, $a=0$.  The following
statements contain no numerical threshold and require no computer-assisted
claim.  Put
\begin{equation}
 B_j(\sigma):=\sum_pp^{-\sigma}(\log p)^j,
 \qquad j=1,2,
 \label{eq:lp-Bj}
\end{equation}
and define
\begin{align}
 \mathfrak f_\sigma(\lambda)
 &:=\sum_{p\geq3}p^{-\sigma}
      \left|1-\lambda\log\frac p2\right|,\qquad \lambda>0,
 \label{eq:lp-fsigma}\\
 \mathfrak m(\sigma)
 &:=\sup_{\lambda>0}
    \frac{2^{-\sigma}-\mathfrak f_\sigma(\lambda)}{\lambda}.
 \label{eq:lp-msigma}
\end{align}
Whenever $\mathfrak m(\sigma)>0$, let $d(\sigma)$ be the unique solution in
$(0,\sigma-1)$ of
\begin{equation}
 d(\sigma)B_2(\sigma-d(\sigma))=2\mathfrak m(\sigma).
 \label{eq:lp-dsigma}
\end{equation}

\begin{theorem}[Derivative and planar hard core]
\label{thm:lp-planar-core}
Let $\omega$ be any unimodular twist and let $\rho=\sigma+it$ be a zero of
$F_{0,\omega}$.  Then
\begin{equation}
 \boxed{|F'_{0,\omega}(\rho)|\geq\mathfrak m(\sigma).}
 \label{eq:lp-derivative-core}
\end{equation}
If $\mathfrak m(\sigma)>0$, the zero is simple and every other zero $\rho'$ of
the same twist satisfies
\begin{equation}
 \boxed{|\rho'-\rho|\geq d(\sigma).}
 \label{eq:lp-planar-core}
\end{equation}
Consequently $K_2^{(0,0)}(s,s')=0$ throughout the corresponding forbidden
region.
\end{theorem}

\begin{proof}
Write $U_p=\omega_pp^{-it}$, $S=\sum_{p\geq3}p^{-\sigma}U_p$, and
$Y=F'_{0,\omega}(\rho)$.  The zero equation gives $|S|=2^{-\sigma}$.  Since
\[
 Y=-\sum_{p\geq3}\log(p/2)p^{-\sigma}U_p,
\]
for every $\lambda>0$,
\[
 2^{-\sigma}=|S|
 \leq\sum_{p\geq3}p^{-\sigma}
       \left|1-\lambda\log\frac p2\right|+\lambda|Y|.
\]
Taking the supremum proves \eqref{eq:lp-derivative-core}.

Let $r=|\rho'-\rho|<\sigma-1$; otherwise
\eqref{eq:lp-planar-core} is automatic.  Taylor's integral formula along the
segment from $\rho$ to $\rho'$ and the bound
$|F''_{0,\omega}(z)|\leq B_2(\Re z)$ give
\[
 \mathfrak m(\sigma)r
 \leq\frac12r^2B_2(\sigma-r).
\]
The function $r\mapsto rB_2(\sigma-r)$ is strictly increasing, so
$r\geq d(\sigma)$.  Since no twist can realize two zeros in the open forbidden
set, its Haar factorial moment measure is zero there; continuity of $K_2$ off
the diagonal makes the density identically zero there.
\end{proof}

For unequal real parts define
\begin{align}
 \Lambda(\sigma_1,\sigma_2)
 &:={\sup}_{\lambda\geq0}\left\{
 \lambda2^{-\sigma_2}
 -\sum_{p\geq3}|p^{-\sigma_1}-\lambda p^{-\sigma_2}|
 \right\},
 \label{eq:lp-Lambda}\\
 H(\sigma_1,\sigma_2)
 &:=\frac{[\Lambda(\sigma_1,\sigma_2)-2^{-\sigma_1}]_+}
          {B_1(\sigma_1)}.
 \label{eq:lp-H}
\end{align}

\begin{theorem}[Symbolic ordinate separation]
\label{thm:lp-ordinate-core}
If $\rho_j=\sigma_j+it_j$ are two zeros of the same unimodular twist and
$\sigma_1\neq\sigma_2$, then
\begin{equation}
 \boxed{|t_2-t_1|\geq H(\sigma_1,\sigma_2).}
 \label{eq:lp-ordinate-separation}
\end{equation}
The same assertion with the indices interchanged may be used, so the maximum
of the two directed bounds is valid.
\end{theorem}

\begin{proof}
Set $h=t_2-t_1$, $U_p=\omega_pp^{-it_1}$, and $V_p=U_pp^{-ih}$.  Rotate the
second zero equation so that, for suitable $x_p\in[-1,1]$,
\[
 \sum_{p\geq3}p^{-\sigma_2}x_p=2^{-\sigma_2}.
\]
Since $|e^{ih\log p}-1|\leq|h|\log p$, the first zero equation implies
\[
 \sum_{p\geq3}p^{-\sigma_1}x_p-|h|B_1(\sigma_1)
 \leq2^{-\sigma_1}.
\]
For every $\lambda\geq0$,
\[
 \sum_{p\geq3}p^{-\sigma_1}x_p
 \geq\lambda2^{-\sigma_2}
      -\sum_{p\geq3}|p^{-\sigma_1}-\lambda p^{-\sigma_2}|.
\]
Optimizing in $\lambda$ and combining the last two inequalities proves
\eqref{eq:lp-ordinate-separation}.
\end{proof}

\begin{corollary}[A symbolic strip hard core]
\label{cor:lp-symbolic-strip-core}
Let $I\Subset(1,\infty)$ be an interval on which
\[
 d_I:=\inf_{\sigma\in I}d(\sigma)>0.
\]
Choose $0<\delta<d_I$, and suppose
\[
 H_{I,\delta}:=
 \inf_{\substack{\sigma_1,\sigma_2\in I\\
                  |\sigma_1-\sigma_2|\geq\delta}}
 \max\{H(\sigma_1,\sigma_2),H(\sigma_2,\sigma_1)\}>0.
\]
Then every two zeros in $I+i\R$ have ordinate difference at least
\begin{equation}
 h_I:=\min\left\{\sqrt{d_I^2-\delta^2},H_{I,\delta}\right\}>0.
 \label{eq:lp-strip-hcore}
\end{equation}
For $0<u<h_I$ the projected Haar process has at most one point in every
interval of length $u$, and hence
\begin{equation}
 \boxed{
 {\mathbb P}_I^0\{G_+>u\}
 ={\mathbb P}_I^0\{R_{\rm ord}>u\}=1,
 \qquad
 V_I(u)=1-c_I^{(0,0)}u.
 }
 \label{eq:lp-exact-core-laws}
\end{equation}
\end{corollary}

\begin{proof}
If $|\sigma_1-\sigma_2|<\delta$, the planar bound gives
$|t_1-t_2|\geq\sqrt{d_I^2-\delta^2}$.  Otherwise apply
Theorem~\ref{thm:lp-ordinate-core}.  Thus an interval shorter than $h_I$
contains at most one projected point.  Its expected count is
$c_I^{(0,0)}u$, so its void probability is exactly one minus that expectation;
the Palm assertions are immediate.
\end{proof}

There is a complementary no-core statement which requires no numerical
construction.

\begin{corollary}[Exact small-gap alternative]
\label{cor:lp-no-core}
For a fixed-order projected process, suppose that
\[
 g_{I,2}^{(m,a)}(0)>0.
\]
Then, as $u\downarrow0$,
\begin{align}
 {\mathbb P}_I^0\{G_+\leq u\}
 &=\frac{g_{I,2}^{(m,a)}(0)}{c_I^{(m,a)}}u+o(u),
 \label{eq:lp-small-forward}\\
 {\mathbb P}_I^0\{R_{\rm ord}\leq u\}
 &=\frac{2g_{I,2}^{(m,a)}(0)}{c_I^{(m,a)}}u+o(u),
 \label{eq:lp-small-nearest}\\
 V_I(u)
 &=1-c_I^{(m,a)}u+\frac12g_{I,2}^{(m,a)}(0)u^2+o(u^2).
 \label{eq:lp-small-void}
\end{align}
In particular the ordinate process has no exact hard core.
\end{corollary}

\begin{proof}
Insert the continuous density $g_{I,2}^{(m,a)}$ into
\eqref{eq:lp-forward-gap}, \eqref{eq:lp-two-sided-gap}, and
\eqref{eq:lp-projected-void}.  Terms involving $g_{I,k}$ with $k\geq3$ are
$O(u^2)$ in the Palm expansions and $O(u^3)$ in the unconditioned expansion,
because the confluent bounds above make these densities locally bounded.
The asserted leading terms follow.
\end{proof}

\begin{remark}[Dominant-prime version at every fixed derivative level]
\label{rem:lp-dominant-prime}
The deterministic argument is not restricted to $m=0$ or to the prime $2$.
Put $a_{m,p}(\sigma)=(\log p)^mp^{-\sigma}$ and fix a candidate dominant
prime $q$.  Then every zero of $F_{m,\omega}$ on $\Re s=\sigma$ satisfies
\[
 |F'_{m,\omega}(s)|\geq
 \sup_{\lambda\in\R\setminus\{0\}}
 \frac{a_{m,q}(\sigma)-
 \sum_{p\neq q}a_{m,p}(\sigma)
 |1-\lambda\log(p/q)|}{|\lambda|}.
\]
Whenever the right side is positive, the Taylor argument gives a planar hard
core, with $B_2$ replaced by
$\sum_p(\log p)^{m+2}p^{-\sigma}$.  This supplies symbolic hard-core
neighbourhoods at any internal dominance edge for which the displayed dual
margin is positive.  A similarly directed linear program, replacing
$p^{-\sigma_j}$ in \eqref{eq:lp-Lambda} by $a_{m,p}(\sigma_j)$ and replacing
$B_1$ by $\sum_{p\neq q}a_{m,p}(\sigma_1)|\log(p/q)|$, gives ordinate
separation.  No claim of uniformity as $m\to\infty$ is made.
\end{remark}

\subsection{What is closed and what remains open}

Theorems~\ref{thm:lp-process}--\ref{thm:lp-palm} close the qualitative fixed-
order local-statistics problem: for every fixed $m$, every target $a$, and
every unimodular twist, the full local process, every factorial correlation,
every finite-window count law, every void law, and every Palm nearest-neighbour
law exist and are universal.  Theorem~\ref{thm:lp-vandermonde} and
Corollary~\ref{cor:lp-quartic-nn} give the corresponding bulk collision law,
while Theorems~\ref{thm:lp-planar-core} and~\ref{thm:lp-ordinate-core} give the
opposite deterministic hard-core regime.

What remains open is effective rather than qualitative: a rate in the vertical
limit, estimates uniform in a derivative order that grows with height, and
closed-form evaluation of the Kac--Rice integrals.  Numerical candidate values
for hard-core transition points and computer-located multiple zeros are useful
evidence, but are not used in any theorem of this section unless accompanied by
independent interval-arithmetic certificates.
\section{Ordinate phase tomography and universal edge rigidity}
\label{sec:edge-rigidity}

The horizontal divisor laws developed above deliberately average out the coefficient
phases.  This section records the complementary fact that the 
\emph{anchored ordinates} of points approaching a right edge retain those phases with
no loss.  Combining this elementary rigidity with the deterministic hard core gives a
universal rooted edge profile, an edge-layer spacing theorem, and a stable inverse
problem.  None of the results below interprets a limiting density as a first-occurrence
height.

Throughout, let
\[
 F_c(s)=\sum_p c_p p^{-s},
 \qquad
 A_c(\sigma)=\sum_p |c_p|p^{-\sigma},
 \qquad \Re s>1,
\]
where $c=(c_p)_p$ is bounded.  The same symbol $A_c$ will denote the
holomorphic Dirichlet series $\sum_p|c_p|p^{-s}$ when its argument is
complex.  We write $\mathcal S_c=\{p:c_p\ne0\}$.

\subsection{The exact phase-defect identities}

The following identity is a quantitative equality case of the triangle inequality.
It is the basic input for every inverse statement in this section.

\begin{theorem}[Exact phase defect at a nonzero target]
\label{thm:exact-phase-defect}
Let $a\ne0$, put $\alpha=a/|a|$, and suppose that
$\rho=\sigma+it$ is an $a$-point of $F_c$.  Define
\[
 \eps_a(\rho)=A_c(\sigma)-|a|\ge0,
 \qquad
 \zeta_p(\rho)=\frac{c_p}{|c_p|}\,\overline\alpha\,p^{-it}
 \quad(p\in\mathcal S_c).
\]
Then
\begin{equation}
 \boxed{
 \sum_{p\in\mathcal S_c}|c_p|p^{-\sigma}
       |1-\zeta_p(\rho)|^2=2\eps_a(\rho).}
 \label{eq:exact-phase-defect}
\end{equation}
Consequently, for every $p\in\mathcal S_c$,
\begin{equation}
 \left|\frac{c_p}{|c_p|}-\alpha p^{it}\right|
 \le
 \left(\frac{2\eps_a(\rho)}{|c_p|p^{-\sigma}}\right)^{1/2}.
 \label{eq:coordinate-phase-defect}
\end{equation}
\end{theorem}

\begin{proof}
The equation $F_c(\rho)=a$ becomes
\[
 \sum_{p\in\mathcal S_c}|c_p|p^{-\sigma}\zeta_p=|a|.
\]
Taking real parts and using $|1-\zeta|^2=2(1-\Re\zeta)$ on the
unit circle proves \eqref{eq:exact-phase-defect}.  Dropping every term but
the $p$th gives \eqref{eq:coordinate-phase-defect}.
\end{proof}

Suppose that $|a|$ lies in the range of $A_c$, and put
\[
 E_c(a)=A_c^{-1}(|a|).
\]
The right-edge theorem gives $a$-points with real parts tending to
$E_c(a)$ from the left.  The preceding equality immediately makes their
ordinates inverse-complete.

\begin{corollary}[Fixed-target recovery]
\label{cor:fixed-target-recovery}
If $\rho_n=\sigma_n+it_n$ are $a$-points and
$\sigma_n\to E_c(a)$, then, for every $p\in\mathcal S_c$,
\begin{equation}
 \boxed{
 \frac{c_p}{|c_p|}=\frac a{|a|}\lim_{n\to\infty}p^{it_n}.}
 \label{eq:fixed-target-recovery}
\end{equation}
In particular, the following two pieces of data determine the complete
coefficient system $c$:
\begin{enumerate}[label=\textup{(\roman*)}]
 \item edge samples $(r_j,E_c(r_j))$ for distinct $r_j$ having one
       accumulation point in the interior of the edge range;
 \item one sequence of points of one fixed nonzero target tending to its
       right edge.
\end{enumerate}
\end{corollary}

\begin{proof}
As $\sigma_n\to E_c(a)$, continuity gives
$\eps_a(\rho_n)\to0$, so \eqref{eq:fixed-target-recovery} follows from
\eqref{eq:coordinate-phase-defect}.  The sparse edge samples recover
$A_c$ and hence every $|c_p|$ by
\Cref{cor:sparse-edge}; the displayed limit then recovers the phases.
\end{proof}

For zeros of a unimodular twist, the least-prime term is separated from
the tail.  Let
\begin{equation}
 S=\sigma_{\rm e},
 \qquad
 D(\sigma)=\sum_{p\ge3}p^{-\sigma}-2^{-\sigma},
 \label{eq:zero-defect}
\end{equation}
so $D(S)=0$ and $D(\sigma)>0$ for $1<\sigma<S$.  Recall that
\[
 2^{-S}=\sum_{p\ge3}p^{-S}.
\]

\begin{corollary}[Exact zero defect and relative phase recovery]
\label{cor:zero-phase-defect}
Let $F_\omega(s)=\sum_p\omega_pp^{-s}$ with $|\omega_p|=1$, and let
$\rho=\sigma+it$ be a zero.  Then
\begin{equation}
 \boxed{
 \sum_{p\ge3}p^{-\sigma}
 \left|\frac{\omega_p}{\omega_2}
       +\left(\frac p2\right)^{it}\right|^2
 =2D(\sigma).}
 \label{eq:exact-zero-phase-defect}
\end{equation}
Hence
\begin{equation}
 \left|\frac{\omega_p}{\omega_2}
       +\left(\frac p2\right)^{it}\right|
 \le \sqrt{2D(\sigma)p^\sigma}
 \qquad(p\ge3).
 \label{eq:zero-coordinate-bound}
\end{equation}
If $\rho_n=\sigma_n+it_n$ and $\sigma_n\to S$, then
\begin{equation}
 \boxed{
 \frac{\omega_p}{\omega_2}
 =-\lim_{n\to\infty}\left(\frac p2\right)^{it_n}}
 \qquad(p\ge3).
 \label{eq:zero-relative-recovery}
\end{equation}
Thus one fixed-level zero sequence approaching the edge recovers the
whole twist modulo its unavoidable common unimodular factor.
\end{corollary}

\begin{proof}
After division of $F_\omega(\rho)=0$ by
$\omega_2 2^{-it}$, the tail has value $-2^{-\sigma}$.  Expanding the
left side of \eqref{eq:exact-zero-phase-defect} therefore gives
\[
 2\sum_{p\ge3}p^{-\sigma}-2\cdot2^{-\sigma}=2D(\sigma).
\]
The remaining assertions follow by dropping terms and letting
$\sigma\uparrow S$.
\end{proof}

\subsection{Universal rooted profiles}

The phase defect controls not merely one coefficient at a time but the
entire analytic function rooted at a near-edge point.

\begin{theorem}[Rooted profile at a nonzero target]
\label{thm:general-rooted-profile}
Fix the coefficient moduli $(|c_p|)_p$, let $a\ne0$, and put
$E=E_c(a)$ and $\alpha=a/|a|$.  The coefficient phases may vary with
$n$.  If $\rho_n=\sigma_n+it_n$ is an $a$-point and
$\sigma_n\to E$, then
\begin{equation}
 \overline\alpha F_{c^{(n)}}(\rho_n+z)-|a|
 \longrightarrow
 H_{c,a}(z):=A_c(E+z)-A_c(E)
 \label{eq:general-rooted-limit}
\end{equation}
locally uniformly on $\Re z>1-E$.  Consequently the translated
$a$-point divisors converge locally vaguely, with multiplicity, to the
zero divisor of $H_{c,a}$.

Moreover,
\[
 H_{c,a}(0)=0,
 \qquad
 H_{c,a}'(0)=A_c'(E)<0,
\]
and all normalized jets at the rooted points converge to the
corresponding jets of $H_{c,a}$.  If the support contains two distinct
primes, then $0$ is the only zero of $H_{c,a}$ on the imaginary axis.
\end{theorem}

\begin{proof}
For each supported prime, \Cref{thm:exact-phase-defect} gives
\[
 \overline\alpha\frac{c_p^{(n)}}{|c_p|}p^{-it_n}\longrightarrow1.
\]
On a compact subset of $\Re z>1-E$, the summands are dominated, for all
large $n$, by a fixed convergent series $C p^{-1-\delta}$.  Dominated
convergence proves \eqref{eq:general-rooted-limit}; convergence of all
derivatives follows either in the same way on a slightly smaller
half-plane or from Cauchy's formula.  Hurwitz's theorem, applied on
compact sets whose boundaries avoid the limiting divisor, gives divisor
convergence.

The derivative formula is immediate.  If $H_{c,a}(iy)=0$, equality in
the triangle inequality forces $p^{-iy}=1$ for every supported prime.
For two distinct primes this is impossible when $y\ne0$, by unique
factorization.
\end{proof}

The zero edge has a different canonical profile because one term opposes
the entire tail.

\begin{theorem}[Universal rooted zero-edge profile]
\label{thm:rooted-zero-edge-profile}
Allow both the twist $\omega^{(n)}$ and the zero to vary with $n$.  Let
\[
 \rho_n=\sigma_n+it_n,
 \qquad \sigma_n\to S,
 \qquad U_{2,n}=\omega^{(n)}_2\,2^{-it_n}.
\]
Then
\begin{equation}
 G_n(z):=\frac{F_{\omega^{(n)}}(\rho_n+z)}{U_{2,n}}
 \longrightarrow
 H_{\rm e}(z):=2^{1-(S+z)}-\cP(S+z)
 \label{eq:rooted-zero-limit}
\end{equation}
locally uniformly on
\begin{equation}
 \Omega_{\rm e}=\{z:\Re z>1-S\}.
 \label{eq:rooted-domain}
\end{equation}
The convergence is uniform over all twists and all zeros with
$D(\Re\rho)\le\eps$ as $\eps\downarrow0$.

Write
\begin{equation}
 \kappa_{\rm e}
 =\sum_{p\ge3}\log(p/2)p^{-S}=-D'(S)>0.
 \label{eq:kappa-edge}
\end{equation}
Then
\begin{equation}
 H_{\rm e}(0)=0,
 \qquad H_{\rm e}'(0)=\kappa_{\rm e},
 \qquad
 \frac{F_{\omega^{(n)}}'(\rho_n)}{U_{2,n}}\longrightarrow\kappa_{\rm e},
 \label{eq:universal-edge-slope}
\end{equation}
and the analogous assertion holds for every higher derivative.

More quantitatively, let $K\Subset\Omega_{\rm e}$ and put
$s_K=\min_{z\in K}\Re(S+z)>1$.  If
$\eps_n=D(\sigma_n)$, then
\begin{equation}
 \sup_{z\in K}|G_n(z)-H_{\rm e}(z)|
 \ll_K
 \begin{cases}
  \eps_n^{(s_K-1)/(S-1)},&1<s_K<(S+1)/2,\\
  \sqrt{\eps_n\log(1/\eps_n)},&s_K=(S+1)/2,\\
  \sqrt{\eps_n},&s_K>(S+1)/2.
 \end{cases}
 \label{eq:rooted-profile-rate}
\end{equation}
At either transition one may instead use an arbitrarily small loss in
the exponent.  In particular, an $O_K(\sqrt\eps)$ estimate on the whole
of \eqref{eq:rooted-domain} is not asserted.
\end{theorem}

\begin{proof}
Put
\[
 r_{p,n}=\frac{\omega^{(n)}_p p^{-it_n}}
                 {\omega^{(n)}_2 2^{-it_n}}
 \qquad(p\ge3).
\]
Then \eqref{eq:exact-zero-phase-defect} is equivalently
\begin{equation}
 \sum_{p\ge3}p^{-\sigma_n}|r_{p,n}+1|^2=2\eps_n.
 \label{eq:weighted-l2-defect}
\end{equation}
For each fixed $p$, this implies $r_{p,n}\to-1$.  Since
$|r_{p,n}+1|\le2$, dominated convergence gives
\eqref{eq:rooted-zero-limit} on every compact subset of
$\Omega_{\rm e}$.  It also shows that the convergence is uniform in the
choice of twist once only the defect is prescribed.

For completeness, the rate follows from a prime split.  If
$s=\min_{z\in K}\Re(\sigma_n+z)$ and $X\ge3$, then
\begin{align}
 \sum_{p\ge3}|r_{p,n}+1|p^{-s}
 &\le
 \sqrt{2\eps_n}
 \left(\sum_{p\le X}p^{\sigma_n-2s}\right)^{1/2}
 +2\sum_{p>X}p^{-s}.                         \label{eq:rooted-split}
\end{align}
Bounding the prime sums by the corresponding integer sums and taking
$X=\eps_n^{-1/(S-1)}$ gives the three cases of
\eqref{eq:rooted-profile-rate}.  Also
\[
 S-\sigma_n=\frac{\eps_n}{\kappa_{\rm e}}+O(\eps_n^2),
\]
so replacing $\sigma_n$ by $S$ contributes only $O_K(\eps_n)$.
The derivative assertions follow from local uniform convergence or
directly from \eqref{eq:weighted-l2-defect}.
\end{proof}

\begin{corollary}[Divisor and rooted Palm collapse]
\label{cor:rooted-palm-collapse}
Let
\[
 \mathscr D_{\omega,\rho}
 =\sum_{F_\omega(\rho')=0}m_{\rho'}\,\delta_{\rho'-\rho}
\]
be the zero divisor viewed from the zero $\rho$.  Uniformly over all
twists and roots with $D(\Re\rho)\le\eps$,
\begin{equation}
 \mathscr D_{\omega,\rho}
 \xrightarrow[\eps\downarrow0]{\rm vague}
 \mathscr D_{\rm e}:=\operatorname{div}(H_{\rm e})
 \quad\hbox{on }\Omega_{\rm e}.
 \label{eq:rooted-divisor-collapse}
\end{equation}
Consequently, any Palm sampling scheme supported on roots in the layer
$D(\Re\rho)\le\eps$, after recentering at its sampled root, converges in
law to the point mass at the deterministic divisor $\mathscr D_{\rm e}$.
\end{corollary}

\begin{proof}
This is Hurwitz's theorem applied to the uniform compact convergence in
\Cref{thm:rooted-zero-edge-profile}.  The Palm statement is immediate
because the convergence is uniform over every admissible root, not only
almost sure with respect to a particular sampling measure.
\end{proof}

The canonical function has a one-sided boundary geometry.

\begin{proposition}[Zero-free face of the rooted profile]
\label{prop:rooted-profile-zero-free-face}
The zero at $0$ of $H_{\rm e}$ is simple, and
\begin{equation}
 H_{\rm e}(z)\ne0
 \qquad(\Re z\ge0,\ z\ne0).
 \label{eq:rooted-right-face}
\end{equation}
If $d_{\rm hc}(\sigma)$ is any valid deterministic planar hard-core
radius for zeros near $S$, continuous from the left at $S$, then
\begin{equation}
 H_{\rm e}(z)\ne0
 \qquad(0<|z|<d_{\rm hc}(S)).
 \label{eq:rooted-hard-core}
\end{equation}
\end{proposition}

\begin{proof}
Simplicity follows from \eqref{eq:kappa-edge}.  The open half-plane in
\eqref{eq:rooted-right-face} is zero-free by the definition of the
projection edge.  If $H_{\rm e}(iy)=0$, equality in the polygon
inequality forces
\[
 e^{-iy\log(p/2)}=1\qquad(p\ge3).
\]
Taking $p=3,5$ and using unique factorization gives $y=0$.

If \eqref{eq:rooted-hard-core} failed, Hurwitz's theorem would produce,
near a nonzero limiting root $z$, a second zero of $F_\omega$ at
distance strictly smaller than $d_{\rm hc}(\sigma_n)$ from the root
$\rho_n$, a contradiction.
\end{proof}

\subsection{Exact boundary attainment and orbit classification}

Right-edge projection is phase blind, but actual attainment of the
edge is maximally phase rigid.

\begin{theorem}[Classification of boundary-attaining phases]
\label{thm:boundary-attaining-orbits}
Let $a=|a|\alpha\ne0$ and $E=A_c^{-1}(|a|)$.  Then
\begin{equation}
 F_c(E+it)=a
 \quad\Longleftrightarrow\quad
 \frac{c_p}{|c_p|}=\alpha p^{it}
 \quad\text{for every }p\in\mathcal S_c.
 \label{eq:a-boundary-classification}
\end{equation}
If $\mathcal S_c$ contains two distinct primes, the boundary ordinate
$t$ is unique.

For a unimodular prime twist,
\begin{equation}
 F_\omega(S+it)=0
 \quad\Longleftrightarrow\quad
 \frac{\omega_p}{\omega_2}
 =-\left(\frac p2\right)^{it}
 \quad(p\ge3).
 \label{eq:zero-boundary-classification}
\end{equation}
The boundary ordinate is unique.  If
\[
 \omega_2^*=1,
 \qquad \omega_p^*=-1\quad(p\ge3),
\]
then the complete family of boundary-attaining zero twists is
\begin{equation}
 \omega_p=\kappa\,\omega_p^*p^{it},
 \qquad |\kappa|=1,\quad t\in\R.
 \label{eq:boundary-orbit}
\end{equation}
This family is dense in the coefficient torus in the product topology
and has Haar measure zero.
\end{theorem}

\begin{proof}
At the boundary the sum of the moduli equals the modulus of the sum, so
all summands must have the target phase.  This proves
\eqref{eq:a-boundary-classification}.  Equality at the zero edge forces
every tail summand to have the phase opposite to the $p=2$ summand,
which gives \eqref{eq:zero-boundary-classification}.  Two different
ordinates in the zero case would imply simultaneously
\[
 (3/2)^{i(t-t')}=(5/2)^{i(t-t')}=1;
\]
unique factorization forces $t=t'$.  The nonzero-target case is the
same argument with any two supported primes.

Equation \eqref{eq:boundary-orbit} is a rearrangement of
\eqref{eq:zero-boundary-classification}.  Density follows from
Kronecker's theorem.  Haar nullity already follows after projection to
the coordinates $2,3,5$: on each bounded $t$-interval the image of
$(\kappa,t)$ is a two-dimensional Lipschitz image in the three-torus,
and the whole image is a countable union of such null sets.
\end{proof}

\subsection{An edge-layer spacing law}

The fixed hard core becomes arbitrarily large in the vertical direction
when both roots are required to lie in a shrinking edge layer.

\begin{theorem}[Divergent ordinate exclusion]
\label{thm:divergent-ordinate-exclusion}
There are constants $\eps_*,c,C>0$, independent of the twist, such that
whenever $0<\eps<\eps_*$ and $\rho_j=\sigma_j+it_j$ are distinct zeros
of the same unimodular twist satisfying
\[
 D(\sigma_1)\le\eps,
 \qquad D(\sigma_2)\le\eps,
\]
one has
\begin{equation}
 \boxed{
 |t_1-t_2|\ge c\log(1/\eps)-C.}
 \label{eq:divergent-gap}
\end{equation}
One may take
\begin{equation}
 c=\frac{\pi}
 {\bigl(\log(3/2)+\log(5/2)\bigr)\log 10},
 \label{eq:elementary-gap-constant}
\end{equation}
after enlarging $C$ and decreasing $\eps_*$.
In particular, the minimum ordinate spacing in the defect-$\eps$ layer
tends to infinity as $\eps\downarrow0$, uniformly over all twists.
\end{theorem}

\begin{proof}
The condition $D(\sigma_j)\le\eps$ forces $\sigma_j\to S$ uniformly as
$\eps\downarrow0$.  The rooted convergence to a function having a
simple zero at $0$ gives a uniform constant $h_0>0$ such that
$|t_1-t_2|\ge h_0$ for all sufficiently small $\eps$.  Equivalently,
one may combine the planar hard core near $S$ with
$|\sigma_1-\sigma_2|\to0$.

Put $h=|t_1-t_2|$, $a_0=\log(3/2)$, and $b_0=\log(5/2)$.  Applying
\eqref{eq:zero-coordinate-bound} at the two roots gives
\begin{equation}
 |e^{iha_0}-1|\le2\sqrt{2\eps\,3^S},
 \qquad
 |e^{ihb_0}-1|\le2\sqrt{2\eps\,5^S}.
 \label{eq:two-prime-recurrence}
\end{equation}
For small $\eps$, choose $m,n\in\mathbb Z$ such that
\begin{align}
 |ha_0-2\pi m|&\le\pi\sqrt{2\,3^S\eps},
 &|hb_0-2\pi n|&\le\pi\sqrt{2\,5^S\eps}.
 \label{eq:integer-phase-approximants}
\end{align}
The fixed lower bound $h_0$ ensures that $(m,n)\ne(0,0)$ after a further
decrease of $\eps_*$.  Hence
\begin{equation}
 |mb_0-na_0|\le C_0\sqrt\eps,
 \qquad
 C_0=\frac{b_0\sqrt{2\,3^S}+a_0\sqrt{2\,5^S}}2.
 \label{eq:log-linear-upper}
\end{equation}
On the other hand,
\[
 mb_0-na_0
 =\log\left(\frac{(5/2)^m}{(3/2)^n}\right)\ne0.
\]
Writing the rational number in lowest terms, its numerator and
denominator are at most $10^M$, where $M=|m|+|n|$.  Thus
\begin{equation}
 |mb_0-na_0|\ge\frac1{10^M+1}\ge\frac1{2\cdot10^M}.
 \label{eq:elementary-log-lower}
\end{equation}
Moreover,
\[
 M\le Ah+1,
 \qquad
 A=\frac{a_0+b_0}{2\pi}.
\]
Combining \eqref{eq:log-linear-upper} and
\eqref{eq:elementary-log-lower}, then absorbing fixed terms into $C$,
gives \eqref{eq:divergent-gap} with
\eqref{eq:elementary-gap-constant}.
\end{proof}

\begin{remark}
The logarithmic rate is elementary and intentionally conservative.
An effective lower bound for the nonzero linear form
$m\log(5/2)-n\log(3/2)$ from the theory of logarithms of algebraic
numbers would give a power-law exclusion with a very small effective
exponent.  No such external strengthening is needed here.
\end{remark}

\subsection{Edge-germ rigidity and stable matching}

Two zero sets have the same \emph{edge germ} if they agree in some strip
$S-\eta<\Re s\le S$.  Equality after an imaginary translation is
defined analogously.

\begin{theorem}[Edge-germ orbit rigidity]
\label{thm:edge-germ-rigidity}
For two unimodular twists $\omega,\omega'$ and $\tau\in\R$, the
following are equivalent:
\begin{enumerate}[label=\textup{(\roman*)}]
 \item the zero edge germ of $F_{\omega'}$ is the translate by $i\tau$
       of the zero edge germ of $F_\omega$;
 \item there exists $\kappa\in\mathbb T$ such that
       \begin{equation}
        \omega'_p=\kappa\omega_p p^{i\tau}
        \qquad\text{for every prime }p.
        \label{eq:zero-germ-orbit}
       \end{equation}
\end{enumerate}
When these conditions hold,
\[
 F_{\omega'}(s)=\kappa F_\omega(s-i\tau),
\]
so the equality extends from the edge germs to the full zero sets.
In particular, an anchored zero edge germ determines the twist modulo
one common unimodular factor.

More generally, let $|a|=|b|\ne0$, and let $c,d$ have the same
coefficient moduli.  If the $b$-point edge germ of $F_d$ is the translate
by $i\tau$ of the $a$-point edge germ of $F_c$, then
\begin{equation}
 \frac{d_p}{|d_p|}
 =\frac ba\frac{c_p}{|c_p|}p^{i\tau}
 \qquad(p\in\mathcal S_c).
 \label{eq:a-germ-orbit}
\end{equation}
The converse holds and gives equality of the complete translated
divisors.
\end{theorem}

\begin{proof}
Choose zeros $\sigma_n+it_n$ in the common germ with
$\sigma_n\to S$.  Their translates have ordinates $t_n+\tau$.
Applying \eqref{eq:zero-relative-recovery} to both twists gives
\[
 \frac{\omega'_p}{\omega'_2}
 =\frac{\omega_p}{\omega_2}
  \left(\frac p2\right)^{i\tau}.
\]
Taking $\kappa=(\omega'_2/\omega_2)2^{-i\tau}$ proves
\eqref{eq:zero-germ-orbit}.  The converse follows by direct
substitution.  The proof of \eqref{eq:a-germ-orbit} is identical, using
\eqref{eq:fixed-target-recovery}; the factor $b/a$ fixes the global
phase.
\end{proof}

The reconstruction is stable on every finite collection of prime
coordinates.  The hard core supplies the otherwise missing canonical
matching of noisy point clouds.

\begin{proposition}[Canonical noisy matching and inverse stability]
\label{prop:noisy-phase-stability}
Let $J\Subset(1,S]$ lie in a region where every zero of every twist has
planar separation at least $d_J>0$.  Suppose two locally finite zero
patterns in $J+i\R$ have Hausdorff distance at most $\eta<d_J/2$ and
ignore points within distance $\eta$ of the boundary of the observation
window.  Then proximity within $\eta$ defines a unique bijective
matching.

If $\rho=\sigma+it$ and $\rho'=\sigma'+it'$ are matched zeros of twists
$\omega$ and $\omega'$, respectively, then for every $p\ge3$,
\begin{align}
 \left|\frac{\omega_p}{\omega_2}
       -\frac{\omega'_p}{\omega'_2}\right|
 &\le \sqrt{2D(\sigma)p^\sigma}
      +\sqrt{2D(\sigma')p^{\sigma'}}
      +2\left|\sin\frac{(t-t')\log(p/2)}2\right|             \notag\\
 &\le \sqrt{2D(\sigma)p^\sigma}
      +\sqrt{2D(\sigma')p^{\sigma'}}
      +\min\{2,\eta\log(p/2)\}.
 \label{eq:noisy-zero-inverse}
\end{align}
For nonzero targets, after normalizing coefficient phases by the target
phase, the same formula holds with the two defects from
\Cref{thm:exact-phase-defect}, the weights $|c_p|p^{-\sigma}$, and
$\log p$ in place of $\log(p/2)$.
\end{proposition}

\begin{proof}
Two candidates for the same matched point would be at mutual distance
at most $2\eta<d_J$, contrary to separation.  Mutual Hausdorff
approximation then makes the matching bijective.  The first two terms of
\eqref{eq:noisy-zero-inverse} are
\eqref{eq:zero-coordinate-bound}; the last is the chordal distance
between $(p/2)^{it}$ and $(p/2)^{it'}$.
\end{proof}

There is also a useful forward stability statement.  Let $m_*(\sigma)>0$
be a deterministic lower bound for $|F_\omega'(\rho)|$ at a zero with
real part $\sigma$, and let $d_{\rm hc}(\sigma)>0$ be its hard-core
radius.

\begin{proposition}[Forward root stability]
\label{prop:forward-root-stability}
Let $F_\omega(\rho)=0$, $\sigma=\Re\rho$, and choose $R>0$ so that
\begin{equation}
 R<\min\{d_{\rm hc}(\sigma),\sigma-1\},
 \qquad
 R\cP''(\sigma-R)\le m_*(\sigma).
 \label{eq:root-stability-radius}
\end{equation}
For another twist $\nu$, put
\[
 \delta_R=\sum_p|\omega_p-\nu_p|p^{-(\sigma-R)}.
\]
If $\delta_R<m_*(\sigma)R/2$, then $F_\nu$ has exactly one zero
$\rho_\nu$ in $|s-\rho|<R$, counted with multiplicity, and
\begin{equation}
 |\rho_\nu-\rho|\le\frac{2\delta_R}{m_*(\sigma)}.
 \label{eq:forward-root-lipschitz}
\end{equation}
\end{proposition}

\begin{proof}
On $|z-\rho|=R$, Taylor's formula gives
\[
 |F_\omega(z)|
 \ge m_*(\sigma)R-\frac12\cP''(\sigma-R)R^2
 \ge\frac12m_*(\sigma)R.
\]
The hard core says that $\rho$ is the only zero of $F_\omega$ in the
disc.  Rouch\'e's theorem gives one zero of $F_\nu$.  Applying the same
Taylor bound at that zero, now with its distance from $\rho$ in place of
$R$, proves \eqref{eq:forward-root-lipschitz}.
\end{proof}

\subsection{Phase-readable Delone divisors}

The preceding results expose a sharp distinction between anchored and
translation-averaged data.

\begin{corollary}[Delone geometry and phase readability]
\label{cor:phase-readable-delone}
Fix a twist $\omega$ and an edge strip
\[
 S-\eta<\Re s<S
\]
whose closure lies in the deterministic hard-core region.  Above a
sufficiently large ordinate, its zero divisor is a Delone set in the
strip: it is uniformly discrete by the hard core and relatively dense
by \Cref{cor:edge-recurrence}.  Its anchored edge germ determines
$\omega$ modulo global phase, and its germ modulo imaginary translation
determines the orbit of $\omega$ modulo global phase and the vertical
Kronecker flow.

Nevertheless, every translation-averaged finite-order correlation law
of these divisors is the same for all unimodular twists.  Thus the family
is stationary-homometric at every finite order but anchored-rigid at the
edge.
\end{corollary}

\begin{proof}
Uniform discreteness and relative density are precisely the hard-core
and bounded-window recurrence conclusions.  Phase readability is
\Cref{thm:edge-germ-rigidity}; phase blindness of the averaged laws is
the joint Kac--Rice universality established earlier.
\end{proof}

\begin{remark}[Logical scope]
The exact defect identities, rooted limits, orbit classifications, and
spacing theorem are deterministic.  They do not require a numerical
value for the onset of the hard-core region.  Any decimal threshold for
that onset should be quoted only after interval certification.  Likewise,
an edge density describes a limiting frequency after the height limit;
without a uniform moving-edge error term it does not determine the first
height at which a prescribed phase accuracy occurs.
\end{remark}
\section{Phase distortion and its exact zero-edge cost}
\label{sec:phase-distortion}

The horizontal divisor law is phase blind, whereas the ordinate of a zero
contains phase information.  In this section we make that contrast
quantitative.  The natural loss function is not introduced ad hoc: at every
zero it is exactly twice the polygonal defect from the outer edge.  Combining
this identity with the logarithmic edge profile gives a universal small-loss
law for every unimodular prime twist.

Let
\[
 F_\omega(s)=\sum_{p\in\PP}\omega_p p^{-s},
 \qquad |\omega_p|=1,
\]
and let $\sigma_{\rm e}$ be the unique solution of
\[
 \sum_{p\geq3}p^{-\sigma_{\rm e}}=2^{-\sigma_{\rm e}}.
\]
For $1<\sigma<\sigma_{\rm e}$ put
\begin{equation}
\label{eq:zero-edge-defect}
 \eps_0(\sigma)
 :=\sum_{p\geq3}p^{-\sigma}-2^{-\sigma}.
\end{equation}
Thus $\eps_0(\sigma)>0$, $\eps_0(\sigma)\downarrow0$ as
$\sigma\uparrow\sigma_{\rm e}$, and
$\eps_0'(\sigma_{\rm e})=D_1(\sigma_{\rm e})<0$ in the notation of
\Cref{sec:exact-edge-profile}.

\begin{definition}[Relative-phase decoder and distortion]
\label{def:relative-phase-distortion}
If $\rho=\sigma+it$ is a zero of $F_\omega$, use the phase of the coefficient
at $2$ as reference and set
\[
 \widehat\vartheta_p(t):=-\left(\frac p2\right)^{it},
 \qquad p\geq3.
\]
The associated weighted chordal distortion is
\begin{equation}
\label{eq:relative-phase-distortion}
 \mathcal D_\omega(\rho)
 :=\sum_{p\geq3}p^{-\sigma}
 \left|\frac{\omega_p}{\omega_2}-\widehat\vartheta_p(t)\right|^2
 =\sum_{p\geq3}p^{-\sigma}
 \left|\frac{\omega_p}{\omega_2}+\left(\frac p2\right)^{it}\right|^2.
\end{equation}
\end{definition}

\begin{proposition}[Exact phase-defect identity]
\label{prop:distortion-defect}
Every zero $\rho=\sigma+it$ of every unimodular twist satisfies
\begin{equation}
\label{eq:distortion-defect}
 \boxed{\mathcal D_\omega(\rho)=2\eps_0(\sigma).}
\end{equation}
Consequently, for every $p\geq3$,
\begin{equation}
\label{eq:coordinate-phase-bound}
 \left|\frac{\omega_p}{\omega_2}
       +\left(\frac p2\right)^{it}\right|
 \leq \sqrt{2\eps_0(\sigma)p^\sigma}.
\end{equation}
In particular, any sequence of zeros with real parts tending to
$\sigma_{\rm e}$ recovers every relative coefficient phase
$\omega_p/\omega_2$.
\end{proposition}

\begin{proof}
This is \Cref{cor:zero-phase-defect}, rewritten in the notation of
\eqref{eq:relative-phase-distortion}.  The coordinatewise estimate follows
by discarding every nonnegative summand except the one indexed by $p$.
\end{proof}

The identity allows the zero divisor itself to be marked by an intrinsic
phase-reconstruction error.  For sufficiently small $d>0$, let $\sigma(d)$
be the unique number near $\sigma_{\rm e}$ such that
\begin{equation}
\label{eq:sigma-of-distortion}
 2\eps_0(\sigma(d))=d,
\end{equation}
and define
\begin{equation}
\label{eq:distortion-intensity}
 \mathcal C(d):=\int_{\sigma(d)}^{\sigma_{\rm e}}G(\sigma,0)\,\dd\sigma.
\end{equation}
Here and below an isolated zero on the boundary line has no effect on a
vertical density.

\begin{theorem}[Universal small-distortion law]
\label{thm:small-distortion-law}
For $\bullet\in\{\mathrm{mult},\mathrm{dist},\mathrm{simp}\}$ let
\[
 M_\omega^\bullet(d;T)
 :=\#_\bullet\{\rho:F_\omega(\rho)=0,\ 0<\Im\rho<T,
                         \ \mathcal D_\omega(\rho)\leq d\}.
\]
For every fixed sufficiently small $d>0$,
\begin{equation}
\label{eq:distortion-counting-law}
 \boxed{
 \lim_{T\to\infty}\frac{M_\omega^\bullet(d;T)}{T}=\mathcal C(d),
 }
\end{equation}
and the limit is independent of $\omega$ and of the counting convention.
By the universal rooted edge profile, after decreasing the allowed range of
$d$, every zero counted here has a twist-independent neighbourhood in which it
is the unique zero and is simple.

Let
\[
 \beta_{\rm e}=\frac1{\sigma_{\rm e}-1},
 \qquad \ell_d=\log\frac2d,
\]
and let $K_{\rm e}$ and $\mathscr I_{\rm e}$ be the constants and rate
function of \Cref{thm:exact-edge-profile}.  Then
\begin{equation}
\label{eq:integrated-rate-sharp}
 \boxed{
 -\log\mathcal C(d)
 =\mathscr I_{\rm e}(\sigma(d))
  +O\!\left(\mathscr I_{\rm e}(\sigma(d))^{1/2}\right),
 }
\end{equation}
and hence
\begin{equation}
\label{eq:integrated-distortion-profile}
 \boxed{
 -\log\mathcal C(d)
 =K_{\rm e}\left(\frac d2\right)^{-\beta_{\rm e}}
  \ell_d^{-(1+\beta_{\rm e})}
 \left(
  1+(1+\beta_{\rm e})\frac{\log\ell_d}{\ell_d}
  +O(\ell_d^{-1})
 \right).
 }
\end{equation}
\end{theorem}

\begin{proof}
By \Cref{prop:distortion-defect}, the condition
$\mathcal D_\omega(\rho)\leq d$ is precisely
$\Re\rho\in[\sigma(d),\sigma_{\rm e}]$.  The phase-blind horizontal
divisor theorem, including its distinct- and simple-point forms, therefore
gives \eqref{eq:distortion-counting-law}.

For the asserted pointwise isolation, choose a disc about $0$ whose closure
lies in the domain of \Cref{thm:rooted-zero-edge-profile} and contains no
other zero of $H_{\rm e}$.  Uniform rooted convergence and Rouch\'e's theorem
then give, for all sufficiently small $d$, exactly one simple zero in that
disc---the rooted zero itself---uniformly in the twist.

It remains to check that integration in $\sigma$ does not alter the exact
logarithmic edge scale.  Write $u=\eps_0(\sigma)$,
$x=d/2$, $\sigma=\sigma(u)$, and
\[
 R(u):=\mathscr I_{\rm e}(\sigma(u)).
\]
Since $\eps_0'(\sigma_{\rm e})\neq0$, the Jacobian
$|\sigma'(u)|$ is bounded above and below by positive constants for small
$u$.  The exact edge theorem gives, uniformly there,
\begin{equation}
\label{eq:edge-in-u}
 -\log G(\sigma(u),0)=R(u)+O(R(u)^{1/2}).
\end{equation}
Moreover $R$ decreases as $u$ increases, tends to infinity at zero, and its
closed asymptotic implies
\begin{equation}
\label{eq:rate-short-interval}
 R\left(x\left(1-\frac1{R(x)}\right)\right)=R(x)+O(1).
\end{equation}
For the upper bound, monotonicity of $R$ and
\eqref{eq:edge-in-u} give
\[
 \mathcal C(d)
 \leq Cx\exp\{-R(x)+O(R(x)^{1/2})\}.
\]
For the lower bound, integrate only over
\[
 x\left(1-\frac1{R(x)}\right)\leq u\leq x.
\]
Equations \eqref{eq:edge-in-u}--\eqref{eq:rate-short-interval} give
\[
 \mathcal C(d)
 \geq \frac{cx}{R(x)}
       \exp\{-R(x)+O(R(x)^{1/2})\}.
\]
Now $|\log(x/R(x))|=O(\log R(x)+\log(1/x))=o(R(x)^{1/2})$.
The two bounds prove \eqref{eq:integrated-rate-sharp}.  Substituting
$x=d/2$ in the closed edge expansion proves
\eqref{eq:integrated-distortion-profile}; the integration error is smaller
than every displayed inverse-logarithmic correction.
\end{proof}

The preceding theorem can be inverted without a probabilistic or mixing
hypothesis.  It is an inversion of limiting intensity, not a statement about
the first zero observed at finite height.  Since $\eps_0$ is strictly
decreasing near $\sigma_{\rm e}$ and $G(\sigma,0)>0$ immediately to its left,
$\mathcal C(d)$ is strictly increasing for all sufficiently small $d>0$ and
therefore has the local inverse used below.

\begin{corollary}[Density--fidelity frontier]
\label{cor:density-fidelity-frontier}
Let $d(c)$ be the local inverse of $\mathcal C(d)$ near $c=0$, and put
$Y=\log(1/c)$.  Then
\begin{equation}
\label{eq:density-fidelity-inverse}
 \boxed{
 d(c)\sim
 2\left(K_{\rm e}\beta_{\rm e}^{1+\beta_{\rm e}}
   \right)^{1/\beta_{\rm e}}
 Y^{-1/\beta_{\rm e}}
 (\log Y)^{-(1+\beta_{\rm e})/\beta_{\rm e}}.
 }
\end{equation}
In particular, since $1/\beta_{\rm e}=\sigma_{\rm e}-1$ and
$(1+\beta_{\rm e})/\beta_{\rm e}=\sigma_{\rm e}$,
\begin{equation}
\label{eq:density-fidelity-geometric}
 d(c)\asymp
 [\log(1/c)]^{-(\sigma_{\rm e}-1)}
 [\log\log(1/c)]^{-\sigma_{\rm e}}.
\end{equation}
\end{corollary}

\begin{proof}
Set
\[
 \eps=A Y^{-1/\beta_{\rm e}}(\log Y)^{-q}.
\]
Substitution in the leading term of
\eqref{eq:integrated-distortion-profile} shows that
$q=(1+\beta_{\rm e})/\beta_{\rm e}$ and
$A=(K_{\rm e}\beta_{\rm e}^{1+\beta_{\rm e}})^{1/\beta_{\rm e}}$.
Since $d=2\eps$, this proves the assertion.
\end{proof}

There is also a useful finite-dimensional consequence.  It must be stated as
a sufficient certificate, rather than as the optimal cost of learning a
finite set of phases.

\begin{corollary}[Certified relative-phase recovery through $Q$]
\label{cor:finite-prime-certificate}
Fix $0<\delta<2$.  If a zero $\rho=\sigma+it$ satisfies
\begin{equation}
\label{eq:finite-prime-certificate}
 \mathcal D_\omega(\rho)\leq \delta^2Q^{-\sigma_{\rm e}},
\end{equation}
then
\[
 \left|\frac{\omega_p}{\omega_2}
       +\left(\frac p2\right)^{it}\right|\leq\delta
 \qquad(3\leq p\leq Q).
\]
The vertical intensity of the defect band certified by
\eqref{eq:finite-prime-certificate} is
\[
 \mathcal C_{Q,\delta}^{\rm cert}
 :=\mathcal C(\delta^2Q^{-\sigma_{\rm e}}),
\]
and, for fixed $\delta$ as $Q\to\infty$,
\begin{equation}
\label{eq:finite-prime-cost}
 \boxed{
 -\log\mathcal C_{Q,\delta}^{\rm cert}
 \sim \kappa_{\rm cert}\delta^{-2\beta_{\rm e}}
       \left(\frac Q{\log Q}\right)^{1+\beta_{\rm e}},
 \qquad
 \kappa_{\rm cert}
 =K_{\rm e}2^{\beta_{\rm e}}
  \sigma_{\rm e}^{-(1+\beta_{\rm e})}.
 }
\end{equation}
Thus, by the prime number theorem,
\begin{equation}
\label{eq:finite-prime-cost-pi}
 -\log\mathcal C_{Q,\delta}^{\rm cert}
 \sim\kappa_{\rm cert}\delta^{-2\beta_{\rm e}}
       \pi(Q)^{1+\beta_{\rm e}}.
\end{equation}
\end{corollary}

\begin{proof}
For $p\leq Q$ and $\sigma\leq\sigma_{\rm e}$,
\Cref{prop:distortion-defect} gives
\[
 \left|\frac{\omega_p}{\omega_2}
       +\left(\frac p2\right)^{it}\right|^2
 \leq \mathcal D_\omega(\rho)p^\sigma
 \leq\delta^2Q^{-\sigma_{\rm e}}Q^{\sigma_{\rm e}}
 =\delta^2.
\]
Substitute $d=\delta^2Q^{-\sigma_{\rm e}}$ into
\eqref{eq:integrated-distortion-profile}.  Since
$\sigma_{\rm e}\beta_{\rm e}=1+\beta_{\rm e}$ and
$\log(2/d)=\sigma_{\rm e}\log Q+O_\delta(1)$, one obtains
\eqref{eq:finite-prime-cost}.  Equation
\eqref{eq:finite-prime-cost-pi} follows from
$\pi(Q)\sim Q/\log Q$.
\end{proof}

\begin{remark}[What the cost law does not say]
\label{rem:not-first-hit}
The condition \eqref{eq:finite-prime-certificate} is sufficient, not
necessary.  The first prime-indexed phases may happen to be recovered accurately while
the distortion is carried by larger primes.  Thus
$\mathcal C_{Q,\delta}^{\rm cert}$ is the intensity of a canonically
certified family, not the optimal sampling rate for finite-dimensional phase
recovery.

Likewise, $1/\mathcal C(d)$ is a reciprocal limiting intensity.  The limit in
\eqref{eq:distortion-counting-law} first sends $T\to\infty$ with $d$ fixed.
It does not locate the first zero of distortion at most $d$, and one may not
replace $d$ by a moving threshold $d(T)$ without a uniform effective divisor
law.  In particular, the present theorem does not prove an extreme-value law
for the rightmost zero.  The hard core prevents clustering at very small
separation, but it does not turn a mean density into a first-hitting estimate.
\end{remark}

\subsection{The exact identity at a nonzero target}
\label{subsec:nonzero-a-distortion}

The elementary phase identity has a fully rigorous nonzero-target form.  We
record it separately because the corresponding logarithmic density cost
requires an outer-edge theorem for $G(\sigma,a)$ that is not part of the
zero-edge theorem proved above.

\begin{proposition}[Exact distortion at a nonzero $a$-point]
\label{prop:nonzero-a-distortion}
Let $a\neq0$, and suppose
$F_\omega(\rho)=a$ with $\rho=\sigma+it$ and $\sigma>1$.  Put
\[
 \eps_a(\sigma)=\cP(\sigma)-|a|,
\]
and define
\[
 \mathcal D_{\omega,a}(\rho)
 :=\sum_p p^{-\sigma}
 \left|\omega_p-\frac a{|a|}p^{it}\right|^2.
\]
Then
\begin{equation}
\label{eq:nonzero-a-exact-distortion}
 \boxed{\mathcal D_{\omega,a}(\rho)=2\eps_a(\sigma).}
\end{equation}
In particular,
\begin{equation}
\label{eq:nonzero-a-coordinate}
 \left|\omega_p-\frac a{|a|}p^{it}\right|
 \leq\sqrt{2\eps_a(\sigma)p^\sigma}
 \qquad(p\in\PP).
\end{equation}
Consequently any sequence of $a$-points approaching the outer edge
$E(a)=\cP^{-1}(|a|)$ recovers all coefficient phases.
\end{proposition}

\begin{proof}
Write $a=|a|e^{i\alpha}$ and put
$\zeta_p=\omega_pe^{-it\log p-i\alpha}$.  The $a$-point equation is
\[
 \sum_pp^{-\sigma}\zeta_p=|a|.
\]
Taking real parts and using $|1-\zeta|^2=2(1-\Re\zeta)$ gives
\[
 \sum_pp^{-\sigma}|1-\zeta_p|^2
 =2(\cP(\sigma)-|a|).
\]
Multiplication of each difference by the appropriate unimodular factor gives
\eqref{eq:nonzero-a-exact-distortion}; the coordinate bound follows by
discarding all other summands.
\end{proof}

\begin{remark}[Conditional nonzero-target extension]
\label{rem:conditional-a-edge}
It is natural to expect the preceding small-distortion law at every fixed
$a\neq0$.  With $E=E(a)$ and $\beta_a=(E-1)^{-1}$, the required input would
be the general outer-edge estimate below, where
\[
 H_0(E)=\int_0^\infty
 \bigl(x-\log I_0(x)\bigr)x^{-1/E-1}\,\dd x
\]
and
\[
 K(E):=\frac{E-1}{E}H_0(E)^{1+\beta_a}E^{-\beta_a}
       (1+\beta_a)^{-(1+\beta_a)}:
\]
\begin{align*}
 -\log G(\sigma,a)
 ={}&K(E)\eps_a(\sigma)^{-\beta_a}L_a^{-(1+\beta_a)}
 \\
 &\times\left(1+(1+\beta_a)\frac{\log L_a}{L_a}
  +O(L_a^{-1})\right),
 \qquad L_a=\log\frac1{\eps_a(\sigma)}.
\end{align*}
The same PNT--Cram\'er mechanism should apply after replacing the
dominant-prime cancellation by alignment of the full prime sum; on the edge
fibre the derivative weight tends to $|\cP'(E)|^2$.  Nevertheless, the
zero-edge theorem as stated proves only $a=0$.  A complete proof of the
full-prime tilted local-density lower bound and of the derivative-weight
comparison must be supplied before the displayed nonzero-target expansion,
or any cost law derived from it, is promoted to a theorem.  No such extension
is used in the proved results of this section.
\end{remark}
\section{Uniform small-target surgery and divisor monodromy}
\label{sec:small-target-surgery}

The weak escape theorem describes what happens to the horizontal divisor measure as
$a\to0$, while the hard-core theorem controls individual zeros in a right-hand band.
Together they give a stronger, pointwise description on every compact hard-core
subband: its bounded $a$-points are uniformly controlled holomorphic perturbations of
the zero divisor, whereas the missing mass is carried by a disjoint logarithmic lattice
at infinity.  These controlled persistent branches and the escaping branches have
different monodromy about the singular target $a=0$.

Throughout this section
\[
 F_\omega(s)=\sum_p\omega_pp^{-s},\qquad |\omega_p|=1,
 \qquad
 \cP^{\prime\prime}(\sigma)=\sum_p(\log p)^2p^{-\sigma}.
\]
For $\lambda>0$ put
\begin{align}
 \mathfrak f_\sigma(\lambda)
   &:=\sum_{p\ge3}p^{-\sigma}\left|1-\lambda\log(p/2)\right|,\label{eq:st-f}\tag{ST.1}\\
 q_\lambda(\sigma)
   &:=\frac{2^{-\sigma}-\mathfrak f_\sigma(\lambda)}{\lambda},
 &
 \mathfrak m(\sigma)&:=\sup_{\lambda>0}q_\lambda(\sigma),\label{eq:st-q}\tag{ST.2}\\
 \kappa_\lambda&:=\lambda^{-1}+\log2.
 \label{eq:st-kappa}\tag{ST.3}
\end{align}
The hard-core region is the open set
\[
 \mathcal H:=\{\sigma>1:\mathfrak m(\sigma)>0\}.
\]
The preceding hard-core theorem gives
$|F_\omega'(\rho)|\ge \mathfrak m(\Re\rho)$ at every zero with real part in
$\mathcal H$.  This agrees with the notation in
\Cref{thm:lp-planar-core}; no numerical description of $\mathcal H$ is needed below.

\subsection{A derivative bound stable under perturbation of the target}

The general nonzero-target estimate obtained by treating all prime terms symmetrically
is useful for large targets.  Near $a=0$ one should instead retain the distinguished
$p=2$ term.  This gives the estimate required for a uniform implicit-function theorem.

\begin{lemma}[Small-target derivative bound]
\label{lem:small-target-slope}
Let $z=\sigma+it$, $\sigma>1$, and suppose that $F_\omega(z)=a$.  Then, for every
$\lambda>0$,
\[
 \boxed{
 |F_\omega'(z)|
 \ge q_\lambda(\sigma)-\kappa_\lambda|a|.
 }
\tag{ST.4}
\]
At $a=0$, taking the supremum over $\lambda$ recovers the hard-core lower bound
$|F_\omega'(z)|\ge \mathfrak m(\sigma)$.
\end{lemma}

\begin{proof}
Write $U_p=\omega_pp^{-it}$, $a_p=p^{-\sigma}$, and
\[
 S=\sum_{p\ge3}a_pU_p,
 \qquad
 D=\sum_{p\ge3}\log(p/2)a_pU_p.
\]
Since $a_2U_2+S=a$,
\[
 |S|\ge 2^{-\sigma}-|a|.
\]
Moreover
\[
 F_\omega'(z)=-(\log2)a-D,
\]
and hence
\[
 S
 =\sum_{p\ge3}\bigl(1-\lambda\log(p/2)\bigr)a_pU_p
   -\lambda F_\omega'(z)-\lambda(\log2)a.
\]
The triangle inequality now gives
\[
 2^{-\sigma}-|a|
 \le \mathfrak f_\sigma(\lambda)+\lambda|F_\omega'(z)|
       +\lambda(\log2)|a|,
\]
which is (ST.4).
\end{proof}

The estimate can be made uniform without using a numerically computed threshold.
If $K\Subset\mathcal H$ is a compact interval, compactness supplies a finite set
$\Lambda_K\subset(0,\infty)$ such that
\[
 \mu_K:=\inf_{\sigma\in K}\max_{\lambda\in\Lambda_K}q_\lambda(\sigma)>0.
 \tag{ST.5}
\]
Indeed, choose one positive witness $q_\lambda(\sigma)>0$ at each point and take a
finite subcover.  Set
\[
 \kappa_K:=\max_{\lambda\in\Lambda_K}\kappa_\lambda.
 \tag{ST.6}
\]

\begin{corollary}[A small-target hard core]
\label{cor:small-target-hard-core}
Let $K=[\alpha,\beta]\Subset\mathcal H$ and let $0\le r<\mu_K/\kappa_K$.
Every $a$-point with $|a|\le r$ and real part in $K$ satisfies
\[
 |F_\omega'(z)|\ge \mu_K-\kappa_Kr>0.
 \tag{ST.7}
\]
In particular all such points are simple.  If $z\ne w$ are two such $a$-points,
then
\[
 |z-w|\ge d_{K,r},
 \tag{ST.8}
\]
where $d_{K,r}\in(0,\alpha-1)$ is the unique solution of
\[
 d_{K,r}\,\cP^{\prime\prime}(\alpha-d_{K,r})
 =2(\mu_K-\kappa_Kr).
 \tag{ST.9}
\]
Thus simplicity and planar hard core persist uniformly under sufficiently small
perturbations of the target.
\end{corollary}

\begin{proof}
The first assertion follows from Lemma~\ref{lem:small-target-slope}, choosing at each
$\sigma$ a maximizing witness in $\Lambda_K$.  Put $h=w-z$.  If
$|h|\ge\alpha-1$ there is nothing to prove.  Otherwise the segment from $z$ to $w$
lies in $\Re s>1$, and Taylor's formula gives
\[
 0=F_\omega(w)-F_\omega(z)
  =hF_\omega'(z)
   +h\int_0^1\bigl(F_\omega'(z+\theta h)-F_\omega'(z)\bigr)\dd\theta.
\]
Consequently
\[
 2|F_\omega'(z)|
 \le |h|\,\cP^{\prime\prime}(\Re z-|h|)
 \le |h|\,\cP^{\prime\prime}(\alpha-|h|).
\]
The left-hand side of (ST.9), viewed as a function of $d$, is strictly
increasing from $0$ to $+\infty$ on $(0,\alpha-1)$.  Equations (ST.7)--(ST.9)
therefore give the claim.
\end{proof}

\subsection{Uniform zero--\texorpdfstring{$a$}{a}-point continuation}

The next theorem is stated with three nested intervals.  This buffer is necessary:
an $O(|a|)$ displacement may cross a prescribed vertical boundary even when the
correspondence itself is perfectly regular.

\begin{theorem}[Uniform holomorphic continuation of the persistent divisor]
\label{thm:uniform-small-target-continuation}
Let
\[
 I\Subset J\Subset K\Subset\mathcal H
\]
be compact intervals.  There are constants $a_0,C_1,C_2>0$, depending only on these
intervals, such that the following assertions hold simultaneously for every unimodular
twist $\omega$ and every $|a|<a_0$.

\begin{enumerate}[label=\textup{(\roman*)}]
\item Every zero $\rho$ of $F_\omega$ with $\Re\rho\in J$ has a unique associated
$a$-point $\rho(a)$ in a fixed disc about $\rho$.  The map $a\mapsto\rho(a)$ is
holomorphic on $|a|<a_0$, and
\[
 |\rho(a)-\rho|\le C_1|a|,
 \qquad
 \left|\rho(a)-\rho-\frac{a}{F_\omega'(\rho)}\right|
 \le C_2|a|^2.
 \tag{ST.10}
\]

\item Every $a$-point $z$ with $\Re z\in I$ is $\rho(a)$ for a unique zero $\rho$
with $\Re\rho\in J$.  Thus, writing
\[
 Z_a(L):=\{z:F_\omega(z)=a,\ \Re z\in L\},
\]
one has the exact buffered identity
\[
 Z_a(I)
 =\{\rho(a):\rho\in Z_0(J),\ \Re\rho(a)\in I\}.
 \tag{ST.11}
\]

\item If $d_0>0$ is a uniform zero hard-core constant on $K$, then, after decreasing
$a_0$ if necessary,
\[
 |\rho_1(a)-\rho_2(a)|
 \ge d_0-2C_1|a|
 \ge \frac{d_0}{2}
 \tag{ST.12}
\]
for distinct zeros $\rho_1,\rho_2$ in $J$.
\end{enumerate}
All constants and conclusions are uniform in the height of the points and in the twist.
\end{theorem}

\begin{proof}
We give explicit safe choices.  Let
\[
 b:=\min\bigl\{
   \operatorname{dist}(I,\R\setminus J),
   \operatorname{dist}(J,\R\setminus K)
 \bigr\}>0,
\]
let $\mu=\mu_K$, $\kappa=\kappa_K$, and put
\[
 M:=\cP^{\prime\prime}(\min K).
\]
Let $d_0=d_{K,0}$ in Corollary~\ref{cor:small-target-hard-core}, and choose
\[
 0<\delta<\frac13\min\left\{b,d_0,\frac{\mu}{M}\right\}.
 \tag{ST.13}
\]
It is enough to take
\[
 a_0<\min\left\{
  \frac{\mu}{2\kappa},
  \frac{\mu^2}{8M},
  \frac{\mu\delta}{4},
  \frac{\mu d_0}{8}
 \right\}.
 \tag{ST.14}
\]

Let $F_\omega(\rho)=0$ and $\Re\rho\in J$.  On $|w|=\delta$,
\[
 \left|F_\omega(\rho+w)-F_\omega'(\rho)w\right|
 \le\frac{M}{2}\delta^2,
\]
whereas $|F_\omega'(\rho)|\ge\mu$.  Equations (ST.13)--(ST.14) and
Rouch\'e's theorem show that $F_\omega-a$ has exactly one zero in
$B(\rho,\delta)$.  Repeating the comparison on the circle
\[
 |w|=\frac{2|a|}{\mu}
\]
gives
\[
 |\rho(a)-\rho|\le\frac{2|a|}{\mu}.
 \tag{ST.15}
\]
The zero is simple by Corollary~\ref{cor:small-target-hard-core}; the implicit-function
theorem and uniqueness in the fixed disc make $\rho(a)$ holomorphic throughout the
target disc.  Taylor's formula once more yields
\[
 \left|\rho(a)-\rho-\frac{a}{F_\omega'(\rho)}\right|
 \le \frac{M}{2\mu}|\rho(a)-\rho|^2
 \le\frac{2M}{\mu^3}|a|^2.
 \tag{ST.16}
\]
Thus one may take $C_1=2/\mu$ and $C_2=2M/\mu^3$ in part~(i).

Conversely, suppose that $F_\omega(z)=a$ and $\Re z\in I$.  By (ST.7),
$|F_\omega'(z)|\ge\mu/2$.  On the circle
\[
 |w|=\frac{4|a|}{\mu},
\]
the linear term $F_\omega'(z)w$ has modulus at least $2|a|$, while the quadratic
remainder has modulus at most
\[
 \frac{M}{2}\left(\frac{4|a|}{\mu}\right)^2<|a|.
\]
Rouch\'e's theorem, now applied to $F_\omega(z+w)$, supplies exactly one zero within
$4|a|/\mu$ of $z$.  By the buffer and (ST.14) this zero lies in $J$.  The $a$-point
$z$ must be the branch already constructed from that zero, because the fixed
$\delta$-disc contains only one $a$-point.  The zero hard core and
$\delta<d_0/3$ make the correspondence injective and prove (ST.11).

Finally (ST.15) and the triangle inequality give
\[
 |\rho_1(a)-\rho_2(a)|
 \ge |\rho_1-\rho_2|-\frac{4|a|}{\mu}
 \ge d_0-\frac{4|a|}{\mu}.
\]
The last condition in (ST.14) gives (ST.12).  None of the estimates involves
$\Im\rho$ or the phases $\omega_p$, proving the asserted uniformity.
\end{proof}

\begin{remark}[Sharp scope of the continuation theorem]
The buffers in Theorem~\ref{thm:uniform-small-target-continuation} cannot simply be
deleted from a statement about raw rectangular counts: a branch can cross either
vertical boundary.  More importantly, below $\mathcal H$ arbitrary twists may have
multiple zeros.  There one retains local convergence of divisors with multiplicity,
but not a uniform decomposition into simple one-point branches.
\end{remark}

\subsection{The zero-free corridor and two-scale degeneration}

We now return to the full prime support.  Put
\[
 A(\sigma)=\sum_pp^{-\sigma},
 \qquad
 B(\sigma)=2^{-\sigma}-\sum_{p\ge3}p^{-\sigma}.
 \tag{ST.17}
\]
Thus $B(\sigma_{\rm e})=0$ and $B(\sigma)>0$ for
$\sigma>\sigma_{\rm e}$.  The phase-annulus criterion says that the real part of an
$a$-point with $r=|a|$ must satisfy
\[
 \max\{B(\sigma),0\}\le r\le A(\sigma).
 \tag{ST.18}
\]

\begin{proposition}[A separating corridor]
\label{prop:small-target-corridor}
Fix $\beta>\sigma_{\rm e}$.  For all sufficiently small $r>0$ there are numbers
\[
 \beta<\ell(r)<u(r)
\]
such that no $a$-point with $|a|=r$ has real part in
$[\beta,\ell(r))$, and every $a$-point to the right of $\beta$ lies in
$[\ell(r),u(r)]$.  The endpoints are characterized on the far-right monotone branch by
\[
 B(\ell(r))=r,
 \qquad
 A(u(r))=r,
 \tag{ST.19}
\]
and
\[
 \ell(r),u(r)
 =\frac{\log(1/r)}{\log2}
   +O\!\left(r^{\log_2 3-1}\right).
 \tag{ST.20}
\]
Consequently the bounded and escaping parts of the $a$-point divisor are separated by
a genuine zero-free corridor whose length tends to infinity.
\end{proposition}

\begin{proof}
Choose $M>\beta$ so large that $B$ is strictly decreasing on $[M,\infty)$; this is
possible because
\[
 B'(\sigma)=-(\log2)2^{-\sigma}
   +\sum_{p\ge3}(\log p)p^{-\sigma}<0
\]
eventually.  Since $B$ is positive on $[\beta,M]$, its minimum there is positive.
For $r$ below that minimum, $B(\sigma)>r$ throughout $[\beta,M]$ and until the unique
far-right solution $\ell(r)$ of $B(\sigma)=r$.  Criterion (ST.18) excludes every
$a$-point in this interval.  The function $A$ is strictly decreasing to zero, so its
unique solution $u(r)$ satisfies $u(r)>\ell(r)$; (ST.18) also excludes points to the
right of $u(r)$.

Finally
\[
 \sum_{p\ge3}p^{-\sigma}=O(3^{-\sigma})
 \qquad(\sigma\to\infty).
\]
Both equations in (ST.19) therefore have the form
$2^{-\sigma}=r\bigl(1+O(r^{\log_2 3-1})\bigr)$, which proves (ST.20).
\end{proof}

Combining Proposition~\ref{prop:small-target-corridor} with
Proposition~\ref{prop:escape-lattice} gives the following divisor-level strengthening
of Theorem~\ref{thm:mass-escape}.  On each compact subband of $\mathcal H$, the bounded
part is the canonical divisor
\[
 \sum_{F_\omega(\rho)=0}
 \left[\rho+\frac{a}{F_\omega'(\rho)}+O(|a|^2)\right],
 \tag{ST.21}
\]
where the statement is interpreted locally and the error is uniform.  The points
beyond the separating corridor form the disjoint family
\[
 \rho_{a,k}
 =\frac{\log(1/|a|)}{\log2}
 +i\frac{\arg\omega_2-\arg a+2\pi k}{\log2}
 +O\!\left(|a|^{\log_2 3-1}\right),
 \qquad k\in\mathbb Z.
 \tag{ST.22}
\]
Thus weak convergence to an atom at infinity reflects an actual two-scale degeneration
of the divisor, not merely a cancellation in an averaged counting formula.

\subsection{Monodromy of the two families}

\begin{theorem}[Identity versus shift monodromy]
\label{thm:small-target-monodromy}
Fix a sufficiently small $r>0$.  Analytically continue (a) any persistent branch
furnished by \Cref{thm:uniform-small-target-continuation} and (b) the far-right
escaping family along the positively oriented loop
$a(\theta)=re^{i\theta}$, $0\le\theta\le2\pi$.

\begin{enumerate}[label=\textup{(\roman*)}]
\item Every persistent branch furnished by
Theorem~\ref{thm:uniform-small-target-continuation} returns to itself.
\item With the labeling in (ST.22), continuation of the escaping family acts by the
bilateral shift
\[
 \boxed{\rho_{a,k}\longmapsto\rho_{a,k-1}.}
 \tag{ST.23}
\]
The reverse loop gives $k\mapsto k+1$.
\end{enumerate}
Hence the controlled persistent sheets extend across $a=0$, whereas the escaping
sheet is an infinite cyclic covering of the punctured target disc.  No assertion is
made here about possible branching of bounded points outside the chosen hard-core
subbands.
\end{theorem}

\begin{proof}
A persistent branch is holomorphic on the full disc $|a|<a_0$, including its centre,
so the monodromy theorem gives (i).  The escaping roots are simple by
Proposition~\ref{prop:escape-lattice}, and hence admit analytic continuation along the
loop.  Increasing the continuous argument of $a$ by $2\pi$ changes the leading
ordinate in (ST.22) from the one labeled $k$ to the one labeled $k-1$.  The error is
uniformly $o(1)$, while consecutive leading ordinates are separated by
$2\pi/\log2$.  For sufficiently small $r$, the disjoint Rouch\'e discs defining the
labels therefore force the terminal permutation to be exactly $k\mapsto k-1$.
\end{proof}

\begin{remark}[Least-prime recovery in three equivalent forms]
For a general infinite support $S$ whose two least primes are $p_0<p_1$, replace
$\log(p/2)$ by $\log(p/p_0)$ in (ST.1).  Lemma~\ref{lem:small-target-slope} becomes
\[
 |F_{S,\omega}'(z)|
 \ge
 \frac{p_0^{-\sigma}-
  \sum_{p\in S,\,p>p_0}p^{-\sigma}|1-\lambda\log(p/p_0)|}{\lambda}
 -\left(\lambda^{-1}+\log p_0\right)|a|.
 \tag{ST.24}
\]
The escaping error exponent is $\log p_1/\log p_0-1$, the vertical lattice step is
$2\pi/\log p_0$, its density is $\log p_0/(2\pi)$, and the monodromy is again a
unit shift.  Thus the escape mass, reciprocal lattice spacing, and monodromy translation
all encode the same least support prime.
\end{remark}

\subsection{An even smooth expansion behind the quantized jump}

The atom at infinity is nonperturbative, but the bounded part has an ordinary smooth
expansion.  This separates the two effects to every finite order.

\begin{corollary}[Quantized jump plus an even smooth jet]
\label{cor:even-escape-jet}
Fix $\sigma_0>1$ and define
\[
 C_a(\sigma_0)=\int_{\sigma_0}^{\infty}G(\sigma,a)\dd\sigma.
\]
For every integer $N\ge0$, as $a\to0$ through nonzero values,
\[
 \boxed{
 C_a(\sigma_0)
 =C_0(\sigma_0)+\frac{\log2}{2\pi}
  +\sum_{j=1}^{N}c_j(\sigma_0)|a|^{2j}
  +O_{N,\sigma_0}(|a|^{2N+2}).
 }
\tag{ST.25}
\]
If $\Delta_a$ denotes the planar Laplacian in the target variable, then
\[
 c_j(\sigma_0)
 =\frac{1}{4^j(j!)^2}
   \int_{\sigma_0}^{\beta}(\Delta_a^jG)(\sigma,0)\dd\sigma,
 \tag{ST.26}
\]
where any fixed $\beta>\sigma_{\rm e}$ lying in the separating corridor for all
sufficiently small $|a|$ may be used.  Formula (ST.25) is a $C^\infty$ asymptotic jet;
no convergence of the infinite Taylor series is asserted.
\end{corollary}

\begin{proof}
Choose $\beta>\max\{\sigma_0,\sigma_{\rm e}\}$ and then restrict $|a|$ so that
Proposition~\ref{prop:small-target-corridor} applies.  On the $2$-dominance side of the
remote component the Jessen function is $J_a(\sigma)=-\sigma\log2$, while beyond the
outer value-support edge it is $J_a(\sigma)=\log|a|$.  The Riesz-measure identity is
$J_a''(\dd\sigma)=2\pi G(\sigma,a)\dd\sigma$ (and
$J_a''(\sigma)=2\pi G(\sigma,a)$ on each smooth piece).  Integration across
the remote component therefore gives its
exact mass
\[
 \frac{0-(-\log2)}{2\pi}=\frac{\log2}{2\pi}.
 \tag{ST.27}
\]
There is no density in the corridor, and $G(\sigma,0)=0$ for
$\sigma>\sigma_{\rm e}$, so
\[
 C_a(\sigma_0)-\frac{\log2}{2\pi}
 =\int_{\sigma_0}^{\beta}G(\sigma,a)\dd\sigma,
 \qquad
 C_0(\sigma_0)=\int_{\sigma_0}^{\beta}G(\sigma,0)\dd\sigma.
\]
The function $G$ is smooth and invariant under $a\mapsto e^{i\theta}a$.  Its Taylor
polynomial at $a=0$ therefore contains only powers of $|a|^2$, uniformly for
$\sigma\in[\sigma_0,\beta]$.  For a smooth radial function, the coefficient of
$|a|^{2j}$ is $(4^j(j!)^2)^{-1}\Delta_a^jG(\sigma,0)$.  Taylor's theorem followed by
integration proves (ST.25)--(ST.26).
\end{proof}

\begin{remark}[What has been strengthened]
Theorem~\ref{thm:mass-escape} gives weak convergence of horizontal measures.
Theorem~\ref{thm:uniform-small-target-continuation} identifies the bounded divisor on
compact subbands of $\mathcal H$ with individual, uniformly separated zero branches;
Proposition~\ref{prop:small-target-corridor}
separates those branches from the escaping lattice; Theorem~\ref{thm:small-target-monodromy}
distinguishes the two parts topologically; and Corollary~\ref{cor:even-escape-jet}
resolves the remaining bounded mass to every finite perturbative order.  None of these
statements requires an uncertified decimal value for the hard-core threshold.
\end{remark}

\section{Limitations and next problems}
\label{sec:limits}

The fixed-level correlation and phase-recovery questions are now resolved at
the level stated above.  The remaining bottlenecks are different and should not
be hidden by those qualitative closures.

\begin{enumerate}
\item \textbf{Effective height errors and a moving edge.}
The divisor and point-process limits use qualitative unique ergodicity of the
Kronecker flow.  No uniform rate is proved as a test function approaches a
diagonal, a hard-core boundary, or the outer edge.  In particular, the limiting
edge density does not determine the first height at which a point of a prescribed
defect occurs.  A moving-edge theorem would require effective
high-dimensional equidistribution with a relative error meaningful when the
expected count is of order one.

\item \textbf{Gaussian prefactors and the remaining edge profiles.}
The logarithmic zero-edge rate and its integrated small-distortion consequence
are exact, but a sharp Gaussian prefactor still requires a tilted local central
limit theorem.  The exact phase-defect identity also holds at every fixed
nonzero target, but the analogous nonzero-target density--fidelity law remains
conditional on a full edge-profile theorem for $G(\sigma,a)$.  Regular internal
band edges should have related log-weighted profiles; proving them requires
uniform tilted density estimates and conditional nondegeneracy of the
derivative weight.  Smooth fixed-parameter Kac--Rice densities do not by
themselves supply this uniformity.

\item \textbf{Explicit local laws and the hard-core transition.}
The Palm nearest-neighbour and hole laws are exact finite
inclusion--exclusion expressions in the $K_k$, not closed elementary functions.
The complete region in which the joint jet value $(0,0)$ is interior to its
support, and hence the exact boundary between positive quadratic coefficient and
hard core, remains to be classified analytically.  Decimal transition values,
double-zero examples, and equal-height examples require reproducible interval
certificates before they may be used as theorem statements.

\item \textbf{Finite-height phase tomography and optimal sampling.}
The defect identity and canonical edge matching give deterministic stability
from a supplied near-edge point or germ.  What remains unavailable is an
effective deterministic bound for how high one must search to find that datum.
Reciprocal limiting intensity is only a frequency scale, not a first-occurrence
theorem.  The stated finite-prime criterion is a sufficient certificate; the
optimal intensity for reconstructing a prescribed finite set of phases is not
determined.  Zero data retain the unavoidable common-rotation ambiguity, and
unanchored data additionally retain the vertical-flow ambiguity.

\item \textbf{Small targets beyond the hard-core subbands.}
Uniform holomorphic continuation and identity monodromy of the persistent
divisor are proved on compact subbands where the symbolic derivative margin is
positive.  Across multiple-zero caustics that labeling can branch, and a global
description of the persistent divisor under arbitrary target loops remains
open.  The escaping lattice and its unit-shift monodromy are controlled only in
the small-target regime; no uniform statement is asserted while the target and
height vary together.

\item \textbf{Crossing $\Re s=1$.}
The endpoint transfer uses the special factorization
$e^{\cP}=\zeta H$.  It neither counts zeros on $\Re s=1$ nor gives a general
twisted continuation theorem below that line.  The singularities of
$\zeta(ks)$ and the required branch choices lie outside the present
almost-periodic framework.

\item \textbf{Uniform moving-prime geometry.}
The internal bands are classified at each fixed derivative order, and the
two-prime lattice theorem fixes the adjacent pair before the order tends to
infinity.  Uniformity for prime pairs growing with the derivative order, together
with interval-certified computations of large birth levels, remains open.
\end{enumerate}

\medskip
\noindent
The integrated structural conclusion is
\[
\boxed{
\begin{gathered}
 \substack{\text{vertical Haar laws and}\\
           \text{joint Kac--Rice densities}}
   \Longrightarrow \text{exact simple counts and the full local process},\\
 \mathfrak m(\sigma)>0
   \Longrightarrow \text{uniform simplicity, hard cores, and persistent branches},\\
 -\log G(\sigma,0)\sim\mathscr I_{\rm e}(\sigma)
   \Longrightarrow \text{the PNT edge profile and the exact fidelity frontier},\\
 \text{near-edge roots and their rooted germs}
   \Longrightarrow \text{relative phases and edge-orbit rigidity},\\
 a\to0
   \Longrightarrow
   \substack{\text{quantized escape of }\log p_0/(2\pi)\\
             \text{and identity/shift divisor monodromy}},\\
 \substack{\text{right-edge moduli}+\\
           \text{anchored near-edge ordinates}}
   \Longrightarrow \text{the complete coefficient system},\\
 \text{component masses and moving lattices}
   \Longrightarrow \text{prime-pair and relative-phase recovery}.
\end{gathered}}
\]

\phantomsection
\addcontentsline{toc}{section}{Acknowledgement and status note}
\section*{Acknowledgement and status note}

The manuscript distinguishes throughout among projection closure, boundary
attainment, positive limiting density, deterministic recurrence, and effective
first occurrence.  It also distinguishes multiplicity, distinct-point, and
simple-point counts, and separates self-inclusive from factorial correlations.
In particular, the smoothed self-inclusive two-point function is Bohr almost
periodic, while deletion of the diagonal leaves an $AP+C_0$ function rather
than, in general, another Bohr almost-periodic function.

The polygonal projection mechanism and multiplicity-weighted Jessen mean theory
are classical foundations.  The exact edge underlying Conjecture~1 of Belovas,
\v{C}epaityt\.{e}, and Sabaliauskas is therefore not claimed as a new general
phenomenon, and the independent specialization of Camargo is explicitly
acknowledged.  The claims of this paper concern the specialized density,
distinct/simple divisor, local-process, hard-core, edge-profile, phase-recovery,
rooted-rigidity, zero-edge rate--distortion, small-target monodromy, and inverse
consequences proved from those foundations.  The rate--distortion theorem is a
limiting-intensity statement at the zero edge: it is neither a deterministic
first-hit estimate nor, in its finite-prime corollary, an optimality theorem.
The nonzero-target defect identity is unconditional, but its matching edge-cost
asymptotic is not asserted without a corresponding nonzero-target edge theorem.
The persistent small-target divisor theorem is restricted to compact hard-core
subbands.  Numerical values are illustrative unless accompanied by an explicit
interval certificate, and no historical ``first'' claim is made beyond the
cited targeted search.  A full MathSciNet and zbMATH priority review remains
appropriate before journal submission.

\phantomsection
\addcontentsline{toc}{section}{References}
\begin{thebibliography}{99}

\bibitem{BakerHarmanPintz2001}
R.~C.~Baker, G.~Harman, and J.~Pintz,
\emph{The difference between consecutive primes. II},
Proc. London Math. Soc. (3) \textbf{83} (2001), no.~3, 532--562.
\href{https://doi.org/10.1112/plms/83.3.532}{doi:10.1112/plms/83.3.532}.

\bibitem{BCS2025}
I.~Belovas, R.~\v{C}epaityt\.{e}, and M.~Sabaliauskas,
\emph{On the zero-free region and the distribution of zeros of the prime zeta
function},
An. \c{S}tiin\c{t}. Univ. Ovidius Constan\c{t}a Ser. Mat. \textbf{33} (2025), no.~2,
27--44.
\href{https://doi.org/10.2478/auom-2025-0017}{doi:10.2478/auom-2025-0017}.

\bibitem{BGT1987}
N.~H.~Bingham, C.~M.~Goldie, and J.~L.~Teugels,
\emph{Regular Variation},
Encyclopedia of Mathematics and its Applications 27, Cambridge University Press,
1987.

\bibitem{BinderPauliSaidak2015}
T.~Binder, S.~Pauli, and F.~Saidak,
\emph{Zeros of high derivatives of the Riemann zeta function},
Rocky Mountain J. Math. \textbf{45} (2015), no.~3, 903--926.
\href{https://doi.org/10.1216/RMJ-2015-45-3-903}
{doi:10.1216/RMJ-2015-45-3-903}.

\bibitem{Bauer1916}
M.~Bauer,
\emph{Zur Theorie der algebraischen Zahlk\"orper},
Math. Ann. \textbf{77} (1916), 353--356.
\href{https://doi.org/10.1007/BF01475865}{doi:10.1007/BF01475865}.

\bibitem{BorchseniusJessen1948}
V.~Borchsenius and B.~Jessen,
\emph{Mean motions and values of the Riemann zeta function},
Acta Math. \textbf{80} (1948), 97--166.
\href{https://doi.org/10.1007/BF02393647}{doi:10.1007/BF02393647}.

\bibitem{Camargo2025}
A.~Pierro de Camargo,
\emph{On the analytic continuation, computation and zeros of the prime zeta function},
preprint, 18 November 2025.
\url{https://www.researchgate.net/publication/397738720}.

\bibitem{Connes1994}
A.~Connes,
\emph{Noncommutative Geometry},
Academic Press, 1994.

\bibitem{Froberg1968}
C.-E.~Fr\"oberg,
\emph{On the prime zeta function},
BIT \textbf{8} (1968), 187--202.
\href{https://doi.org/10.1007/BF01933420}{doi:10.1007/BF01933420}.

\bibitem{JessenTornehave1945}
B.~Jessen and H.~Tornehave,
\emph{Mean motions and zeros of almost periodic functions},
Acta Math. \textbf{77} (1945), 137--279.
\href{https://doi.org/10.1007/BF02392225}{doi:10.1007/BF02392225}.

\bibitem{JessenWintner1935}
B.~Jessen and A.~Wintner,
\emph{Distribution functions and the Riemann zeta function},
Trans. Amer. Math. Soc. \textbf{38} (1935), 48--88.
\href{https://doi.org/10.1090/S0002-9947-1935-1501802-3}
{doi:10.1090/S0002-9947-1935-1501802-3}.

\bibitem{LandauWalfisz1920}
E.~Landau and A.~Walfisz,
\emph{\"Uber die Nichtfortsetzbarkeit einiger durch Dirichletsche Reihen definierter
Funktionen},
Rend. Circ. Mat. Palermo \textbf{44} (1920), 82--86.

\bibitem{Lee2014}
Y.~Lee,
\emph{On the zeros of Epstein zeta functions},
Forum Math. \textbf{26} (2014), 1807--1836.
\href{https://doi.org/10.1515/forum-2012-0057}{doi:10.1515/forum-2012-0057};
\href{https://arxiv.org/abs/1204.6297}{arXiv:1204.6297}.

\bibitem{Levin1996}
B.~Ya.~Levin,
\emph{Lectures on Entire Functions},
Translations of Mathematical Monographs 150, American Mathematical Society, 1996.

\bibitem{MontgomeryTenenbaum2017}
H.~L.~Montgomery and G.~Tenenbaum,
\emph{On multiplicative compositions of integers},
Mathematika \textbf{63} (2017), no.~3, 1081--1090.
\href{https://doi.org/10.1112/S0025579317000249}
{doi:10.1112/S0025579317000249}.

\bibitem{Maynard2015}
J.~Maynard,
\emph{Small gaps between primes},
Ann. of Math. (2) \textbf{181} (2015), no.~1, 383--413.
\href{https://doi.org/10.4007/annals.2015.181.1.7}
{doi:10.4007/annals.2015.181.1.7}.

\bibitem{Nakamura2017PrimeZeta}
T.~Nakamura,
\emph{On zeros of derivatives of the prime zeta-function},
invited lecture, Number Theory Week 2017, 7 September 2017; listed in the
KAKENHI project 16K05077 activity record.
\url{https://kaken.nii.ac.jp/en/grant/KAKENHI-PROJECT-16K05077/}.

\bibitem{SepulcreVidal2022}
J.~M.~Sepulcre and T.~Vidal,
\emph{On the real projections of zeros of analytic almost periodic functions},
Carpathian J. Math. \textbf{38} (2022), 489--501.
\href{https://doi.org/10.37193/CJM.2022.02.18}{doi:10.37193/CJM.2022.02.18};
\href{https://arxiv.org/abs/1805.02041}{arXiv:1805.02041}.

\bibitem{Simon2005}
B.~Simon,
\emph{Trace Ideals and Their Applications}, 2nd ed.,
Mathematical Surveys and Monographs 120, American Mathematical Society, 2005.
\href{https://doi.org/10.1090/surv/120}{doi:10.1090/surv/120}.

\bibitem{Titchmarsh1986}
E.~C.~Titchmarsh, revised by D.~R.~Heath-Brown,
\emph{The Theory of the Riemann Zeta-Function}, 2nd ed.,
Oxford University Press, 1986.

\end{thebibliography}

\end{document}
